%% ISCA 2027 submission sample — ACM acmart (sigconf), following ISCA 2025 practice.
%% \documentclass[sigconf, screen, review]{acmart}
\documentclass[sigconf, screen, review]{acmart}

%% Submission helpers (ISCA 2025): page numbers + confidentiality banner
\settopmatter{printfolios=true}

%% Rights block suppressed for anonymous review submission
\setcopyright{none}
\acmConference[ISCA '27]{54th IEEE/ACM International Symposium on Computer Architecture}{June 5--9, 2027}{Atlanta, Georgia, USA}
\acmYear{2027}
\acmISBN{978-x-xxxx-xxxx-x/YY/MM}
\acmDOI{XXXXXXX.XXXXXXX}
\settopmatter{printacmref=false}

%%%%%%%%%%%---SETME-----%%%%%%%%%%%%%
\newcommand{\iscasubmissionnumber}{443}
%%%%%%%%%%%%%%%%%%%%%%%%%%%%%%%%%%%%

%% Rightsform / CCS can be filled for camera-ready; omit CCS for this sample
\keywords{Do, Not, Us, This, Code, Put, the, Correct, Terms, for, Your, Paper}

% anna packages
\usepackage{xspace}
\newcommand{\toolname}{SimSpect\xspace}
\usepackage{tikz}
\usetikzlibrary{arrows.meta, positioning, calc, bending}
\usepackage{array}
\usetikzlibrary{fit,backgrounds,shapes.geometric}
\usetikzlibrary{shapes.symbols,shapes.misc,decorations.pathmorphing}% pipeline figures: load globally (loading inside a figure's group breaks a second figure)
\usepackage{colortbl}
\input{tikz/sscolors}% figure palette (2026-11-03): blue = defense-specific, tan = speculative, lavender = protection set, red = transmitter/bug; delete this line for grayscale figures
\definecolor{ctarget}{gray}{0}% legacy B&W colors (2026-11-02), used only by hidden content: the old bug table under \iffalse in results.tex, sections/pendingfigures.tex and unused tikz files; new figures use tikz/sscolors; was HTML B3322C
\definecolor{copen}{gray}{0.35}% legacy B&W (unused by live figures); was HTML 8E5A9B
\definecolor{calloy}{gray}{0.3}% legacy B&W (unused by live figures); was HTML 2C5AA0
\definecolor{camulet}{gray}{0.62}% legacy B&W (unused by live figures); was HTML 8C6D4B
\definecolor{cclean}{gray}{0.15}% legacy B&W (unused by live figures); was HTML 2E8B57
\definecolor{cpipe}{gray}{0.45}% legacy B&W (unused by live figures); was HTML B8860B
\newcommand{\anna}[1]{\textcolor{violet}{\footnotesize[Anna: #1]}}
\newcommand{\caroline}[1]{\textcolor{blue}{\footnotesize[Caroline: #1]}}
\newcommand{\todo}[1]{\textcolor{orange}{\footnotesize[todo: #1]}}
\definecolor{cfigtodo}{HTML}{C2185B}% figure TODOs (new figures, 2026-11-02)
\newcommand{\figtodo}[1]{\textcolor{cfigtodo}{\footnotesize[fig todo: #1]}}
\newcommand{\reviewnote}[1]{\textcolor{red}{\footnotesize\textbf{[Note]} #1}}% review notes: suggested cuts, rewords and checks
\newcommand{\reviewtag}[1]{\textcolor{red}{#1}}% red [TODO: ...] tags marking text that was changed and should be checked
\definecolor{caufact}{HTML}{1565C0}\definecolor{cauclaim}{HTML}{00897B}\definecolor{cauterm}{HTML}{8E24AA}\definecolor{caufig}{HTML}{2E7D32}\definecolor{cauwrite}{HTML}{6D4C41}\definecolor{caucite}{HTML}{9E9D24}\definecolor{caudefense}{HTML}{D84315}\definecolor{curgent}{HTML}{FF1744}%
\DeclareRobustCommand{\auditnote}[3]{\textcolor{cau#1}{\footnotesize\textbf{[#2]} #3}}% fine-grained audit notes (2026-11-03), one color per audit: blue fact, teal claim, green figure, brown writing, olive cite (purple term and orange defense are defined but unused; vocabulary findings are the VOCAB urgent notes)
\DeclareRobustCommand{\labnote}[1]{\textcolor{gray}{\footnotesize\itshape[Note: #1]}}% visible note for lab readers (2026-11-03)
\newcommand{\notebr}{\par\noindent}% review mode: each note in body text on its own line (2026-11-03); empty when notes are hidden
\DeclareRobustCommand{\urgentnote}[1]{\textcolor{curgent}{\footnotesize\textbf{
[URGENT]} #1}}% big concerns to fix before sending (2026-11-03)
\newcommand{\urgentbr}{\par\noindent}\newcommand{\labbr}{\par\noindent}% urgent and lab notes each start on their own line even when review notes are hidden; emptied by switches (3) and (4)
% ===== NOTE SWITCHES: each line below HIDES one group of notes. Put % at the start of a line to SHOW that group again; remove the % to hide it. =====
\renewcommand{\todo}[1]{\unskip}\renewcommand{\todo}[1]{\unskip}\renewcommand{\figtodo}[1]{\unskip}% (1) hides your notes: [Anna: ...], [todo: ...], [fig todo: ...]
% \renewcommand{\anna}[1]{\unskip}\renewcommand{\todo}[1]{\unskip}\renewcommand{\figtodo}[1]{\unskip}% (1) hides your notes: [Anna: ...], [todo: ...], [fig todo: ...]
\renewcommand{\reviewnote}[1]{\unskip}\renewcommand{\reviewtag}[1]{\unskip}\renewcommand{\auditnote}[3]{\unskip}\renewcommand{\notebr}{}% (2) hides review notes: red [Note] notes, red [TODO: ...] tags, and the colored audit notes
% \renewcommand{\urgentnote}[1]{\unskip}\renewcommand{\caroline}[1]{\unskip}\renewcommand{\urgentbr}{}% (3) hides urgent notes ([URGENT], including VOCAB) and Caroline notes
% \renewcommand{\labnote}[1]{\unskip}\renewcommand{\labbr}{}% (4) uncomment to also hide the gray [Note: ...] notes meant for lab readers




\begin{document}

\title{SimSpect: Automated and Thorough Litmus Testing of  Simulated Spectre Defenses}

%% Banner in author position (double-blind); same pattern as ISCA 2025 sample
\author{ISCA 2027 Submission
  \textbf{\#\iscasubmissionnumber} -- Confidential Draft -- Do NOT Distribute!!}
\renewcommand{\shortauthors}{ISCA 2027 Submission \#\iscasubmissionnumber}

\begin{abstract}
Spectre attacks exploit hardware mispredictions to coerce victim programs into leaking secrets over hardware side channels via transient (i.e., uncommitted) instructions. To provide software with a specific security guarantee against these attacks at minimal performance cost, many Spectre defenses require changes to conventional processor microarchitectures. These changes are usually prototyped in microarchitectural simulators like gem5, which can evaluate a defense implementation's performance and functional correctness by running sizable programs, but offer no analogous mechanism for checking whether it upholds its security guarantee. As a result, subtle security bugs can go undetected, voiding a defense's security guarantee and potentially understating its performance overhead.

We propose \toolname{}, a three-part framework for automated, thorough security validation of simulated Spectre defenses. First, a new specification language enables formalizing a Spectre defense's security guarantee as a predicate identifying which (leaky) program executions the defense forbids. Second, a synthesizer uses such a specification to generate a compact, targeted suite of \textit{security litmus tests}---small programs, each paired with a target execution (i.e., specific microarchitectural control-/data-flows) each exhibiting forbidden leakage. Third, an executor actively steers each test toward its target execution on a simulated defense prototype (here, in gem5) and checks whether the defense prevents it. Applied to gem5 prototypes of five publicly available Spectre defenses and one obtained from its authors, \toolname{} uncovers all bugs that were found by prior state-of-the-art simulator tester AMuLeT, and are in-scope for the defense, plus six new bugs. Compared to AMuLeT, SimSpect needs at least $23\times$ fewer tests to find each known bug.

% We propose \textbf{SimSpect}, a rigorous, targeted, and automated approach for validating early-stage simulator prototypes of hardware support for side-channel defenses. We present a formal language intermediate representation for specifying leakage prohibited by a defense as well as a validator that generates and executes a comprehensive litmus test suite from this specification, combining into a pipeline for rigorous defense validation and iteration. 
% Applying SimSpect to all accessible simulator prototypes of state-of-the-art defenses, we replicate all in-scope previously known security bugs, and find the following previously unknown bugs: 
% (1) STT mistakenly fails to taint one register, 
% ReCon's underlying implementation of STT (2) skips a taint check on store-to-load forwarding and (3) tracks only the oldest taint source as opposed to all taint sources, 
% SpecLFB (4) misimplements its ROB walk order and (5) fails to protect under the shadow of a jmp instruction, and finally 
% (6) there is a visibility-point/squash timing race shared by STT and SPT.
% Compared to prior work, SimSpect decreases the number of tests needed to find known bugs on average by a factor of 6.5\anna{update}.


% find a bug in STT where it fails to taint an FPR, uncover a previously unknown visibility-point/squash timing race shared by STT and SPT, and find two bugs in ReCon's underlying implementation of STT: store-to-load forwarding bypasses its taint check, and an instruction that merges two tainted operands keeps only the older operand's taint. Finally, we show that our test generation corpus takes less time to rediscover known bugs, for instance, on SpecLFB, SimSpect reaches the known first-load bug in 6.5x fewer tests than AMuLeT.



% Hardware side-channel attacks pass secrets to instructions with data-dependent optimizations that leak victim data via timing or other indicators. Proposed defenses against these attacks include hardware modifications that are prototyped in computer architecture simulators such as gem5. However, these prototyped mechanisms frequently contain security bugs that invalidate their reported performance results, and there exists no comprehensive early-stage validation method for their proposed security guarantees. 
% Many hardware defenses have been proposed to mitigate the dangers posed by side-channel attacks, particularly Spectre attacks but they present performance guarantees for their hardware via simulator prototypes. There exists no comprehensive validation method for the security guarantees of such simulators, and most have been shown to contain security-violating bugs that invalidate performance results. We develop SimSpect, a rigorous, targeted, and automated approach for validating early-stage simulator prototypes of hardware support for side-channel defenses (SCDs). We present a formal language intermediate representation for specifying hardware SCDs and proving security properties as well as a validator that generates and executes a comprehensive litmus test suite from this specification, together forming a pipeline for rigorous defense validation and iteration.


% security bugs \anna{that violate their defense's proposed security guarantees} as well as \textit{six new ones}.
% \urgentbr
% \urgentnote{``All known security bugs'' needs its scope and count. AMuLeT reports nine bugs on these defenses (fu-2025-amulet.pdf p1); five are inside each defense's own guarantee (UV1, UV3, UV4, UV5, UV6) and SimSpect finds all five (UV3 only with the read-modify-write corpus and a cold cache, 407 of 590 tests vs 0 of 590 warm; UV5 in one test so far). The other four are out of scope by each paper (UV2, KV1, KV2, KV3). On ``six new ones'': you decided SpecLFB's walk order is UV6 reaching further than AMuLeT published, yet the abstract and the intro count it as new bug (4). Options: count five new bugs plus a wider form of UV6, or keep six and call (4) a new form of UV6. \textbf{PROPOSED TEXT} (edits the sentence just before this note; recommendation, using the first option): ``\ldots{} \toolname{} uncovers all five previously known bugs that violate their defense's own guarantee, as well as five new bugs and a wider form of a known one.'' \textit{Source:} simspect wiki/findings/mdsToRead/AMULET\_BUGS\_DETECTION\_ACCOUNTING.md sec.~2; URGENT\_TODO.md walk-order decision (2026-11-09).}
% \notebr
\end{abstract}

\maketitle

\section{Introduction}
% \textit{What is a hardware side channel, why are they bad, what is spectre in particular}
In a \textit{hardware side-channel attack} (HSCA), an unsafe instruction in a victim program creates operand-dependent hardware resource usage (e.g., load address operand-dependent cache hits/misses), and an attacker observes this resource usage (e.g., through its effect on program execution time) to infer the operand value. \textit{Constant-time} (CT) programming was long thought to offer a complete defense: it rewrites code to avoid passing secrets to unsafe instruction operands. However, \textit{Spectre attacks}\cite{spectre} (a class of HSCA) pass secrets to \textit{transient} (partially executed and eventually squashed due to hardware control-/data-flow mispredictions) unsafe instructions, revealing that secret program data can leak via hardware side channels even if it is never leaked or accessed architecturally. This means that even code that avoids passing any secrets to \textit{transmitters}, or instructions that leak one or more operands via an HSCA, can leak these secrets through transmitters under speculation. Spectre attacks circumvent pervasive software-level (e.g., CT programming or sandboxing) and hardware-level (e.g., privilege isolation) security policies, establishing HSCAs as a threat to all programs that hold secrets in architectural state.

% \textit{HSCDs exist for HSCAs }
Defending secret-processing programs against Spectre attacks is a prerequisite to running differently trusted code in the same process or on the same machine. Simply disabling the optimizations that enable these attacks incurs prohibitive performance overheads. Hence, numerous \textit{Spectre defenses} have been proposed that modify software, hardware, or both, with the goal of preventing speculative leakage of secrets. To prevent behavior from occurring speculatively without disabling optimization extensively, it is necessary to process information about speculation at runtime, so many Spectre defenses implement hardware mechanisms. We consider the class of defenses with some hardware component — such defenses can delay or taint-track speculative instructions until they are safe~\cite{yu2019stt,spt,nda,specshield,loughlin2021dolma,sakalis2019delay,recon,doppelganger,protean,cheng2024speclfb}, make speculative cache state invisible~\cite{yan2018invisispec,khasawneh2019safespec,ainsworth2020muontrap,ainsworth2021ghostminion}, undo state updates after a squash~\cite{saileshwar2019cleanupspec}, or execute speculative instructions data-obliviously~\cite{sdo}.

The performance implications of these hardware modifications are assessed via early-stage (pre-RTL) prototyping in computer architecture simulators, especially gem5\cite{gem5}. These prototypes aim to maximize performance while enforcing some advertised security guarantee \anna{should we use}. However, due to the largely inherent tradeoff between performance and security, subtle bugs that undermine a defense’s security may lead to severe underestimates of its performance overhead relative to an insecure baseline, rendering a prototype useless for predicting defense viability. 
Such bugs are common in state-of-the-art Spectre defenses: testing four simulated defenses from 2018 to 2024 (InvisiSpec, CleanupSpec, STT and SpecLFB), AMuLeT uncovered three known and six previously unknown bugs and vulnerabilities within three hours of testing~\cite{amulet}\cite{nickscommittoSTT}.

% not cycle accurate

Simulated Spectre defenses in hardware must be “correct enough” to yield reliable performance estimates to avoid wasted RTL development effort on nonviable proposals. At this stage of the design process, correctness is not formally proven but rather validated through testing, as formal proof would require a large investment of time and resources and designs are subject to change. However, \textit{security} correctness in Spectre defenses is a \textit{microarchitectural} property requiring knowledge of side channels, rather than a purely \textit{architectural} function, like functional correctness. While functional correctness can be reasonably validated by running benchmarks and checking their architectural outputs, it is more difficult to rigorously validate side-channel security guarantees \cite{spec2017}. There are three main reasons for this. First, there is no established formal framework for specifying the security properties of a defense within the space of possible Spectre attacks, nor terminology to specify how a violation occurs. Second, the existing concept of a litmus test is a short program, which by its nature can only specify architectural behavior, so is unable to consistently induce potential violations. Third, observing violating microarchitectural leakage requires probing microarchitectural state, which is not straightforward. 
% \todo{cite the spectre benchmarks}

% First, there is a lack of formal specification to describe the security guarantees a particular defense should uphold, making it difficult to test for this property. Second, there is no good suite of tests that exercises these security guarantees. \anna{is this a weird framing, I'm trying to set up the need for formal description} \anna{i need a consistent term for the overall security guarantee contract thing}

% \anna{it's important to note this is early in the design process => validation not verification, want to iterate, so we are working with test suites and checking, not with formal proofs}
There exists no principled way to validate security guarantees of defenses at this stage, which is necessary to do fast high-level iteration.
Existing validation either tests fabricated CPUs against generic speculation contracts~\cite{revizor,oleksenko2023hideseekspectresefficient}, verifies or fuzzes RTL designs~\cite{fadiheh2023exhaustive,contractshadow,rtl2upath,hur2022specdoctor,ghaniyoun2021introspectre}, or fuzzes simulator prototypes against a sequential baseline~\cite{amulet}; none checks a prototype against the defense's own security guarantee.
\notebr
\reviewnote{\textbf{Citation check} (asked for in an earlier source comment here): every key in this sentence and in the benchmarks sentence below exists in references.bib and is the intended paper: Revizor (ASPLOS'22), Hide and Seek with Spectres (S\&P'23), Fadiheh et al.'s exhaustive RTL detection (IEEE TC'23), Contract Shadow Logic (ASPLOS'25), RTL2M$\mu$PATH (MICRO'24), SpecDoctor (CCS'22), IntroSpectre (ISCA'21), AMuLeT (ASPLOS'25), Kocher's 15 Spectre v1 examples (2018) and Canella et al.\ (USENIX Security'19). AMuLeT's abstract confirms ``3 known and 6 unknown bugs and vulnerabilities, within 3 hours of testing'' (fu-2025-amulet.pdf p1). Three small points: (1) Pensieve~\cite{pensieve} (ISCA'23), which model-checks early-stage defense designs against speculative non-interference (yang-2023-pensieve.pdf abstract), is missing here and is the closest prior work for ``at this stage''; (2) \cite{spec2017} in the previous paragraph follows the security sentence but supports ``running benchmarks and checking their architectural outputs''; (3) in the second paragraph, \cite{nickscommittoSTT} (Mosier's STT commit 18e304f) follows the AMuLeT count, but that bug is not one of AMuLeT's nine; give it its own clause if it is meant as another example.}
\notebr
% When generating a set of litmus tests to validate some
% - we need litmus test suite
% - security different than functional correctness, because edge case could be very niche and that is very bad (can we make this argument), hard to randomly get to instantiation of one side channel but once the side channel is found it is infinitely exploitable
% - people have done this with randomized test generation, but it does not target potentially violating tests
% - people have benchmarks of tests that trigger spectre attacks, but they are not specialized to the 
Benchmarks of tests that trigger certain Spectre attacks have been compiled~\cite{kocher2018msvc,canella2019systematic}%
, but they are not specialized to the security guarantees of a particular defense, as different defenses enforce different rules on what hardware can leak, and cover a small subset of all relevant leakage scenarios. 
One piece of prior work — AMuLeT, has tested simulator implementations of Spectre defenses. AMuLeT uses randomized differentiated testing (running two randomly generated architecturally identical tests and checking for discrepancies in a limited portion of microarchitectural state) \cite{amulet}. 
\urgentbr
\urgentnote{The parenthetical misdescribes AMuLeT: it runs one random program on two inputs whose contract traces match and flags a difference in the microarchitectural trace (its Def.~2.1, fu-2025-amulet.pdf \S2.3), not ``two \ldots{} architecturally identical tests''. \textbf{PROPOSED TEXT} (edits the sentence just before this note): AMuLeT uses random relational testing: it runs a randomly generated program on pairs of inputs whose contract traces match and checks for differences in a limited part of the microarchitectural state~\cite{amulet}.}
\notebr
\anna{I think maybe drop the security property vs leakage predicate discussion here, and pick that up in related work. the main point being that it is random}.
% AMuLeT checks that a simulated defense satisfies a hardware--software contract~\cite{guarnieri2021contracts}: any two runs of a program whose contract traces match must also produce matching microarchitectural traces~\cite{amulet}. Such contracts generalize speculative non-interference (SNI)~\cite{guarnieri2020spectector}, which requires that speculatively executed instructions leak no more into microarchitectural state than non-speculative execution does. SNI and contracts are end-to-end, relational properties: they compare pairs of whole executions and do not say which instruction leaked or which rule of the defense it broke, so every violation must be root-caused by hand.
\anna{this paragraph abt fuzzing vs targeting? also in related work?}
% It also does not bias for generation of programs that are able to induce leakage. Security testing is sensitive, because just one leaky edge case that might be very difficult to randomly instantiate could be a potent attack vector, so it is important to generate a test suite that has a wide breadth of leaky instances. There exists no existing work that can generate targeted test sets for Spectre defense validation in a principled manner. 


In this paper we propose \toolname, the first rigorous, targeted, and automated approach for validating early-stage prototypes of hardware support for Spectre defenses against their promised security guarantees. \toolname provides (1) a \textbf{formal language} for specifying security guarantees of an arbitrary Spectre defense, (2) a \textbf{test generator} for generating a comprehensive set of tests (given a bound on test size) capturing all leakage scenarios blocked by a specified defense, and (3) a \textbf{test validator} harness for running these tests on gem5 implementations of the defense to detect bugs. \anna{stt running mini example figure 1?? bring the figure upwards otherwise}
\input{tikz/intro_threelayer_v2_8pt}% Fig. 1 right after the proposal (your note: bring the figure upwards), 2026-11-03

% \toolname enables validation of the hardware implementation Spectre defense against a defense security guarantee, and the specification of HSCAs against which it protects. This specification is precise to the cases blocked by a defense, and can be used to generate test sets; using it to prove adherence to a higher-level security property is future work.
% \notebr
\reviewnote{\textbf{ADD} ``(Fig.~\ref{fig:intro})'' to the sentence above: Fig.~1 is never referenced in the text, and this sentence is what it shows.}

\toolname{} tests defense implementations to their guarantees at the level of the defense's \emph{leakage predicate}: a condition on a single execution that names the transmitter, the protected operand it may not expose, and the speculative microarchitectural state in which the defense forbids the transmission. 

This means that an arbitrary natural language security description can be translated into formal components, which can then be used to generate all possible violating sequences of (abstract) instructions. These sequences are then used to generate a litmus test suite of assembly programs and corresponding \textit{microarchitectural annotations} that together induce the violating sequences. Finally, \toolname provides a harness that instruments the simulator to follow these speculation annotations, inducing the modeled violations.

The test sets generated by \toolname are complete in including all  \textit{leakage backbones}, or fundamental leakage structures, up to an instruction bound, and are diversified across timing, instruction type, and speculation direction during their concretization to assembly programs. 
\urgentbr
\urgentnote{Completeness holds only in part. Stage~1a (minimal graphs, up to six instructions) runs to UNSAT for STT and ReCon and their store-bypass variants, but SPT's Stage~1a does not converge (its four-instruction partition hits the 200,000 cap), and Stage~1b and Stage~2 are capped or sampled by design. \textbf{PROPOSED TEXT} (edits the sentence just before this note): For STT and ReCon, the test sets include every \textit{leakage backbone} up to six instructions (for SPT, up to a cap), and are diversified \ldots{} \textit{Source:} simspect wiki/state/EVAL\_MANIFEST.md `Stage 1' and `Stage 1.5' sections.}
\notebr
The harness that they run on both enforces the microarchitectural constraints of the test execution, and induces microarchitectural behavior that is likely to precipitate the test. 

In summary, \toolname provides not only a simple formal language in which to define a defense, but also uses it to target \textit{all possible} types of violations stopped by the defense, probing leakage in a principled and efficient manner and providing better assurance of design correctness.

We have expressed 10 hardware defenses and their variants in the \toolname{} language; STT, SpecShield, NDA, DoM, SDO, SPT, ReCon, InvisiSpec, SpecLFB, and CleanupSpec. We can therefore generate comprehensive litmus test suites up to specified bounds for each defense. 
\urgentbr
\urgentnote{Not supported as written. Only STT, ReCon and SPT exist as Alloy models (STT and ReCon also in store-bypass variants); InvisiSpec, CleanupSpec and SpecLFB are tested with STT's test set and a leakage criterion written from each defense's paper; SpecShield, NDA, DoM and SDO have no model. \textbf{PROPOSED TEXT} (edits the sentence just before this note): We have expressed three hardware defenses in the \toolname{} language, STT, ReCon and SPT, and test three more, InvisiSpec, CleanupSpec and SpecLFB, with STT's test set and a defense-specific leakage criterion. \textit{Source:} simspect STAGE1\_alloy/models/ (STT\_6, Recon\_6, SPT\_6\_oneLP\_v2 and their proposed/ultimate versions); specs/leakage/\{invisispec,cleanupspec,speclfb\}.json.}
\notebr

% We have run the \toolname{} testing pipeline with an instruction bound of six on the gem5 implementations of STT, SPT, and ReCon, and on the AMuLeT builds of InvisiSpec, SpecLFB, and CleanupSpec

We apply \toolname{}, with an instruction bound of six, to every simulator prototype of a state-of-the-art defense we could obtain: the gem5 implementations of STT, SPT, and ReCon, and the AMuLeT builds of InvisiSpec, CleanupSpec, and SpecLFB. Checking each defense against only the leakage its own paper claims to prevent, \toolname{} replicates every previously known bug in that scope: InvisiSpec's speculative L1D eviction, CleanupSpec's three cleanup bugs (an uncleaned speculative read-modify-write, an uncleaned line-crossing access, and cleaning too much), and SpecLFB's unprotected first load.
\notebr
\reviewnote{Evidence behind this sentence, for the counts you add later: AMuLeT reports nine bugs on these defenses; five fall inside each defense's own guarantee and SimSpect finds all five: UV1 (InvisiSpec, 2,995 of 3,000 tests), UV3 (only with the read-modify-write corpus and a cold cache: 407 of 590 tests cold, 0 of 590 warm), UV4 (38 tests), UV5 (one test so far, inst-039102; whether to count it is still your open decision) and UV6 (about 1,140 of 1,500 tests per mode). Out of scope, each by its paper: UV2 (MSHR contention, InvisiSpec \S VII), KV1 (instruction fetch, InvisiSpec fn.~2), KV2 (unXpec cleanup timing: CleanupSpec \S4(b) p8 acknowledges that cleanup time varies and leaves constant-time cleanup to future work, and AMuLeT observes it through the L1I, which CleanupSpec \S2.1 p3 excludes) and KV3 (speculative store fills the D-TLB: STT assumes stores act only at retirement, \S5 p5). \textit{Source:} simspect wiki/findings/mdsToRead/AMULET\_BUGS\_DETECTION\_ACCOUNTING.md sec.~2; URGENT\_TODO.md (KV2, KV3, UV3, UV5 items).}
\notebr

\toolname{} also finds six previously unknown bugs. (1) STT fails to taint one register: a check that skips register index 16 ignores the register class, so a secret copied through the low half of \texttt{xmm4} loses its taint. In ReCon's underlying implementation of STT, (2) a load that receives its value by store-to-load forwarding completes before the taint check runs, and (3) an instruction that merges two tainted operands keeps only the oldest taint source, so its result becomes untainted once that source's branch resolves. SpecLFB (4) misimplements its ROB walk order, so loads after the first can also go unprotected, and (5) fails to protect loads in the shadow of a direct \texttt{jmp} or \texttt{call}, which erases an older unresolved branch's shadow. Finally, (6) STT and SPT share a visibility-point/squash timing race: a wrong-path transmitter is released when its branch resolves and accesses the cache before the squash removes it.
\urgentbr
\urgentnote{(6) leaves out the condition the race needs: it leaks only when commit's squash reaches the issue stage at least two to three cycles late (gem5's commitToIEWDelay), and never at stock timing. SPT Fence on 24 race tests: stock 0, delay 2 $\to$ 3, delay 3 $\to$ 14; STT leaks at delay 3 and 0 of 1,200 at stock. This is a global latency setting, not one of the test hooks, so say so. \textbf{PROPOSED TEXT} (edits the last sentence before this note): Finally, (6) STT and SPT share a visibility-point/squash timing race: when the squash signal reaches the issue stage a few cycles after commit (a longer commit-to-IEW delay than gem5's default), a wrong-path transmitter is released when its branch resolves and accesses the cache before the squash removes it. Also, (5) is reached by the generated tests' own jumps, not only by the harness's call. \textit{Source:} simspect wiki/findings/bug/BUG\_TAXONOMY.md A4, A5; project log 2026-11-09 15:20 (bug-investigate-a5) and 20:10 (bug-investigate-e14new).}
\notebr

Compared to prior work, \toolname{} needs [N]$\times$ fewer tests on average to find the known bugs. Every test it generates contains a protected transmitter, while only 17\% of AMuLeT's randomly synthesized tests contain the components necessary for leakage.

% NO

% Because the leakage predicate is formal, it can be related to the end-to-end properties. A bounded Alloy check that every pair of executions allowed by a defense's leakage predicate satisfies SNI, or a given contract, would turn passing \toolname{} tests into evidence for that end-to-end claim, such as STT's non-interference theorem~\cite{yu2019stt}.
\section{Background}
\subsection{Hardware-side-channel defenses}\label{sec:bg-defenses}
Many HSCDs exist that include some hardware component, as viable performance necessitates HSCAs to be mitigated at least partially in hardware. These defenses aim to protect some set of architectural state from being leaked through some side channel formed by a data-dependent optimization. Many of these defenses function by designating a set of architectural state that contains secrets, and tracking how information flows from an secret piece of state to a speculative side channel. STT, or Speculative Taint Tracking, is a defense that stalls transmitters, or instructions that create data-dependent behavior, until they are non-speculative when they depend on a tainted, or secret, piece of state \cite{yu2019stt}.
\notebr
\auditnote{cite}{proposed}{\textbf{ADD (new sentence, SPT; insert here, before ReCon)}: ``SPT~\cite{spt} delays every transmitter until it can prove that the transmitter's operands are leaked by the program's non-speculative execution anyway, using a dynamic information-flow analysis to declassify data as it leaks.'' \newline \textbf{REWORD (optional, the STT sentence above)}: end it with ``\ldots until the instruction that produced the tainted data becomes non-speculative, i.e., reaches its visibility point~\cite{yu2019stt}'', which sets up the visibility-point race in the results. \newline \textit{Source:} choudhary-2021-spt.pdf abstract (``delays execution of every transmitter until it can prove that the transmitter's operands leak during the program's non-speculative execution''); yu-2019-stt.pdf abstract (``the earliest time is when the instruction(s) producing the protected data become non-speculative'').}
\notebr
ReCon builds on such defenses by shrinking what they protect: it detects data that has already leaked during non-speculative execution, limiting detection to pairs of directly dependent loads that dereference a pointer, records that knowledge in the coherence protocol, and lets a defense such as STT treat that data as safe to use under speculation \cite{recon}.
\notebr
\reviewnote{Anna, check this ReCon sentence: it replaces the NDA sentence (as you asked) and is paraphrased from ReCon's abstract in references.bib.}
\notebr
Many more hardware HSCDs exist, some of which improve the performance of STT by actively eliminating items from the STT protection set. All of these defenses protect against different combinations of leakage components, including different protection sets, and different types of side channels.
\notebr
\auditnote{cite}{proposed}{\textbf{ADD (new sentences; insert after this paragraph, or in place of its last sentence)}: ``Other defenses let speculative loads execute and instead control their effect on the cache: InvisiSpec~\cite{yan2018invisispec} keeps speculative loads invisible to the cache hierarchy and exposes each load only once it is safe, CleanupSpec~\cite{saileshwar2019cleanupspec} lets speculative loads change the cache and undoes those changes on a mis-speculation, and SpecLFB~\cite{cheng2024speclfb} holds speculative misses in the cache's line-fill buffer and installs them only once they are safe. All of these defenses target Spectre attacks~\cite{spectre}, including store bypass (Spectre v4), in which a load speculatively reads stale data before an older store's address is known~\cite{horn2018ssb}.'' \newline \textit{Source:} fu-2025-amulet.pdf \S4.5 (InvisiSpec: ``speculative loads are invisible to the caches; when loads become safe, an expose operation makes the load visible''), \S4.6 (CleanupSpec: ``allows speculative loads to modify cache state and cleans up these state changes on a mis-speculation''), \S4.7 (SpecLFB: ``adds security checks to the cache Line-Fill Buffer\ldots blocks speculative cache misses from being installed into the cache until they are safe''); the store-bypass wording follows your own pipeline.tex description (check it against horn2018ssb). \newline \textbf{Two caveats (2026-11-09):} keep the SpecLFB clause as a description of the paper's design only: the gem5 SpecLFB build bounces unsafe misses and replays them, and its park-in-the-LFB path is unreachable, so do not carry it into the Results. Our STT tests are branch-only and the store-bypass tests run only on the Kaxiras-group builds, so write ``All of these designs target \ldots'' rather than implying that every defense is tested for store bypass. \textit{Source:} simspect URGENT\_TODO.md (``SpecLFB parks unsafe misses in the LFB'' correction); wiki/state/EVAL\_MANIFEST.md STL section.}
\urgentbr
\urgentnote{\textbf{VOCAB}: \textbf{ADD (new paragraph: leakage contracts)} here: this is where you separate our terms from Guarnieri et al.'s and AMuLeT's. \textbf{PROPOSED TEXT}: ``Guarnieri et al.\ specify what a CPU may leak as a \emph{hardware--software contract}~\cite{guarnieri2021contracts}: an \emph{observation clause} says what each executed instruction exposes to an attacker (for example, the addresses of loads and stores), and an \emph{execution clause} says which instructions may execute speculatively. Revizor and AMuLeT test CPUs and simulators against such contracts~\cite{revizor,amulet}, which Revizor calls speculation contracts and AMuLeT calls leakage contracts. We use \emph{leakage contract} in the narrower sense of SynthLC~\cite{rtl2upath}: the set of (transmitter, operand) pairs through which a design can leak. Our leakage contract is therefore closest to their observation clause, though it names the operands that are fundamentally leaked rather than the attacker-visible state; our execution contract is closest to their execution clause; and our protection set makes explicit which data must not reach those operands, which a contract leaves implicit (everything its traces do not expose).'' Check the last clause: it is an interpretation of contract satisfaction, not a statement from the papers. Also cite here: Spectector~\cite{guarnieri2020spectector} for SNI, LCM~\cite{lcms}, and the RISC-V follow-ups~\cite{mohr2024synthesizing,synthlc,leave}. \textit{Source:} fu-2025-amulet.pdf Sec.~2.3 (``A leakage contract captures the expected leakage of a CPU under test'') and Table~1 (CT-SEQ's observation and execution clauses); Revizor's title; RTL2M$\mu$PATH (MICRO 2024) on SynthLC's ``formally verified leakage contracts''.}
\notebr

\subsection{HSCD litmus testing}\label{sec:bg-litmus}
Prior work that validates advertised security guarantees against hardware support for hardware SCDs analyzes fabricated designs\cite{revizor} or RTL, due to the wide availability\cite{rtl2upath}.
\notebr
\auditnote{cite}{cite, high}{The sentence before this note still ends in the unfinished clause `due to the wide availability', and Revizor tests commercial black-box CPUs against speculation contracts, not a defense's advertised guarantee. The intro's prior-work sentence (``Existing validation either tests fabricated CPUs\ldots'') now says this with the full cite list, so point back to it or reuse it here. \newline \textbf{PROPOSED TEXT} (edits the sentence just before this note): Prior work validates side-channel security on fabricated CPUs~\cite{revizor,oleksenko2023hideseekspectresefficient} and on RTL designs~\cite{rtl2upath,contractshadow,hur2022specdoctor}. \newline \textit{Source:} references.bib entries revizor, oleksenko2023hideseekspectresefficient (S\&P 2023), rtl2upath (MICRO 2024), contractshadow (ASPLOS 2025), hur2022specdoctor (CCS 2022); sections/intro.tex prior-work sentence.}
\notebr
Only AMuLeT\cite{amulet} analyzes simulator implementations. AMuLeT is a fuzzer that generates validation tests to detect security bugs in simulated Spectre defense prototypes—one of our goals. It does random test generation and conducts relational testing by running programs in simulation that, when supplied with two different inputs, are indistinguishable in their architectural side-channel leakage, flagging differentiated microarchitectural state. However, it attempts to validate high-level security properties, and has no ability to do directed testing for leakage.

\anna{talk abt uarch state}

\subsection{Axiomatic execution models}\label{sec:bg-axiomatic}
Prior work has formalized processor memory consistency models (MCMs) into both axiomatic and operational models \cite{herdingcats}\anna{cite a bunch of other stuff, arm + simner stuff}. Axiomatic MCM models represent an execution as a graph of memory events connected by relations such as program order (po), which relates an instruction to its subsequent instruction in thread-order, reads-from (rf) which relates a write to a read that uses its value, and coherence (co), which relates two writes to the same address. Lustig et. al. use this format of model along with a set of minimality requirements to generate a comprehensive set of MCM-violating execution graphs and concretize them into litmus tests for testing an MCM implementation \cite{lustig}. They do this by using the Alloy Analyzer \cite{jackson2002alloy} to generate all tests that fail according to their MCM and are minimal
\notebr
\auditnote{cite}{cite, high}{The sentence before this note has no final period, and Fig.~3 names the Kodkod model finder, which has no bib entry yet (jackson2002alloy is now cited here). \newline \textbf{PROPOSED TEXT} (edits the sentence just before this note): They use the Alloy Analyzer~\cite{jackson2002alloy,kodkod} to generate every minimal test that their MCM forbids. New bib entry \texttt{kodkod}: E.~Torlak and D.~Jackson, ``Kodkod: A relational model finder,'' TACAS 2007 (as cited in lustig-2017-automated-litmus-synthesis.pdf and trippel-2018-checkmate.pdf ref.~[33]).}
\notebr
\anna{replace this sentence, cite wickerson}. Minimality, as defined in the test generation procedure, means that a graph violates the desired predicate, and cannot be relaxed, i.e. have a component removed or demoted, without becoming a non-violating graph. The relaxations considered by Lustig et. al. include removing instructions and dependencies, as well as demoting fences. Minimality allows prioritization of axes for test diversity.

Axiomatic modeling has been used for formalization of speculation restrictions \anna{cite checkmate, all of the speculation formalization, LCM}. 

\anna{add way more about modeling executions}
% A similar such axiomatic definition of side channel leakage would allow generation of a similar test suite satisfying a minimality criterion. 
\notebr
\auditnote{cite}{proposed}{\textbf{WRITE MORE HERE (axiomatic modeling)}: write this out more here (you said you would). Proposed thread and citations: memory models as relations over executions~\cite{herdingcats,alglave2010fences,batty2011mathematizing,sarkar2011power,pulte2018simplifying,lahav2017repairing}; Alloy and relational model finding~\cite{jackson2002alloy}; comparing and synthesizing models and litmus tests~\cite{memalloy,lustig,memsynth,alglave2011litmus,mcmutants}; microarchitectural happens-before ($\mu$hb) models~\cite{lustig2014pipecheck,manerkar2015ccicheck,lustig2016coatcheck,trippel2017tricheck,lustig2015armor,hossain2020transform}; full ISA and system semantics~\cite{simner2020ifetch,simner2022vm,simner2025exceptions,armstrong2021isla,armstrong2019sail}; and the step to security: CheckMate's $\mu$hb exploit patterns~\cite{checkmate}, Cats vs.\ Spectre~\cite{ponce2022cats} and LCM~\cite{lcms}. End with what SimSpect takes from this line of work (execution graphs with rf/spo-style relations, minimal litmus enumeration) and what it adds (speculation flags and a defense's leakage predicate). Since 2026-11-08 the enumeration has two stages, so say ``minimal litmus enumeration (Stage~1a, RI-minimal graphs), then a capped flesh-out under loosened minimality (Stage~1b)'' rather than only ``minimal litmus enumeration''.}
\notebr


\section{Formal model of defense leakage} \label{formalmodel}
We adopt a formal language for specifying the security guarantees of a \anna{comprehensive??} Spectre defense.
% , that allows us to generate guarantee-violating execution graphs and concretize them into litmus tests to do targeted testing on an implementation. 
We take inspiration from prior work to 
We survey the literature of hardware defenses to Spectre attacks and conclude that a formal Spectre defense specification should consist of three components: an execution contract (what is speculative), a protection set (what is protected and how protection propagates), and a leakage contract (which operands can transmit).
\begin{enumerate}
    \item \textbf{Leakage contract:}
A leakage contract is the set of pairs of \textit{transmitters} (instructions that leak one or more operands via a hardware side channel) and their unsafe operand(s), as recognized by the defense.
\item \textbf{Execution contract:} 
an execution contract characterizes what speculative control- and data-flow is allowed in the $\mu$arch by identifying which \textit{speculation primitives} (instructions that initiate speculative executions, e.g., a conditional branch) the defense considers. 
\item \textbf{Protection set:}
A protection set is the set of architectural state (memory and registers) that the defense promises to not leak architecturally and/or speculatively. Along with the protection set we encode a protection set update rule, that represents the updates that a given protection set undergoes with each instruction.
\end{enumerate}
We then construct a per-defense \textit{leakage predicate}.
A defense's \emph{leakage predicate} combines its three components: an execution violates it when data in the protection set reaches, through rf and ddi edges, an operand in the leakage contract of an instruction that is speculative under the execution contract.
% \textbf{1. Leakage contract}
% A leakage contract is the set of pairs of \textit{transmitters} (instructions that leak one or more operands via a hardware side channel) and their unsafe operand(s), as recognized by the defense.
% \textbf{2. Execution contract}
% An execution contract characterizes what speculative control- and data-flow is allowed on the microarchitecture by identifying what \textit{speculation primitives} (instructions that initiate speculative executions, e.g., a conditional branch) the defense considers. 
% \textbf{3. Protection set}
% A protection set is the set of architectural state (memory and registers) that the defense promises to not leak architecturally and/or speculatively. Along with the protection set we encode a protection set update rule, that represents the updates that a given protection set undergoes with each instruction.
% todo these components sufficient for now etc etc 

\subsection{Execution Syntax}
\input{tikz/execvocab_v2_8pt}

\auditnote{fact}{fact}{\textbf{About the caption of Table~\ref{tab:execvocab}:} The vocabulary is shared, but more than the defense block differs. Recon\_6 and SPT\_6\_oneLP lack rf\_from\_most\_recent\_writer. SPT\_6\_oneLP rewrites skeleton code: operand-typed a and last\_committed\_protset\_p, secure\_speculation\_scheme\_p, tag\_xm, and an exactly-one-leak-point run predicate. Recon\_6 changes a skeleton call site to hardware\_protection\_policy[p]. This is open as a DECIDE in URGENT\_TODO.md:31.  \textbf{PROPOSED TEXT} (edits the sentence just before this note): The defense block (speculation contract, protection set, leakage function) is what defines each defense; the models also differ in a few skeleton facts (Section\~{}\textbackslash{}ref\{...\}).  \textit{Source:} diff STT\_6.als vs Recon\_6.als (rf fact; :229) and vs SPT\_6\_oneLP.als (:220-284, 313-314); URGENT\_TODO.md:31}
\notebr
\auditnote{fact}{fact}{\textbf{Update to the note above (2026-11-09):} two of its differences are gone. Every in-use model now has rf\_from\_most\_recent\_writer, and the production SPT model (SPT\_6\_oneLP\_v2, approved 2026-11-08) asks for an ordinary violation (at least one leaking operand). Still different outside the defense block: SPT tracks its protection set per operand, not per state (a, last\_committed\_protset\_p, secure\_speculation\_scheme\_p, tag\_xm), and Recon\_6 passes the perturbation into hardware\_protection\_policy. The ultimate models (STT\_6\_u approved for the clean runs on 2026-11-09) put every defense block on one shared skeleton with per-operand protection; once they are the evaluated models, the shared-skeleton claim is exact. \newline \textbf{PROPOSED TEXT} (edits the sentence before the note above): The defense block (speculation contract, protection set, leakage function) defines each defense; the rest of each model is shared, except that SPT tracks its protection set per operand. \newline \textit{Source:} STAGE1\_alloy/models/Recon\_6.als:163, :239; SPT\_6\_oneLP\_v2.als:6-17, :184, :254-256, :346; STAGE1\_alloy/models/proposed/ultimate/SKELETON\_u.als; URGENT\_TODO.md decision U1}
\notebr
To express a leakage predicate distinguishing leaky from non-leaky executions (e.g., the load-transmitter leakage in Figure~\ref{fig:ddi-rf}), we model program executions as directed happens-before graphs, inspired by MCM analyses~\cite{herdingcats,lustig}\anna{add citations}, and define leakage as a predicate over them, as in \ref{formalmodel}.
% \figtodo{cite the basic execution-graph example, Fig.~\ref{fig:ddi-rf} (the dataflow figure, redrawn in the new notation), here; the violation definition below says I3, the figure now says i3}
Generated tests consist of a total program order
\urgentbr
\urgentnote{\textbf{VOCAB}: \textbf{ADD DEFINITION} (new text, insert here): \emph{spo}, the speculative program order, the total order on instructions.}
\notebr
of \textbf{Instruction} nodes, typed as transmitters (Load, Store, Branchx, Otherx) or non-transmitters (Branchn, Othern). Input/output \textbf{Operands} (inmem, inreg, inaddr, outreg, outmem) refer to \textbf{State} objects (Reg, Mem). Instructions carry $\mu$arch flags for committed/resolved status
\urgentbr
\urgentnote{\textbf{VOCAB}: \textbf{ADD DEFINITION} (new text, insert here): \emph{resolved} (a branch's outcome, or with store bypass a store's address, is known) and \emph{committed} (retired; implies resolved; a prefix of spo).}%
\urgentbr
\urgentnote{\textbf{VOCAB}, replaces the definition just above: \emph{resolved} (the defense knows a branch's outcome, or with store bypass a store's address; a branch that has executed but not yet squashed its wrong path still counts as unresolved, as decided 2026-11-08) and \emph{committed} (retired; implies resolved; a prefix of spo).}%
\notebr
, enforced in simulation to replicate execution.
\notebr
\auditnote{fact}{fact, low}{Only resolution is actively enforced, by the branch force, resolve stall and STL hold. Commit status is measured, not enforced. The evaluated fixed configs set the fnc commit stall to 0, and the checkers only require lc to retire and the signal to come before fnc retires. The Gem5 section says this correctly ('passively measuring'). \newline \textbf{PROPOSED TEXT} (edits the sentence just before this note): Instructions carry \$\textbackslash{}mu\$arch flags for committed/resolved status; resolution is enforced in simulation and commit status is checked against the trace. \newline \textit{Source:} runconfigs/fixed\_20261030/*.jsonc (fnc\_commit\_stall\_cycles 0); STAGE3\_gem5/check\_ld.py:95-114}
\notebr
Data-dependence edges (rf, ddi)
\urgentbr
\urgentnote{\textbf{VOCAB}: \textbf{ADD DEFINITION} (new text, insert here): \emph{rf} (reads-from: an output operand to an input operand of a later instruction on the same location) and \emph{ddi} (data-dependence internal: an input to an output of the same instruction).}
\notebr
track information flow to transmitters. First-order axioms constrain legal edge placements. We choose the instructions, operands, and state that we use because we find it to be the minimal set of state needed to represent the distinctions in security guarantees in the body of Spectre defenses we study. Figure \ref{tab:execvocab} contains the execution graph syntax.
\notebr
\auditnote{fig}{figure, low}{This prints 'Figure I'; the target is a table. \newline \textbf{PROPOSED TEXT} (edits the sentence just before this note): Table\~{}\textbackslash{}ref\{tab:execvocab\} contains the execution-graph syntax. \newline \textit{Source:} p.4 left column, last line of the Execution Syntax paragraph.}
\notebr
% \figtodo{new vocabulary table, Table~\ref{tab:execvocab}: reference it here; the prose above can shrink to one sentence (suggested text in simspect paper\_figures/FIGURES.md)}

A \textit{violation} is defined as
\notebr
\textit{a data-dependence path from a protection-set state element as updated by the most recently committed instruction,
\auditnote{fact}{fact, high}{No model updates the protection set only up to the last committed instruction. last\_committed\_protset\_p applies the update rule to every instruction in spo order, committed or not. STT's rule removes the memory read by any load with no older unresolved branch. The commit-based rule is commented out. The dataflow figure caption already states the resolution-based rule, so the text and that caption disagree.  \textbf{PROPOSED TEXT} (edits the sentence just before this note): a data-dependence path from a protection-set state element (after the defense's update rule has been applied along speculative program order) to a leaked operand of a transmitter, ...  \textit{Source:} STAGE1\_alloy/models/STT\_6.als:232-243 (fold over f..f.spo\^{}6), :329-331 (rule), :330 (commit rule commented out); tikz/dataflow\_8pt.tex caption}
\auditnote{fact}{fact, high}{\textbf{Update (2026-11-09):} the ultimate STT model, STT\_6\_u (STT\_6\_v2 semantics), was approved for the clean runs. It taints the value read by every unsafe load, forwards taint along rf and ddi, and untaints only by the speculation contract, with a per-operand protection set. If the paper's results come from it, define a violation as a data-dependence path from a value read by a speculative load to a leaked operand of a speculative transmitter, and drop ``protection-set state element''. \textit{Source:} URGENT\_TODO.md decision U1 and the STT\_6\_u DECIDE; STAGE1\_alloy/models/proposed/ultimate/STT\_6\_u.als}
to a leaked operand of a transmitter, where the transmitter is in the leakage contract and speculative as per the execution contract.} In Figure~\ref{fig:ddi-rf}, the transmitter load in \texttt{I3} has an address dependency on the load in \texttt{I1}
\notebr
\auditnote{fig}{figure, low}{The text names the instructions I1, I2, I3 in typewriter, but Fig. 2 now labels them i1, i2, i3. \newline \textbf{PROPOSED TEXT} (edits the sentence just before this note): ... the transmitter load i3 has an address dependency on the load i1 through the ALU op i2 ... \newline \textit{Source:} p.4 text vs Fig. 2 row labels.}
\notebr
through \texttt{I2} arithmetic, and the memory is in the protection set, creating a violation. 

\input{tikz/dataflow_8pt}% redrawn basic example; old tikz/alloyviolation kept, 2026-11-02
\input{tikz/pipeline_full_8pt1col}
\notebr
\auditnote{fig}{figure}{\textbf{About the caption of Fig.~\ref{fig:pipeline}:} 'Non-invasive' (also in the caption) overstates the harness: the directives change prediction and timing on purpose (BTB force, resolve stall, fnc stall). The 'selects hooks' arrow from transmission mechanisms into the harness has no code behind it.  \textbf{PROPOSED TEXT} (edits the sentence just before this note): "+ Minimal harness (defense logic untouched): probes; speculation and timing directives"; drop the 'selects hooks' arrow, or label it 'design time'.  \textit{Source:} BIGFIG\_FACT\_AUDIT.md rows 30 and 39, decisions D6 and D8; pipeline.tex gem5-hooks paragraph describes the stalls as active control hooks.}
\notebr
\auditnote{fig}{figure}{\textbf{Update to the note above (2026-11-09):} the `selects hooks' arrow stays (decided 2026-11-02), but the caption or text must say the hooks are chosen at design time: no code passes the transmission mechanism to gem5 at run time. `Non-invasive' still overstates the harness: the branch force, the resolve stall and the store-address hold change prediction and timing on purpose, and the CleanupSpec and SpecLFB runs start each test with a hot or cold cache at random (decided 2026-11-09); only the defense logic is untouched (wording still open). \newline \textbf{PROPOSED TEXT}: ``+ Minimal harness (defense logic untouched): probes; speculation and timing directives''; in the caption, ``\ldots{} under a minimal harness whose hooks are chosen at design time from the defense's transmission mechanisms''. \newline \textit{Source:} URGENT\_TODO.md (pipeline figure decisions; cache-setup decision 2026-11-09); wiki/design/RESOLUTION\_HOOKS.md sec.~1; runconfigs/twostage\_fv2\_20261109/E13\_STT\_6\_fv2.jsonc, E14 (cache\_setup random)}
\notebr
\auditnote{fig}{figure, high}{\textbf{About the caption of Fig.~\ref{fig:pipeline}:} No evaluated enumeration reached UNSAT; every testset stops at max\_instances (100k), and Sec. V's OTB paragraph says the shape is 'not reached within its instance cap'. Fig. 3 calling the output 'bounded-exhaustive' contradicts the paper's own results.  \textbf{PROPOSED TEXT} (edits the sentence just before this note): "Minimal leaking execution graphs, up to the bound and instance cap" (or keep 'bounded-exhaustive' only if a run reaches UNSAT; decision D2).  \textit{Source:} BIGFIG\_FACT\_AUDIT.md rows 7, 9 and D2; STAGE1\_alloy/models/MANIFEST.md (100k cap); results.tex sec:bug-otb 'does not reach this shape within its instance cap'.}
\notebr
\auditnote{fig}{figure, high}{\textbf{Update to the note above (2026-11-09):} `Bounded-exhaustive' now holds in part. In the two-stage runs, Stage~1a (up to six instructions) reaches UNSAT for STT and ReCon and their store-bypass variants, but SPT's Stage~1a does not converge: its four-instruction partition is capped at 200,000 graphs. The June numbers come from flat test sets that stop at their instance cap. \newline \textbf{PROPOSED TEXT}: ``Bounded-exhaustive (STT, ReCon) set of minimal violating execution graphs'', or keep the label once SPT's Stage~1a is settled. \newline \textit{Source:} wiki/state/EVAL\_MANIFEST.md `Stage 1' section}
\notebr
\auditnote{fig}{figure, high}{\textbf{About the caption of Fig.~\ref{fig:pipeline}:} The caption calls Stage 1b's output minimal, but the figure itself says Stage 1b loosens minimality (slab 2'). Also, per the fact audit, Stage 1b is an STT-only prototype that produced none of the reported results, and neither the figure nor the caption says so.  \textbf{PROPOSED TEXT} (edits the sentence just before this note): Stage\~{}1 uses Alloy to enumerate minimal leaking execution graphs up to a bound; Stage\~{}1.5 (prototype, not used for the results below) re-enumerates with loosened minimality, pinned to each Stage\~{}1 graph.  \textit{Source:} Fig. 3 slab '2'. Loosened minimality requirements'; BIGFIG\_FACT\_AUDIT.md row 10 and decision D1 (no campaign used Stage 1b).}
\notebr
\auditnote{fig}{figure, high}{\textbf{Update to the note above (2026-11-09):} Stage~1b is no longer a prototype. It runs one pinned Alloy run per Stage-1a graph, in parallel, at exactly six instructions, keeps RR, RC and RS but not RI, and caps each graph's flesh-out (about 300,000 graphs per defense); the clean-run results will come from it, no number in the paper does yet. The caption still calls its output minimal. \newline \textbf{PROPOSED TEXT}: Stage~1a uses Alloy to enumerate minimal violating execution graphs of up to six instructions; Stage~1b grows each one into six-instruction graphs under loosened minimality, a capped number per Stage-1a graph. \newline \textit{Source:} wiki/state/EVAL\_MANIFEST.md `Stage 1.5' section; STAGE1\_alloy/models/proposed\_twophase/STT\_6\_core\_ri.als:289-310}
\notebr
\auditnote{fact}{fact}{\textbf{About the caption of Fig.~\ref{fig:pipeline}:} The harness changes prediction and timing on purpose. It forces branch directions, holds branch resolution and store addresses, and on the resolve-hook builds wakes the CPU at release. Only the defense logic is untouched. The same word appears in the figure body and in pipeline.tex ('non-invasively instrument'). The wording is an open item in URGENT\_TODO.md:44.  \textbf{PROPOSED TEXT} (edits the sentence just before this note): Stage\~{}3 runs them on the defense's gem5 build under a minimal harness that leaves the defense logic untouched  \textit{Source:} wiki/design/RESOLUTION\_HOOKS.md sec.1 (behaviour changes; edit 4); gem5-hooks/README.md (STL hooks); paper\_figures/plans/BIGFIG\_FACT\_AUDIT.md row 30; URGENT\_TODO.md:44}
\notebr
\auditnote{fig}{figure, low}{\textbf{About the caption of Fig.~\ref{fig:pipeline}:} The naming is inconsistent: Fig. 3 says 'Inspector' and 'Pipeline info artifact', while the text and the Fig. 7 caption say 'checker', 'per-transmitter pipeline predicate' and 'gem5 execution trace'. Stage 1a is 'Abstract test generation' and Stage 1b 'Abstract test enumeration', which implies a difference that does not exist.  \textbf{PROPOSED TEXT} (edits the sentence just before this note): Use one name for the checker across Fig. 3, Fig. 7 and the text, and 'Pipeline trace' for the artifact; e.g. Stage 1a 'Minimal-graph enumeration', Stage 1b 'Graph expansion'.  \textit{Source:} Fig. 3 vs Fig. 7 caption ('How the checker decides') and pipeline.tex Transmission Requirement; BIGFIG\_FACT\_AUDIT.md rows 1, 35.}
\notebr
\auditnote{fig}{figure}{\textbf{About the caption of Fig.~\ref{fig:pipeline}:} The caption explains only blue. It never explains the tan, lavender and light-red slabs 3a-3c (which carry the same meaning as in the graph figures), the dashed 'Same as Stage 1a' ghost slabs, the + and prime marks, the numbered chevrons, or the dashed 'SimSpect goal' arrow.  \textbf{PROPOSED TEXT} (edits the sentence just before this note): Blue marks defense-specific components; slabs 3a-3c use the speculation (tan), protection-set (lavender) and transmitter (red) colors of the graph figures; dashed slabs are unchanged from Stage 1a, + marks an added and ' a changed constraint.  \textit{Source:} p.4 Fig. 3 render; VISUAL\_LANGUAGE.md CURRENT RULES (pipeline 3a tan, 3b lavender, 3c light red).}
\notebr

\subsection{Example contract}
\notebr
\urgentbr\urgentnote{\textbf{RENAME (Anna, 2026-11-09): call the two enumeration stages of Fig.~\ref{fig:pipeline} Stage~1a and Stage~1b, not Stage~1 and Stage~1.5.} In tikz/pipeline\_full\_8pt1col.tex: the slab titles (``\textbackslash{}textbf\{Stage 1\}'' $\to$ ``Stage 1a'', ``Stage 1.5'' $\to$ ``Stage 1b''), the two ``Same as Stage 1'' ghost labels (``Same as Stage 1a''), and the caption sentence ``Stages~1 and~1.5 use Alloy'' $\to$ ``Stages~1a and~1b''. Stage~2 and Stage~3 keep their numbers. The same renaming applies wherever the text says Stage~1 or Stage~1.5 for these two: results.tex (the Coverage item, ``Stage-1 graphs'' / ``Stage-1.5 graphs''), pipeline.tex, intro.tex, discussion.tex, and the other notes in this file. The other pipeline variants (tikz/pipeline\_full.tex, pipeline\_full\_8pt.tex) are unused; rename there only if you bring them back. \textbf{Status (2026-11-11):} the figure and its caption are renamed, every note now says Stage~1a/1b, and STAGE NAMES notes in pipeline.tex mark the prose to change (the stage walk-through near the top of the section, the Test Generation subsubsection, the interleave caption). The hidden to-do lists in results.tex and discussion.tex carry a comment. Once those are done, delete this note.}
% scope of the types of defenses we can represent
For example, we define the defense specification for Speculative Taint Tracking, or STT, a defense that prevents leakage by delaying 'tainted' transmitters. 
\begin{enumerate}
    \item \textbf{Leakage contract:}
\verb|(load, load address)|, \verb|(branch, condition)|
\item \textbf{Execution contract:} 
An instruction is speculative when it is uncommitted and there are older, unresolved branches. 
\item \textbf{Protection set:}
All memory; a load with no older unresolved branch removes the location it reads from the protection set.
\end{enumerate}
% STT's leakage function is (load, load address), (branch, input register), and (reg-op transmitter, input address). STT does not consider stores to transmit any operands. STT's execution contract can be set as either that an instruction is non-speculative upon commit, or that an instruction is non-speculative when there exists no older unresolved branches. STT's protection set is all memory, with no protection set updates. 

% \subsection{Execution contracts in practice}
% \subsection{Scope of defenses}
% \begin{wrapfigure}{r}{.5\textwidth}
%   \begin{center}
%     \includegraphics[width=.5\textwidth]{images/examplefigure.png}
%   \end{center}
%   \caption{We generate a test from STT using our pipeline.}

%   \label{fig:pipeline}
  
% \end{wrapfigure}
% \begin{figure}[h]
%   \centering
%   \includegraphics[width=0.4\textwidth]{images/examplefigure.png}
%   \caption{We generate a test from STT using our pipeline.}
%   \label{fig:pipeline}
% \end{figure}

\section{Test Generation Pipeline}
We can use a formal leakage predicate for a defense as input to a generation pipeline to produce a set of litmus tests that induce exactly the $\mu$arch leakage scenarios the defense implementation ostensibly prevents. We will then run this set on the defense hardware implementation \anna{hardware under test??}, to ensure that it blocks all such leaks. We implement a pipeline that produces a set complete across high-level types of leakage with additional random fanout to test implementations robustly while avoiding test set explosion. We  run small experiments and analyze examples to ground our selection of variables to focus on during test generation. 
\anna{delete commented i think:}
% We designate the design dimensions we identify as either \textbf{1. comprehensively enumerated, 2. semi-enumerated, or 3. fixed}, where (2) and (3) are heuristically determined, and expose all dimensions of the pipeline to tuning. 


Fundamentally, our generation pipeline is comprehensive across types of leakage backbones, that is, it will generate all possible leakage backbones as blocked by the defense, but is heuristic-based in the concretization phase, e.g. it does not test every combination of eligible x86 .asm concretization for each instruction. 

The pipeline takes in a defense leakage predicate,  uses it to constrain an execution model and generate violating execution graphs (\ref{stage1}), then concretizes those execution graphs to program-annotation pairs
\urgentbr
\urgentnote{\textbf{STAGE NAMES (Fig.~\ref{fig:pipeline} now says Stage~1a, 1b, 2, 3; renamed 2026-11-11).} This sentence is the one place the text walks through the stages, so name all four here. It also has an undefined reference (\texttt{\textbackslash{}ref\{stage1\}}, no such label; it prints ??) and is unfinished. \textbf{PROPOSED TEXT} (replaces the sentence just before this note): The pipeline takes a defense's leakage predicate and uses it to constrain an execution model in Alloy: Stage~1a enumerates the minimal violating execution graphs, and Stage~1b grows each one into larger graphs under loosened minimality (Fig.~\ref{fig:pipeline}). Stage~2 concretizes each graph into program-annotation pairs, and Stage~3 runs them on the defense's gem5 build and checks for transmission.}
\notebr


% \input{tikz/pipeline_full}
\anna{define leakage backbones}


% Once we have a formal model of a defense's security guarantee, we can use this specification as input to a test generation pipeline to build a set of litmus tests \textit{that are tuned to exhibit microarchitectural leakage} according to those guarantees.
\subsection{Abstract test generation}

Since side-channel leakage is a $\mu$arch behavior, to be comprehensive across the leakage space we enumerate tests first as abstract execution graphs, with abstract instructions and $\mu$arch state annotations that represent a time snapshot of a processor. We aim to enumerate \textit{all possible fundamental forms of Spectre leakage} that are blocked by the target defense, and then concretize these abstract \textit{leakage backbones} 
into (1) assembly tests and (2) corresponding speculation and other $\mu$arch directives with which to non-invasively instrument our simulator to elicit leakage. A \emph{leakage backbone} is a minimal violating execution graph: removing or relaxing any of its instructions, operands, state elements or flags makes it non-violating. First, we must have a format for comprehensive enumeration of the forms of Spectre leakage.

\subsubsection{Test Generation (Alloy)}
We define constraints on execution graphs in the Alloy relational modeling language, so that we can use the Alloy Analyzer SAT-based finite model-finder to solve for all satisfying graphs up to a bound~\cite{jackson2002alloy}. As shown in figure \ref{fig:pipeline}, there are 3 types of constraint:
\urgentbr
\urgentnote{\textbf{STAGE NAMES}: this subsubsection describes Stage~1a, so say so. \textbf{PROPOSED TEXT} (edits the sentence just before this note): As shown in Stage~1a of Fig.~\ref{fig:pipeline}, there are three types of constraint: (and, after the list, the Stage~1a note below already proposes the UNSAT sentence).}
\notebr
\begin{enumerate}
    \item \textbf{Execution well-formedness constraints}: we must enforce that an arbitrary graph of instructions can arise from a program. E.g. the speculative program order (spo) relation must be a total order on instructions. These constraints are defense-agnostic.
\notebr
    \item \textbf{Minimality requirements:} we enforce that our graph is minimal along certain dimensions, that is all components specified in the graph are load-bearing to the execution violating the leakage predicate. This is measured by requiring every \textit{relaxation} of the graph, i.e. removal or weakening of an element, to be non-violating.
\notebr
\auditnote{write}{writing, high}{Decided 2026-11-08: Stage~1a adds instruction minimality (RI) to the models' own RR, RC+RR and RS checks; RO is defined but enforced in neither stage; Stage~1b drops RI again. The text should say which stage checks which relaxation. \newline \textbf{PROPOSED TEXT} (edits the sentence just before this note): In Stage~1a, a graph is minimal if removing any instruction (RI) or any state element (RS), resolving any unresolved instruction (RR), or committing and resolving the first uncommitted instruction (RC) makes it non-violating; Stage~1b keeps RR, RC and RS but not RI. \newline \textit{Source:} STAGE1\_alloy/models/proposed\_twophase/STT\_6\_core\_ri.als:289-302 (gen\_useful\_litmus, RI at :301); STT\_6.als:301-311; URGENT\_TODO.md 'Stage 1 minimality' decision (2026-11-08)}
\notebr
    \item \textbf{Defense leakage predicate:} we construct the defense leakage predicate from the three components in section \ref{formalmodel}, and require our execution to violate the predicate (i.e. be an example of leakage)
\notebr
\end{enumerate}
% The model describes a space of executions partitioned by the leakage predicate, which allows us to 
Alloy synthesizes all execution graphs up to a size bound that satisfy these constraints by iteratively model-finding and blocking discovered graphs until the problem is UNSAT.
\notebr
\auditnote{write}{writing, high}{Partly out of date. Stage~1a now enumerates up to six instructions (`but 6') by model-finding and blocking until UNSAT, and reaches UNSAT for STT and ReCon and their store-bypass variants (STT\_6: 129 cores; STT\_6\_stl: 306; every `rest' coverage partition UNSAT). SPT's Stage~1a does not converge: its four-instruction partition hits the 200,000 cap. Stage~1b enumerates a capped number of graphs per Stage-1a graph (about 300,000 per family). So only the Stage-1a cores are exhaustive, and only for STT and ReCon. The same claim is in the intro (`complete in including all leakage backbones') and in the knobs subsection (`coverage is complete up to the instruction bound'). \newline \textbf{PROPOSED TEXT} (edits the sentence just before this note): Stage~1a enumerates every minimal violating graph of up to six instructions by repeatedly model-finding and blocking each graph found until the problem is UNSAT (for SPT, up to a cap); Stage~1b then enumerates up to a fixed number of six-instruction graphs around each one. \newline \textit{Source:} wiki/state/EVAL\_MANIFEST.md `Stage 1' and `Stage 1.5' sections; STAGE1\_alloy/models/proposed\_twophase/STT\_6\_core\_ri.als:308-310}
\notebr
\subsubsection{Relaxations}
We define a minimality criterion on execution graphs in order to limit the problem solutions so that we can synthesize all \textit{leakage backbones}, or fundamental forms of leakage. Then we can later incorporate complexity by concretizing one backbone into a program in multiple ways, with the assurance that we have at least some representative program for all backbones. This also allows us to group executions that exhibit the same form of leakage, and restrict initial synthesis to the smallest example of each class, to keep our test set compact.
 % and make claims about fundamental bounds on leakage backbone size and quntity
% . We want to group executions that exhibit the same form of leakage, to allow us to reason about the simplest form of this leakage and build out in a principled manner.
% For this reason we follow MCM work~\cite{lustig} to define a minimality criterion restricts synthesis to the smallest example of each fundamental leakage backbone, keeping the initial test set compact while covering every leakage type. We synthesize all \textit{minimal, violating} executions up to a specified instruction bound. 
Our minimality criteria are established by a set of \textbf{relaxations}, which are perturbations on our execution instance, under which the instance should transition to a non-violating state. This ensures that all components of a generated execution are necessary for leakage, allowing us to make claims about the fundamental bounds on leakage backbone size and quantity.
\urgentbr
\urgentnote{\textbf{TODO (your note from the old draft, sections/oldtechnicaldetails.tex:113): add a natural-language intuition proof that every leakage backbone fits within a bounded length}, for STT first, using the running example. What the runs already show: Stage~1a reaches UNSAT at six instructions for STT and ReCon, so every minimal violating graph of up to six instructions is found. They do not show that no backbone needs seven or more instructions; that is what the proof would add. A cheaper empirical check is to rerun Stage~1a with a bound of seven and see whether any new minimal graph appears. Until then, ``claims about the fundamental bounds on leakage backbone size'' should say ``up to the instruction bound''; for SPT, whose Stage~1a stops at its cap, even the up-to-six claim does not hold yet. \textit{Source:} wiki/state/EVAL\_MANIFEST.md `Stage 1' section; the Stage~1a note in the Test Generation subsection above.}
\notebr

Figure~\ref{fig:relaxation} shows examples of tests that get minimized to the same test.
\input{tikz/relaxation}% the one relaxation figure (your four panes, redrawn), 2026-11-03
\notebr
\auditnote{fig}{figure, low}{\textbf{About Fig.~\ref{fig:relaxation}:} Pane (c) puts 'unresolved' on a Load, which breaks the rule that resolution is shown only on branches and stores. The reader is never told that an unflagged load (pane d) counts as resolved. The caption also leaves the dashed rf edge in (b) unexplained. \emph{Proposed:} Add: 'Non-branch rows without a flag word are resolved and uncommitted; dashed = a removed edge.' \emph{Source:} Fig. 4 p.5 render; VISUAL\_LANGUAGE.md CURRENT RULES.}
\notebr
\auditnote{fig}{figure, high}{\textbf{About Fig.~\ref{fig:relaxation}:} With RI-only Stage~1a (decided 2026-11-08) the caption's claim holds for Stage~1a: RI removes (b) and RR removes (c). (a) is removed too, but by RS (r1 appears only in the unneeded operand, and removing a state drops every operand on it), because RO is enforced nowhere. Stage~1b re-admits graphs like (b) by design; in the flat test sets behind the June numbers, (b) was never removed. This comes from reading the model, not from an Alloy run. \newline \textbf{PROPOSED TEXT}: relabel (a)'s arrow RS, or keep RO and add `RO is defined but not enforced; RS removes (a)'; end the caption with `\ldots{} so Stage~1a does not synthesize them.' \newline \textit{Source:} STAGE1\_alloy/models/STT\_6.als:214 (o\_p drops operands on a removed state), :301-311; proposed\_twophase/STT\_6\_core\_ri.als:301}
\notebr
We define the following relaxations:
\begin{enumerate}
    \item \textbf{RI}: remove instruction
    \item \textbf{RO}: remove operand
    \item \textbf{RS}: remove state
    \item \textbf{RC}: shift commit: flip the first uncommitted instruction to committed
\notebr
\auditnote{fact}{fact}{The code applies RC to the first uncommitted instruction, the one right after the commit boundary. 'Last-uncommitted' reads as the youngest instruction. RC is also never used alone: it is tested together with RR on that same instruction. When nothing is committed, first\_uncommitted falls back to the first instruction. \newline \textbf{PROPOSED TEXT} (edits the sentence just before this note): \textbackslash{}textbf\{RC\}: shift commit: mark the first uncommitted instruction (the commit boundary) committed and resolved \newline \textit{Source:} STT\_6.als:306-308 (secure[RC-\textgreater{}first\_uncommitted + RR-\textgreater{}first\_uncommitted]), :139 (fallback to first\_instr)}
\notebr
    \item \textbf{RR}: mark an unresolved instruction as resolved
\end{enumerate}

\subsubsection{Overlapping interleaving}
Once we have defined minimality criteria for an execution, we are then able to utilize all our relaxations to make claims about fundamental bounds on leakage backbones, as discussed in the results. However, we recognize a need for more complex tests, so that we can stimulate exercise of protection set tracking implementation, as well as other implementation mechanisms. For instance, a correct implementation of STT blocks the leak in Figure~\ref{fig:ddi-rf}, while the test in Figure~\ref{fig:interleave} adds a chain of instructions after the committed prefix that merges a second speculative source into an ALU op; this shape exposes the ReCon \texttt{getOldestTaint} bug (Sec.~\ref{sec:bug-otb}), where the implementation lets it leak.
\input{tikz/interleave_v4_8pt}
\notebr
\urgentbr
\urgentnote{\textbf{STAGE NAMES}: this figure's caption (in tikz/interleave\_v4\_8pt.tex) says ``no Stage~1 instance contains''; change it to ``no Stage~1a instance contains''. The claim holds for Stage~1a, whose graphs are minimal; Stage~1b drops instruction minimality, so it is not promised there.}
\notebr
\auditnote{fig}{figure}{\textbf{About the caption of Fig.~\ref{fig:interleave}:} With RI-only Stage~1a, `no Stage~1 instance contains' is now true by construction: the second source is redundant for the leakage predicate, so no minimal graph has it. The mechanism is what is out of date. Chain interleaving is off in every evaluation config, put an ALU on a load transmitter's address path in 0 of 1,200 variants (that address is always an Alloy-locked register, and chains write only free registers), and cannot hook onto a Stage-1a core. SimSpect reproduces the getOldestTaint bug with the OTB merge of fanout v2 instead: a second load under an earlier, resolved branch, xor-merged into the transmitter's address (flat test set, ReCon-STT build: 100 of 100 tests hit, 0 of 100 without the merge). The merge is on in the 2026-11-09 two-stage evaluation configs, but it applies to few two-stage programs (it needs a resolved branch before the unresolved one); where it applies, the ReCon-STT row hits the getOldestTaint drop in 11 of 24 OTB-transmitter tests per mode at the 2\% sample. \newline \textbf{PROPOSED TEXT}: redraw the figure and caption around the merge the evaluation uses, or caption this figure as the chain-interleaving capability rather than the route behind the result. \newline \textit{Source:} wiki/design/VOCAB\_TABLES.md (fanout v2 OTB); runconfigs/twostage\_fv2\_20261109/E2\_STT\_6\_fv2.jsonc (interleave off, fanout\_v2.otb on); URGENT\_TODO.md `OTB on two-stage tests'}
\notebr
\auditnote{fig}{figure}{\textbf{About the caption of Fig.~\ref{fig:interleave}:} Branch flags are inconsistent across the graph figures. Fig. 4, 6, 8, 9 and Table II write 'unresolved' on the branch row; Fig. 5 (i2) and Fig. 10 (i2) write nothing; Fig. 2 puts it only in the note column. Captions of Fig. 5, 6 and 10 also use subscripts (\$f\_0\$, \$r\_0\$, \$r\_1\$) where the figures print f0, r0.  \textbf{PROPOSED TEXT} (edits the sentence just before this note): Write the italic 'unresolved' flag on every unresolved branch (Fig. 2 i0, Fig. 5 i2, Fig. 10 i2), and use r0/f0 without subscripts in the captions.  \textit{Source:} Renders p.4-p.11; VISUAL\_LANGUAGE.md CURRENT RULES ('Resolution status ... shown only on branches and stores'; 'write r0, not \$r\_0\$').}
\notebr
\auditnote{fig}{figure, high}{\textbf{About the caption of Fig.~\ref{fig:interleave}:} The figure marks i0 and i1 as committed, so c0 (placed after them and before i2) is not speculative at all under STT's contract. As drawn, the merge has only one tainted source, and the caption's getOldestTaint mechanism cannot occur. The OTB gadget needs the injected load under a branch that is resolved early but still uncommitted.  \textbf{PROPOSED TEXT} (edits the sentence just before this note): Mark i0/i1 'resolved' instead of 'committed', or explain as Fig. 6 does that i0 is committed only at the observation time and resolves at a swept time.  \textit{Source:} Fig. 5 p.6 render (i0, i1 'committed'); OTB\_GENERATION\_GAP.md 'What OTB needs' (ld\_inj under br\_r, a resolved branch that resolves early); STT\_6.als speculation\_contract\_p.}
\notebr
Therefore, we interleave instructions in two rounds. The first round happens in Alloy, by removing some of the relaxations: the models we evaluate enforce only RR, RC and RS, and leave RI and RO defined but unchecked. The second round happens during concretization (Sec.~\ref{sec:conc-interleave}) and builds merges like the one in Figure~\ref{fig:interleave}.
\urgentbr
\urgentnote{Out of date since the two-stage decision (2026-11-08): the evaluated Stage~1a also enforces RI, and Stage~1b is the round that drops it; RO is enforced in neither. \textbf{PROPOSED TEXT} (edits the sentence just before this note): The first round happens in Alloy: Stage~1a enumerates minimal graphs (RI, RR, RC and RS), and Stage~1b grows each one with RI turned off, so it may add instructions that share the backbone's state (RO is defined but not enforced in either stage). \textit{Source:} STAGE1\_alloy/models/proposed\_twophase/STT\_6\_core\_ri.als:289-302; wiki/state/EVAL\_MANIFEST.md `Stage 1.5'.}
\notebr
The first round allows us to emit tests that include instructions not on the leakage backbone, but to minimize the load on the Alloy Analyzer, we still require all instructions to share state that is necessary to the leakage backbone. We will do further interleaving in the concretization phase for the non-state-sharing instructions. 
% Figure~\ref{fig:statetable} shows the types of instructions and state enumerated in each phase. To enable our first round of interleaving, we turn off the remove instruction relaxation, and the remove operand instruction, and require all operands to use essential state.

\subsection{Concretization Phase}
Given our abstract execution graphs which include $\mu$arch state, we must generate programs and simulator directive combinations to elicit this execution state. Fig.~\ref{fig:pipeline} (Stage~2) lists the concretization passes; we describe them in that order.

\subsubsection{Instruction Choice}
\reviewtag{[TODO: Applied your requests in this subsection; check it]}
First, we concretize our abstract instructions to assembly. We can tune our assembly translation by defining instruction tables which contain all possible candidates for each abstract instruction form. The SimSpect pipeline selects from the candidates that use enough operands to specify the abstract instruction as given in the Alloy execution. Each chosen instruction is emitted as an inline-assembly block with a global PC label
\notebr
\auditnote{fact}{fact, low}{Two qualifications. Branches are not picked at random: pass 2.5 rewrites every non-transmitter branch to br\_cond and the transmitter to br\_bez. Also, for a branch only the condition set-up is inline asm with the PC label. LLVM emits the compare and the jcc, so LLVM's block layout decides which jcc direction is the speculative path. That is why the harness derives the forced direction from the binary. \newline \textbf{PROPOSED TEXT} (edits the sentence just before this note): Each chosen instruction is emitted as an inline-assembly block with a global PC label (for a branch, the label sits on the inline-asm condition set-up and LLVM emits the compare and jump) \newline \textit{Source:} STAGE2\_compilation/parsexml.py:789-795 (branch rewrite), :2426-2427 and :2484 (icmp + br emitted by LLVM); wiki/design/RESOLUTION\_HOOKS.md sec.2}
\notebr
(Table~\ref{tab:instrtable}).
We select uniformly at random among the valid candidates, though we could exhaustively enumerate the cross product of valid concretizations.

\subsubsection{Speculation Primitive Instrumentation}
Table~\ref{tab:instrtable} lists, for each speculation primitive, the program Stage~2 emits (including whether each branch is mispredicted and which way it is taken) and the setting the Stage~3 harness applies.
The concretization pipeline must realize speculation primitives in a way that ensures that a test experiences the correct speculation as seen in the abstract execution. In the model, we introduce resolved and committed flags. At some observation time during execution, the run must match the speculative program order and resolved and committed flags in the modeled instance. An unresolved branch is always mispredicted, in both directions (architecturally taken and mispredicted not taken, or the reverse); its constant condition sends the architectural path to the exit, so the rest of the test runs only speculatively. A resolved branch is generated only as correctly predicted and not taken
\notebr
\auditnote{fact}{fact, low}{There is one exception. The transmitter branch is always emitted as br\_bez: architecturally taken, forced to mispredict, and given a resolve stall of 0, whatever its Alloy resolved flag. Table row 5 says this but the text does not. \newline \textbf{PROPOSED TEXT} (edits the sentence just before this note): A resolved branch is generated only as correctly predicted and not taken; the transmitter branch is the exception and is always mispredicted (Table\~{}\textbackslash{}ref\{tab:instrtable\}, row\~{}5). \newline \textit{Source:} STAGE2\_compilation/parsexml.py:738-776 (xm override); pipeline.py:404-421 (\_inject\_stalls: xmit stall 0)}
\notebr
(Table~\ref{tab:instrtable}, rows 1--2). Every case can induce leakage while the branch is unresolved, even if it is not ultimately mispredicted, but we do not generate correctly predicted unresolved branches at this time for simplicity.
We also sweep the resolution time of resolved branches, by setting a delay from the time they receive their direction and target and the time they transmit ``resolved,'' which we will direct in the gem5 implementation. This allows testing of multiple branches with different resolution orders, which can induce different behaviors (Fig.~\ref{fig:resolutiontiming}); the ReCon bug of Fig.~\ref{fig:otb} is one such case (sweep values: Table~\ref{tab:instrtable}, row~2). For unresolved branches, we direct them to resolve with a fixed 5000-cycle stall, and later enforce that this stall lasts until after our observation time.
\input{tikz/instrtable_v2_8pt}
\notebr
\auditnote{fact}{fact}{\textbf{About the caption of Table~\ref{tab:instrtable}:} No evaluated model can produce an Other\_x transmitter. Otherxs is commented out of the leakage function in STT\_6, Recon\_6 and SPT\_6\_oneLP, and tag\_xm requires the xmit to own a leakage-function operand. So this row, check\_other.py and the divide-fault promotion never ran in the reported campaigns. The same applies to pipeline.tex's 'Faulting / long-latency other transmitter' bullet.  \textbf{PROPOSED TEXT} (edits the sentence just before this note): (caption) ... Other\textbackslash{}\_x rows are supported by the pipeline, but no evaluated defense model lists Other\textbackslash{}\_x as a transmitter.  \textit{Source:} STT\_6.als:327-328, Recon\_6.als:320, SPT\_6\_oneLP.als:344 (leakage\_function); STT\_6.als:270-273 (tag\_xm); STAGE3\_gem5/gem5\_common.py:772-852; OTB\_GENERATION\_GAP.md ('no other\_x transmitters in any STT\_6 testset')}
\notebr
\auditnote{fig}{figure, low}{\textbf{About the caption of Table~\ref{tab:instrtable}:} Table II says 'first uncommitted' where the text and Fig. 7 say 'first non-committed', and it adds an 'uncommitted' flag word that Table I does not define. The 'Target transmitter' row covers only Branch\_x and Other\_x (calling the latter 'Divide'), so readers cannot see that Load/Store transmitters need no special handling, or that Other\_x is a divide or zero-skip multiply.  \textbf{PROPOSED TEXT} (edits the sentence just before this note): Use 'first non-committed'; add to row 5: 'Load, Store: no special code'; say 'Other\_x (divide)'.  \textit{Source:} Table II p.6 render; plans/fig\_d\_instrtable.md Row 5 (kinds: load, store, branch, zsm/divide).}
\notebr
\input{tikz/resolutiontiming}% your figure, redrawn in the new notation, 2026-11-03
\notebr
\auditnote{fig}{figure}{\textbf{About the caption of Fig.~\ref{fig:resolutiontiming}:} The graph violates STT's leakage predicate under both options; what changes is whether a getOldestTaint defense releases i5. Labeling only option 2 'violation' (in red) conflicts with the meaning of 'violation' in Sec. IV and with Fig. 2.  \textbf{PROPOSED TEXT} (edits the sentence just before this note): "(1) i0 resolves before i4 executes: ReCon holds i5" / "(2) ... after: ReCon releases i5 (bug)".  \textit{Source:} Fig. 6 p.7 render; pipeline.tex violation definition; Fig. 6 caption ('the run commits without violating the contract').}
\notebr
% \urgentnote{\textbf{VOCAB}: In Fig.~\ref{fig:resolutiontiming} the graph is a violation under both options; what differs is whether the run leaks. Change the figure's notes ``no violation'' to ``no leak'' and ``violation'' to ``leak'', and in its caption ``the run commits without violating the contract'' to ``the run does not leak''.} 
\urgentbr
\urgentnote{Answer to the comment above (``this makes no sense''): in the paper's terms, \emph{violation} is a property of the abstract graph, and this graph breaks STT's leakage predicate under both timings: $m1$ reaches i5's address through i3 and i4 while i2 is unresolved. The two options differ only in what the gem5 run does. In option~1 the defense holds i5 (no hit); in option~2 the getOldestTaint tracking lets i5 send its access (a hit: the bug). So the figure's notes should name what the run does, not ``no violation''/``violation''. The commented line proposed ``no leak''/``leak''; the hidden figure note above proposes ``ReCon holds i5''/``ReCon releases i5 (bug)''. Either works; pick one and delete the commented line. \textit{Source:} tikz/resolutiontiming.tex:104-105 (the two notes), :109 (caption).}
\notebr
\auditnote{fig}{figure}{\textbf{About the caption of Fig.~\ref{fig:resolutiontiming}:} Fig. 6 and Fig. 10 draw the same gadget (per the source comment) in different ways: Branch\_n vs Branch\_x, r2 \textless{}- r1,r0 vs r0 \textless{}- r1,r0, protected m1 vs m0 with no protection mark. Fig. 10's caption calls i1 'speculative under branch i0' while marking it committed, without Fig. 6's explanation.  \textbf{PROPOSED TEXT} (edits the sentence just before this note): Draw Fig. 10's graph with Fig. 6's kinds, names and protection mark, or say in Fig. 10's caption that i0/i1 are committed only at the observation time (as in Fig. 6).  \textit{Source:} Fig. 6 (p.7) vs Fig. 10 (p.11) renders; resolutiontiming.tex header comment 'the same two-branch merge shape'; RESULTS\_FIGURES.md §7 ('m0 is not gray').}
\notebr
For data-flow speculation, we instrument store bypass~\cite{horn2018ssb} (a store-to-load-forwarding misprediction) only in the false negative direction, where a load mispredicts no alias with an older store to the same address, whose address we hold unresolved, so the load bypasses it and reads the stale value until the memory-order squash. For this reason, we also keep unresolved stores unresolved for 5000 cycles, and give the load and the store the same address.

\subsubsection{Interleaving and State Assignment}\label{sec:conc-interleave}
\reviewtag{[TODO: Applied your requests in this subsection; check it]}
We then add a second round of interleaving: we inject one or two synthetic chains of loads, stores and ALU ops at random positions after the committed prefix, with every chain operand drawn from a shared free pool of unspecified state (Fig.~\ref{fig:interleave}). Because a chain's output can feed an unspecified operand of the backbone, interleaving builds transmitters fed by two independent speculative sources, which no single Alloy instance can express;
\notebr
\auditnote{fact}{fact}{This is a capability the evaluated pipeline does not use. Chain interleaving is off in every evaluation config, and on 300 STT\_6 load-transmitter instances it put an ALU on the transmitter's address path 0 of 1,200 times, because that address is always an Alloy-locked register and chains write only free registers. `Which no single Alloy instance can express' is also wrong: the shape is satisfiable when forced; minimality drops the redundant second source, so no Stage-1a core contains it. SimSpect reproduces the getOldestTaint merge with the OTB merge of fanout v2 (flat test set, ReCon-STT build: 100/100 tests hit, 0/100 without it), which is on in the 2026-11-09 evaluation configs. \newline \textbf{PROPOSED TEXT} (edits the sentence just before this note): interleaving can feed a transmitter from two independent speculative sources, a shape that minimal graphs never contain; \newline \textit{Source:} wiki/design/VOCAB\_TABLES.md (interleave, fanout v2 OTB); wiki/findings/mdsToRead/OTB\_GENERATION\_GAP.md (forcing is SAT); runconfigs/twostage\_fv2\_20261109/*.jsonc}
\notebr
chains off the backbone also add program length and timing variety.
We have two sets of state elements — those determined in Alloy, and those left unspecified. We handle the predetermined state by tying physical registers and memory locations to Alloy register and memory atoms. We then reserve those state locations, and draw every unspecified operand from the shared free pool, the same pool the interleaved chains use
\notebr
\auditnote{fact}{fact}{Only unspecified register operands share the pool. Each unspecified memory operand gets a fresh, never-aliased slot. An unspecified address register of a load or store is not emitted: a constant offset is used instead. Without interleaving, the pool holds no chain registers. \newline \textbf{PROPOSED TEXT} (edits the sentence just before this note): draw every unspecified register operand from the shared free pool, the same pool the interleaved chains use (unspecified memory operands get fresh, unaliased slots) \newline \textit{Source:} STAGE2\_compilation/parsexml.py:1601 (pass3: fresh mem slot), :2146 (pass5: only rf-specified inaddr used), pass\_interleave returns empty pools when disabled}
\notebr
so an unspecified input can read an injected chain's output. All registers employed in memory accesses take the form of offsets from a base pointer to prevent segmentation faults due to random accesses to unallocated memory.
Finally, we convert each program to SSA form, emit LLVM IR with a label at every PC, compile it to x86 assembly, and write the annotation that Stage~3 reads (Fig.~\ref{fig:pipeline}, Stage~2, passes~6--7).
\urgentbr
\urgentnote{\textbf{PROPOSED TEXT (how a test survives compilation; insert after the sentence above).} The compiler must not drop, merge, reorder or re-register any instruction the model specifies. Each instruction is an \texttt{asm sideeffect} block, which LLVM never removes or reorders, and every register the model fixes is pinned to its physical register by the block's register constraint (for a branch, only the condition set-up is inline assembly; LLVM emits the compare and the jump). Memory is a single array whose slots are zero-initialised with volatile stores. We compile with clang at \texttt{-O2}: at \texttt{-O0} the register allocator spills every live value at each basic-block boundary despite the pinning, and the extra loads and stores cause spurious memory-order violations on the speculative path. After linking, the harness reads each PC label's address back from the binary, and the direction of each compiled branch jump, and records them in the annotation, so every directive targets the instruction the model names. \textit{Notes:} two details of the old draft (oldtechnicaldetails.tex:88) are out of date: the code no longer uses \texttt{fake.use}, and the register pool is no longer caller-saved only (it adds r12--r15, which LLVM saves and restores). \textit{Source:} STAGE2\_compilation/compile\_annotate.py:86-105 (the -O2 rationale), :137-222 (PC labels from nm, branches from objdump); STAGE2\_compilation/parsexml.py:1428-1439 (register pool), :2097-2100 (volatile zero-initialisation), :2191 (pinned register write); the instruction-choice note above (branches).}
\notebr

\subsection{Gem5 Execution and Instrumentation}
\reviewtag{[TODO: Applied your requests in this subsection; check it]}
Our tests and generated speculation directives must finally be run on the gem5 simulator. The final phase of testing has a few imperatives to ensure that each hit's gem5 trace matches the modeled violation (violation = the model's graph; hit = what gem5 flags). 

It enforces resolution and speculative program order by using the white-box nature of the simulator to actively set this state. However, it enforces transmitter transmission at a time that matches committed annotations by passively measuring the pipeline, and this is how we detect a violation.
\urgentbr
\urgentnote{\textbf{VOCAB}: ``this is how we detect a violation'' $\to$ ``this is how we detect a hit''.}
\notebr
\urgentbr
\urgentnote{\textbf{PROPOSED TEXT (how we use the white-box simulator; from the old draft, oldtechnicaldetails.tex:42 and :98-101).} White-box access lets the harness set each condition of the modeled execution directly instead of coaxing it. (1) \emph{Speculation}: instead of training the predictors, it overrides them at each annotated PC, forcing a branch to mispredict or to predict correctly, and on the store-bypass build forcing a load to predict no alias with an older store. (2) \emph{Timing}: it holds a branch's resolution, and on the store-bypass build a store's address, for a set number of cycles, so the window is wide and its length is known. (3) \emph{Observation}: it reads the outcome from the simulator's debug traces, which give the tick at which each instruction executes, obtains data, resolves and is squashed, so a hit is read off the trace rather than inferred (see Transmission Requirement). The defense's own logic is never modified. \textit{Placement:} this half of old line 42 belongs here, with the hooks. The other half (realizing the same directives on a target without white-box access by pretraining the predictor) is about generality and goes in the Discussion; there is a note there. \textit{Source:} the gem5-hooks paragraph of this section; gem5-hooks/README.md; wiki/design/RESOLUTION\_HOOKS.md sec.~1.}
\notebr

\subsubsection{Speculation Directives}
The simulator must be made to actually realize the modeled speculative scenario --- the intended branches and dataflow must mispredict, the specified branches must stay unresolved through the transmitter's action, and the commit boundary must stay where is specified \emph{without} altering anything the defense does. If the Alloy model describes an abstract scenario (``a branch mispredicts, opening a window; a load executes speculatively inside it''); a stock CPU will not necessarily realize that exact scenario for a given litmus binary (the predictor may guess right, or the window may be too short for the transmitter to fire, even if the defense might permit this). We bridge this gap with a small, fixed set of hooks, partitioned into \emph{control} hooks (C++ source in the gem5, which realize the scenario), \emph{config glue} (Python, which drives the hooks from the per-test annotation), and \emph{observation} (existing debug flags, which expose the transmitter, as described below). The governing constraint is that we touch only these hooks and never the defense's taint/protection logic, and each control hook is paired with a control run --- the fnc-stall on vs.\ off, the branch-annotation on vs.\ off, and the \texttt{Unsafe}-scheme baseline --- proving the hook did not itself manufacture the result.
\notebr
\reviewnote{\textbf{CHECK}: of the three controls, only the Unsafe baseline is in the evaluation (row E1, run on the 2026-11-09 two-stage sample). The fnc commit stall is 0 in every evaluation config, so its on/off pairing does not apply, and no run turns the branch annotation off. Reword to ``can be paired'', or name only the Unsafe baseline. \textit{Source:} runconfigs/twostage\_fv2\_20261109/*.jsonc (fnc\_commit\_stall\_cycles 0); wiki/state/EVAL\_MANIFEST.md rows E1-E14.}
\notebr

% \paragraph{Directives passed from concretization.}
The concretization phase emits, per test, an annotation (\texttt{.ann.json}) carrying: the BTB target each forced branch will receive (where it will predict to); the resolution time of every speculation primitive (the per-branch resolve delay, and the long stall held on unresolved branches and stores); and the (load, store) PC pairs that should incorrectly alias for store bypass.
The same annotation file both tells gem5 how to reproduce the window and tells the checker what to look for.
\notebr
\reviewnote{\textbf{FIX}: Stage~2 writes no timing, and the pairs do not ``incorrectly alias'': the load is forced to mispredict no alias. \newline \textbf{PROPOSED TEXT}: ``\ldots{} carrying each element's PC, each branch's mispredict mode and BTB target, the transmitter's kind, and for store bypass the (store, load) PC pair whose load is forced to mispredict no alias. Stall lengths come from the run configuration.''}
\notebr

\subsubsection{gem5 hooks}
Three source control hooks realize the modeled scenario, each acting only on control or timing (prediction, redirect latency, commit pacing) and never on the defense's data path. A \emph{BTB-forced target} overrides the predictor at an annotated branch PC to force the modeled mispredicted (taken) path.
\notebr
\reviewnote{\textbf{UPDATE}: on the hooked builds every annotated branch gets a forced prediction in either direction: unresolved branches are forced to mispredict, resolved branches to predict correctly (Table~\ref{tab:instrtable}, rows 1--2).}
\notebr
A \emph{per-branch resolve-stall} delays a branch's resolution to slide the speculative window
over the transmitter. A \emph{commit-stall} holds the first uncommitted instruction --- the speculation boundary --- at the ROB head to widen the window. Rather than fix these stalls, we sweep them so that a reported leak is shown not to hinge on any single window setting:
\notebr
\auditnote{fact}{fact, high}{The code does the reverse. A test is reported as leaking if any grid point triggers (OR over the sweep), so a reported leak can depend on a single setting. The per-point outcomes are saved in the *\_sweep.json sidecars (triggered\_points), so robustness could be reported, but the verdict does not mean that. \newline \textbf{PROPOSED TEXT} (edits the sentence just before this note): We sweep these stalls and report a test as leaking if any setting triggers it; per-setting outcomes are kept, so for each leak we also report how many settings trigger it. \newline \textit{Source:} pipeline.py:330-339 (comment: 'aggregates per-test as OR over grid points'), :500-558 (\_aggregate\_sweep\_results: any\_hit, triggered\_points)}
\notebr
the resolved-branch stall over $s_R \in \{0, 500, 2500\}$ cycles, the unresolved-branch stall $s_U$ (a long $5000$-cycle value masks leaks by pushing the transmitter past \texttt{fnc\_retire}, so runs are launched bare at $s_U = 0$), and the commit-stall over $\{0, 150\}$ cycles.
\urgentbr
\urgentnote{These are the pre-fix settings and contradict Table~\ref{tab:instrtable} and the speculation-directives paragraph. The evaluated configs sweep $s_R \in \{0, 500, 2500\}$ (footprint rows and the AMuLeT STT row: 0 only), hold every unresolved branch for $s_U$ cycles (5000 in the 2026-10-30 fixed configs; the 2026-11-09 two-stage configs draw $s_U$ per branch from $\{2500, 5000, 10000\}$), and set no commit stall. The claim that a 5000-cycle $s_U$ masks leaks predates the resolve-stall fix. \textbf{PROPOSED TEXT} (edits the sentence just before this note): ``\ldots{} the resolved-branch stall over $s_R \in \{0, 500, 2500\}$ cycles; unresolved branches are held for $s_U$ cycles, long enough that they stay unresolved past the transmitter, and the commit stall is off (Table~\ref{tab:instrtable}).'' \textit{Source:} runconfigs/fixed\_20261030/STT\_6\_sttmaster\_fixed\_latest.jsonc:47-57; runconfigs/twostage\_fv2\_20261109/E1\_STT\_6\_fv2.jsonc (sweep).}
\notebr
Finally, for data-flow speculation we instrument store bypass in the false-negative direction --- a load fails to find a store in its store set --- by holding the store unresolved and forcing the aliasing load and store to the same address.
\notebr

Realizing these directives requires a small, fixed set of source sites in every defense build, which is also the checklist for porting the pipeline to a new defense: \texttt{pred/BranchPredictor.py} declares the annotation parameters (forced branch PCs, BTB targets, resolve-stall cycles); \texttt{pred/bpred\_unit.\{cc,hh\}} initializes those maps and overrides the branch target and resolve-stall at predict time; \texttt{o3/\{O3CPU,BaseO3CPU\}.py} declares the fnc-commit-stall parameters; \texttt{o3/cpu.\{hh,cc\}} carry the stall into the CPU; \texttt{o3/commit\{.hh,\_impl.hh\}} add the stall loop, before the squash check, that holds the first uncommitted instruction at the ROB head; \texttt{common/Options.py} exposes the four CLI flags; and \texttt{example/se.py} loads the per-test annotation and wires the targets and stalls into each CPU's branch predictor. Every build also carries two later hook sets. The \emph{resolve hooks} hold a branch's resolution (its execute completion, hand-off to commit and squash broadcast) for the annotated number of cycles, force its prediction, wake the CPU when the hold is released, and print the branch's predict and resolve events. On the store-bypass build, the \emph{store-bypass hooks} hold the annotated store's address generation for the configured number of cycles and force the paired load's memory-dependence prediction to no alias.
\reviewtag{[TODO: Added the resolve and STL hooks (gem5-hooks/README.md; RESOLUTION\_HOOKS.md; STL\_SPECULATION.md)]}
\notebr
\reviewnote{\textbf{FIX}: not every build carries both later hook sets. The resolve hooks are on all ten hooked builds; the store-bypass hooks are only on the store-bypass build, which predates the resolve hooks (rebuild planned), and its scorer is a draft, not a window checker. \newline \textbf{PROPOSED TEXT}: ``Every hooked build also carries the \emph{resolve hooks} \ldots{}; the store-bypass build adds the \emph{store-bypass hooks} \ldots'' \textit{Source:} wiki/state/EVAL\_MANIFEST.md STL section; project log 2026-11-08 21:30 (recon-ruby) and 2026-11-09 13:25 (stl-campaign).}

Detecting transmission requires the pipeline to actually \emph{print} the values the checker reads, so we run gem5 with its existing \texttt{O3PipeView}, \texttt{LSQUnit}, and \texttt{Commit} debug flags --- which emit, respectively, per-instruction issue/complete/squash ticks, the load's execute/packet/spec-read ticks, and the pc$\to$squash-tick map --- and parse that trace. No cache or defense code is instrumented to \emph{see} a leak, only to \emph{realize} the scenario;
\notebr
\auditnote{fact}{fact}{SimSpect does add instrumentation whose only job is observation. It adds print-only BR predict/resolve lines in bpred, IEW and commit on all ten hooked builds (the eight of RESOLUTION\_HOOKS.md plus stt-master and recon-ruby3), and cache, TLB and buffer event prints in the Ruby and classic cache code of the InvisiSpec, CleanupSpec and SpecLFB builds. The resolve hooks also change behaviour: the release now wakes the CPU. Also in this paragraph: O3PipeView prints fetch-to-retire ticks but no squash ticks. Squash ticks come from LSQUnit and Commit. \newline \textbf{PROPOSED TEXT} (edits the sentence just before this note): Observation hooks are print-only: they add trace lines (branch prediction and resolution, and on the footprint builds cache and buffer events), verified to leave the rest of the trace unchanged; no defense logic is modified. \newline \textit{Source:} wiki/design/RESOLUTION\_HOOKS.md sec.1 (BR lines; edit 4 wakeCPU); wiki/design/SIMSPECT\_HOOKS.md header; gem5-hooks/README.md table; STAGE3\_gem5/gem5\_common.py (parse\_pipeview) (parse\_pipeview fields)}
\notebr
where a build does not already print a needed signal, we add a minimal print of the missing tick rather than touch any defense logic.

\subsubsection{Transmission Requirement}
Given a target transmitter defined by our leakage backbone, we must specify how to detect transmission. We do this passively: unlike resolution and speculative program order, which we enforce by actively setting state in the white-box simulator, transmission is detected by \emph{measuring} the pipeline, so that we never perturb the defense's own data path. We define detection on a \textbf{per-transmitter pipeline predicate} which takes the gem5 execution trace (the existing \texttt{O3PipeView}, \texttt{LSQUnit}, and \texttt{Commit} debug streams, parsed by our harness) and analyzes it for the transmitter's characteristic signal. If that signal occurs during the \textbf{violation window}, we report a hit.
Fig.~\ref{fig:checkerwindow} shows this decision on a real gem5 trace.
\urgentbr
\urgentnote{\textbf{PROPOSED TEXT (no attacker or probe needed; from the old draft, oldtechnicaldetails.tex:107).} Because we read the simulator's trace directly, detection needs no attacker: we neither time a probe nor compare two runs. The attacker model is already part of the specification. The leakage contract names the transmitters and operands the defense treats as observable, so a hit is simply one of those transmissions inside the violation window: for a load, the transmitter obtained its data; for the footprint defenses, a squashed instruction changed a structure the defense protects. AMuLeT, by contrast, runs a program on pairs of inputs whose contract traces match and flags a difference in the microarchitectural state it observes~\cite{amulet}. \textit{Source:} oldtechnicaldetails.tex:107; background.tex:23 and the AMuLeT note in intro.tex (fu-2025-amulet.pdf \S2.3); STAGE3\_gem5/check\_ld.py (data-obtained gate); the footprint note on Fig.~\ref{fig:checkerwindow} below.}
\notebr

\notebr
\reviewnote{Pick one window name: the checker figure's caption calls it the unresolved window and ends it at the branch's resolution.}
\notebr

\input{tikz/checkerwindow_v4_8pt}% moved up to where it is first cited, 2026-11-03
\urgentbr
\urgentnote{\textbf{VOCAB}: In Fig.~\ref{fig:checkerwindow} (file tikz/checkerwindow\_v4\_8pt.tex), relabel ``Unresolved window'' to ``Violation window''. In its caption, change ``How the checker decides a violation'' to ``How the checker decides a hit'', ``the unresolved window runs from'' to ``the violation window runs from'', and ``reported as a violation'' to ``reported as a hit''.}
\notebr
\auditnote{fig}{figure, high}{\textbf{About the caption of Fig.~\ref{fig:checkerwindow}:} This example is now in scope. Since 2026-11-08 each footprint criterion checks only the structures its defense's paper claims (for InvisiSpec, MSHR contention and the I-cache are out, the L1D is in), and UV1 is found in 2,995 of 3,000 tests on the hooked InvisiSpec build. What is out of date is the results section's scope paragraph, if it still calls UV1 collateral. \newline \textbf{PROPOSED TEXT}: keep this example, caption it as InvisiSpec's footprint criterion, and update the scope paragraph in the results. \newline \textit{Source:} wiki/findings/mdsToRead/AMULET\_BUGS\_DETECTION\_ACCOUNTING.md sec.~2 (UV1 row); specs/leakage/invisispec.json}
\notebr
\auditnote{claim}{claim}{\textbf{About the caption of Fig.~\ref{fig:checkerwindow}:} check\_footprint counts a footprint by any squashed instruction younger than lc, not just the transmitter. In this trace seqs 566, 573 (xmit) and 576 all qualify, and the earliest in-window footprint is seq 566's L2 coherence change at cycle 16, so the verdict does not depend on xmit's eviction. The fixture also predates the 2026-10-30 harness fix (5000-cycle resolve stall).  \textbf{PROPOSED TEXT} (edits the sentence just before this note): xmit's L1D eviction (\$\textbackslash{}star\$) is one of the in-window footprints that meet InvisiSpec's criterion (the criterion counts any squashed load younger than lc), so the test is reported as a violation.  \textit{Source:} project log 2026-11-02 'InvisiSpec footprint checker attributes leaks to any squashed load'; paper\_figures/plans/fig\_e\_checkerwindow.md:140,199; URGENT\_TODO.md InvisiSpec footprint DECIDE}
\notebr

The violation window
\notebr
is the interval during which the \emph{last committed} instruction (\texttt{lc}) has committed but the \emph{first uncommitted} instruction (\texttt{fnc}) has not. Because commit is in order, this is the interval \texttt{(lc\_retire, fnc\_retire)}, with the additional requirement that all branches guarding the transmitter are still unresolved at the transmission tick.
\notebr
\reviewnote{\textbf{REWORD} ``all branches guarding the transmitter'' to ``the branches the model marks unresolved'': that is the gate the checker applies.}
\notebr
A transmission inside this window, on the modeled mis-speculated path, is a violation;
\urgentbr
\urgentnote{\textbf{VOCAB}: ``is a violation;'' $\to$ ``is a hit;'' (a gem5 result; it becomes a leak only once validated).}
\notebr
the same action after the window closes is not (Figure~\ref{fig:checkerwindow}).
We call this the \emph{violation window} because it is where every $\mu$arch condition the instance specifies still holds (which instructions have committed and which branches are unresolved), not only the commit status of \texttt{lc} and \texttt{fnc}; a transmission inside it is exactly what the leakage predicate forbids.

% \input{tikz/xmitcriteria_8pt}% hidden 2026-11-03 (Anna); file kept
% \figtodo{transmission-criteria table, Table~\ref{tab:xmitcriteria}: reference it in the transmission-requirement text; its content is still under review}

\urgentbr
\urgentnote{\textbf{VOCAB}: \textbf{ADD DEFINITION} (new text, insert here): ``We use three terms for three stages. A \emph{violation} is an abstract execution graph that breaks the defense's leakage predicate (Sec.~III). A \emph{hit} is a test whose transmitter the checker sees transmit inside the violation window on the gem5 build. A \emph{leak} is a hit we confirm is real, and a \emph{bug} when the defense should have blocked it.''}
\notebr
What counts as ``a transmission'' is transmitter-specific, and crucially \emph{different transmitter types transmit through different pipeline events, at different pipeline stages, and therefore at different times}. We thus need a distinct signal per transmitter:
\begin{itemize}
  \item \textbf{Load transmitter.} The signal is the load's cache touch --- the tick at which its access reaches the cache
\notebr
\auditnote{fact}{fact}{check\_ld times the load by its latest LSQ `Executing load' tick, which fires before the defense decides whether to send the access; the pipeview complete tick is the fallback. The packet and spec\_read ticks only feed the yes/no data-obtained gate and are never compared with the window. If a multi-micro-op transmitter (a read-modify-write) never executes its load micro-op, the fallback scores the first micro-op's complete tick. The footprint defenses (InvisiSpec, CleanupSpec, SpecLFB) do not use this predicate: their runs use the footprint check over SimSpect cache-event lines with a per-defense criterion, and the memaccess and expose checkers are legacy. \newline \textbf{PROPOSED TEXT} (edits the sentence just before this note): The signal is the tick at which the LSQ executes the load (\texttt{Executing load}; \texttt{O3PipeView} \texttt{complete} if absent), and the load counts only if it then obtained data \ldots; the footprint defenses use their own footprint criteria. \newline \textit{Source:} STAGE3\_gem5/check\_ld.py:10-21, 83-87, 119-122; runconfigs/footprint/run\_config\_STT\_6\_speclfb\_footprint.jsonc:44 (check\_footprint), :68 (legacy\_checker)}
\notebr
--- read from the \texttt{LSQUnit} trace (the \texttt{Executing load} / packet / \texttt{spec\_read} ticks; on newer builds an additive \texttt{memaccess} tick for SpecLFB and an \texttt{expose} tick for InvisiSpec) together with the \texttt{O3PipeView} \texttt{complete} tick. A load counts as transmitting only if it is in-window \emph{and} it actually \textbf{obtained data} --- either it \emph{reached the cache} or it received a \emph{store-to-load forward}. Squash timing never gates a load: one that obtained data in the window has already transmitted even if it is later squashed.
   % This data-obtained gate is essential: store-to-load forwarding lets a speculative load complete and leak \emph{without} a cache packet, so a cache-reach-only signal would miss it.

  

  \item \textbf{Branch transmitter.} The signal is the branch's redirect, keyed on the \emph{squash} tick (the tick at which the misprediction broadcasts a redirect), read from the \texttt{Commit} trace as a pc~$\rightarrow$~squash-tick map. A branch transmits later in the pipeline than a load (at resolution/redirect rather than at execute), which is precisely why it needs its own predicate and its own timing.

  \item \textbf{Faulting / long-latency ``other'' transmitter.} A speculative transmitter that faults (e.g.\ a divide) transmits through its timing: the signal is its completion tick, and a fault after a mispredicted branch is scored as a transmission (Table~\ref{tab:instrtable}, row~5).
\reviewtag{[TODO: Finished this sentence]}
  % ; a speculative divide fault is promoted to a leak, which is what makes a guard such as the \texttt{idivq} guard meaningful.

  \item \textbf{Store transmitter.} A store transmits through its address-dependent effects --- its TLB/cache access and, for the store-to-load-forward channel, the forward it provides to a younger load. A speculative store does not install a cache line until commit, so its in-window signal is the address-dependent access/forward rather than a line install.
\notebr
\reviewnote{There is no store checker: kinds without their own checker fall back to the load check (pipeline.py:313-320), which scores a store by its completion tick. Only our SPT model makes stores transmit (SPT\_6\_oneLP\_v2.als, leakage function), and store transmitters are out of scope for STT (decided 2026-11-08: STT assumes stores act only at retirement). CUT this bullet, or \newline \textbf{PROPOSED TEXT}: ``Store transmitters, used only by our SPT model, are scored by the load predicate's fallback, the store's completion tick.''}
\notebr
\end{itemize}
Because every reported leak must be a transmission inside the window on the modeled path, the predicate is paired with the window gate and the ''branches still unresolved'' gate.
\notebr
\reviewnote{typo: the opening quote should be two backticks.}
\notebr

% , and (for a \emph{demonstrable} leak rather than a dataflow signal) optionally with a data-taint gate requiring the touched address or victim to depend on the tainted operand.

\subsection{The Leakage-prompting $\mu$arch Knobs}
The dimensions the pipeline exposes fall into three classes. Enumeration is \emph{comprehensive} over the \emph{forms of leakage} --- the minimal, necessary leakage backbones the defense recognizes --- so coverage is complete up to the instruction bound. The concretization choices that merely \emph{realize} a backbone --- x86 instruction templates, the two branch-misprediction modes, and the interleaved chains --- are \emph{semi-enumerated}, sampled rather than exhausted. The leakage predicate is fixed: the Alloy predicates that define a violation, and, per campaign, the defense under test. 
\notebr
% --- is fixed, while every knob that \emph{probes} it, from the window stalls to the micro-architectural timing of the fragility sweep, is reported as a swept variable, never tuned to hit or avoid a bug.
\anna{idk}
% (leakage-prompting knobs: written up as a subsection in the Test Generation Pipeline section, sections/pipeline.tex)
% Reorganized 2026-11-10 (Anna's request): each bug is described once, under Known Bugs Replicated or Unknown bugs found.
% The later detail sections ("Bugs found in the gem5 defenses" and its three subsubsections, "vs. AMuLeT: scope") are
% commented out below, unchanged; their content is merged into one PROPOSED TEXT note per bug, and their figures and
% \label{}s moved up to UV2, UV3 and UV6. Redundant older notes were merged into those notes.
\section{Results}
\labbr
\labnote{results are not finalized.}
\notebr
\iffalse % LAB-REVIEW CUT (2026-11-03): placeholder to-do list ([N] numbers)
{\color{cfigtodo}\footnotesize\noindent[fig todo: results to write up (numbers are placeholders)]
\begin{itemize}
\item Re-run every number on the fixed harness (2026-10-30); every number below is a placeholder [N].
\item \textbf{STT}: visibility-point/squash race (Fig.~\ref{fig:race}): the transmitter's cache packet lands [N] cycles after its branch resolves and [N] before the squash; [N] instances; 0 at stock \texttt{commitToIEWDelay}=1. Open: whether this counts under the window (R1); option: report post-resolution, pre-squash transmission as its own category.
\item \textbf{SPT-Fence}: same race (re-run with branch forces); [N] untriaged branch-transmitter hits (post-fix sample: [N] not-taken / [N] taken) are a possible new finding: triage.
\item \textbf{Recon} (STT scheme): getOldestTaint/OTB (Fig.~\ref{fig:otb}) confirmed by a hand-built test; show the cache access with LSQ tracing; Stage~1.5 and interleaving should reach it. Store-to-load forwarding bypass (Fig.~\ref{fig:reconstl}): confirmed (Anna, 2026-11-03); 1 of [N] load hits.
% STAGE NAMES (2026-11-11): in the item above, Stage~1.5 is now Stage~1b (Fig.~pipeline renamed).
\item \textbf{SpecLFB}: UV6 (first speculative load left unprotected): \toolname{} first detects it at test [N] vs.\ AMuLeT's [N]; hit density [N]\% vs.\ [N]\%; footprint criterion in Table~\ref{tab:xmitcriteria}.
\item \textbf{InvisiSpec}: footprint criterion (L1D, L2, directory, D-TLB events of squashed loads; Table~\ref{tab:xmitcriteria}); UV1 L1D eviction on a real trace (Fig.~\ref{fig:checkerwindow}); [N] hits; collateral channels sit outside the defense's own guarantee.
\item \textbf{CleanupSpec}: per-line undo-balance criterion; post-squash residue (UV3--UV5) and unXpec timing; [N] hits; out of scope of the defense's guarantee.
\item \textbf{Targeting}: 100\% of \toolname{} tests carry a protected transmitter vs.\ [N]\% of AMuLeT's random tests (Fig.~\ref{fig:target}); independent of the harness.
\item \textbf{Cost}: per-test gem5 time [N]\,s vs.\ [N]\,s; total simulation to first detection [N]$\times$ less.
\item \textbf{Coverage}: per defense, [N] Stage-1 graphs, [N] Stage-1.5 graphs, [N] programs, [N] runs.
% STAGE NAMES (2026-11-11): Stage-1 graphs -> Stage-1a graphs, Stage-1.5 graphs -> Stage-1b graphs (Fig.~pipeline renamed).
\item Timing sweeps (Fig.~\ref{fig:timing}) predate the fix: re-run or drop.
\end{itemize}}
\fi

We run \toolname{} end-to-end on the gem5 implementations of STT, SPT and ReCon, and on the
AMuLeT artifact builds of InvisiSpec, CleanupSpec, and SpecLFB.
\urgentbr
\urgentnote{\textbf{PROPOSED TEXT} (replaces the sentence just before this note): We run \toolname{} end-to-end on the gem5 implementations of STT and SPT, on ReCon's implementation of STT, and on the AMuLeT artifact builds of InvisiSpec~\cite{yan2018invisispec}, CleanupSpec~\cite{saileshwar2019cleanupspec}, and SpecLFB~\cite{cheng2024speclfb}, which we drive with tests from our STT model. Why: the three AMuLeT defenses are not cited at first mention; our ReCon runs never exercise ReCon's reveal (classic caches, or the authors' Ruby setup with tests that cannot reach the reveal), so they test ReCon's implementation of STT; and the footprint checker counts a footprint left by any squashed load, not only the targeted transmitter (counts in the note under Known Bugs Replicated). \textit{Source:} simspect wiki/state/EVAL\_MANIFEST.md (E6 deviation paragraph; rows E12--E14).}
\iffalse The result is the set of bugs in the simulated defenses seen in Table~\ref{tab:bugs}.\fi% LAB-REVIEW CUT: bug table hidden
\urgentbr

We refer to vulnerabilities discussed in the AMuLeT paper as A-KVn and A-UVn for known and unknown vulnerabilities respectively, and denote unknown vulnerabilities found by SimSpect as UVn.

\subsection{Known Bugs Replicated}
We check each defense only against the leakage its own paper claims to prevent (see the Discussion). AMuLeT~\cite{amulet} reports nine bugs in these defenses; five fall in that scope, and \toolname{} replicates all five. On InvisiSpec~\cite{yan2018invisispec}, a speculative load evicts a live L1D line (A-UV1) in 2{,}995 of 3{,}000 sampled tests. On CleanupSpec~\cite{saileshwar2019cleanupspec}, a speculative read-modify-write's line is not cleaned (A-UV3, next paragraph), a line-crossing access is not cleaned (A-UV4, 38 tests), and the cleanup invalidates a line that an older committed access had used (A-UV5, [1] test). On SpecLFB~\cite{cheng2024speclfb}, the first speculative load is left unprotected (A-UV6) in 1{,}143 and 1{,}142 of 1{,}500 tests per branch direction. The other four lie outside the guarantees we check: InvisiSpec excludes MSHR contention (A-UV2,~\cite[\S VII]{yan2018invisispec}) and the instruction cache (A-KV1, its footnote~2) from its threat model, CleanupSpec's unXpec (A-KV2) is a timing channel, which we do not check, and STT's speculative store that fills the TLB (KV3) needs a store transmitter, which we scope out.
\urgentbr
\urgentnote{\textbf{PROPOSED TEXT (Known Bugs Replicated; merges the earlier notes here).} (1) A-UV3: ``(A-UV3, next paragraph)'' points to the read-modify-write paragraph below, which is commented out. Either restore it, or replace the clause with: ``a speculative read-modify-write's line is not cleaned (A-UV3): when the test's memory is flushed after setup, CleanupSpec leaves the RMW's line installed in 407 of 590 sampled RMW tests, against 0 of 590 without the flush''. Optional, pending your decision: the same RMWs also upgrade lines that older squashed loads brought in to the modified state, and CleanupSpec then skips those loads' cleanup (73 of 590 tests with the flush); add it as A-UV3 reaching further once its dependence on the secret is checked. (2) A-UV6: add after its sentence, ``SpecLFB's ROB walk also unprotects later loads, a separate bug (UV4 below).'' (3) If the commented details paragraph below is restored, fix its A-UV6 cause: the LSQ walk's start one entry after the head never decides the outcome; the ROB walk marks loads as having no older unsafe load up to and including the first unsafe load itself (AMuLeT calls it ``an undocumented optimization in SpecLFB's implementation'',~\cite{amulet} p12 and Fig.~8). Its last sentence becomes: ``On SpecLFB (A-UV6), the walk that marks loads with no older unsafe load includes the first unsafe load itself, so that load is never held (AMuLeT's Fig.~8); this happens in 76\% of our tests.'' (4) Footprint counts include footprints left by any squashed load, not only the transmitter: on an earlier InvisiSpec sample only other squashed loads (mostly the wrong-path function epilogue) leave a footprint in 802 of 2{,}883 hits, and on SpecLFB the transmitter is among the offending loads in 911 of 1{,}300 hits (not-taken) and 937 of 1{,}452 (taken); decide whether to report all squashed-load footprints, as now, or only the transmitter's. Counts refresh from the final run. \textit{Source:} simspect URGENT\_TODO.md (PAPER FIX ``the UV6 cause sentence''); wiki/state/EVAL\_MANIFEST.md (row C1).}
% \subsubsection{Read-modify-write instructions} UV3 needs a read-modify-write (RMW) instruction rather than a plain store. CleanupSpec assumes that a store requests its cache line only when it commits~\cite[\S 4]{saileshwar2019cleanupspec}, so a speculative store never changes the cache, and, as in our other specifications, we do not model speculative stores as transmitters. An x86 RMW such as \texttt{or \$0, (mem)} is different: its load half issues speculatively and requests the line with write permission, and CleanupSpec keeps no record with which to undo that request. Every RMW transmitter our model admits is a load transmitter with the load relabelled (an Alloy check finds all 150 minimal RMW cores of up to four instructions to be relabelled load cores), so we derive an RMW corpus from the load corpus by replacing each transmitter load with an RMW to the same address (124{,}198 tests). One more condition matters: our tests zero their memory before the speculative window, so the RMW's line is already cached and there is nothing to undo. When the test's memory is flushed after that setup, CleanupSpec leaves the RMW's line installed in 407 of 590 sampled tests, against 0 of 590 without the flush. The other defenses protect an RMW's load half like a plain load.
\anna{following paragraph has more details on known vuln probs not necessary commenting out rn}
% Each known bug appears in \toolname{}'s traces with a signature that names its cause. On InvisiSpec (A-UV1), a speculative load evicts a live L1D line, and the victim is written back to the L2; with AMuLeT's fix for this bug enabled (\texttt{--apply-spec-eviction-patch}), the eviction disappears, which confirms the cause. It appears under both of InvisiSpec's threat models (Spectre and Futuristic). On CleanupSpec, A-UV4 appears as a cleanup that is skipped because the squashed access crossed a cache-line boundary. Such accesses arise naturally from stack operations in our tests, so no special stimulus is needed, but the runs only complete with enough L1 miss entries: with the 2 that AMuLeT uses, 62\% of our runs deadlock, so we use 8, the smallest setting without deadlocks. A-UV5 appears once: an older, committed line-crossing load hits a line that a squashed load installed, and CleanupSpec still invalidates that line and restores the evicted one. CleanupSpec's paper requires lines used by older accesses to be kept~\cite[\S 3.4]{saileshwar2019cleanupspec}, but the implementation decides only from the squashed load's own state (\texttt{Sequencer.cc:1300--1346}). On SpecLFB (A-UV6), the first speculative load sits at the head of the load queue, and the walk that flags unsafe loads starts one entry after the head (\texttt{lsq\_unit.cc:877}), so that load is never held; this happens in 76\% of our tests.

\subsection{Unknown bugs found by \toolname{}}
\urgentbr\urgentnote{\textbf{HOW THIS SECTION IS ORGANIZED NOW.} Each bug keeps its current text and gets one PROPOSED TEXT note that merges in the detail from the commented-out ``Bugs found in the gem5 defenses'' section further down (mechanism with file:line, the test shape that reaches it, the count and its control, what we do not claim) plus the items still open. Paste a note over the paragraph above it to accept it. The three trace figures moved here (Fig.~\ref{fig:reconstl} to UV2, Fig.~\ref{fig:otb} to UV3, Fig.~\ref{fig:race} to UV6), and so did the labels sec:bug-stlf, sec:bug-otb and sec:bug-race. Numbering stays UV1--UV6: the walk order counts as a new bug (your decision, 2026-11-09).}
\subsubsection{UV1: STT fails to taint one register}
STT's taint propagation skips any source register with index 16 (\texttt{rob\_impl.hh:229}), a check meant for the zero register that ignores the register class~\cite{yu2019stt}. Floating-point index 16 is the low half of \texttt{xmm4}, so a speculatively loaded secret copied through \texttt{xmm4} and back loses its taint, and the dependent transmitter accesses the cache inside the speculative window. \toolname{} exposes it once its test generator routes values through floating-point registers: of 100 load-to-load tests that STT blocks in their integer form, 98 leak when the accessing load's value passes through \texttt{xmm4}, against 0 through \texttt{xmm3} and 0 in the integer form. The check is still present in the latest upstream STT.
\urgentbr\urgentnote{\textbf{PROPOSED TEXT (UV1):} only one change: ``inside the speculative window'' $\to$ ``inside the violation window'' (the paper's term for the window; same change in UV2 and in Every reported leak below).}
\subsubsection{UV2: ReCon: store-to-load forwarding skips the taint check}\label{sec:bug-stlf}
% [lead for Sec.~\ref{sec:bug-stlf}] (2026-11-10: the detail subsubsection is commented out below; its label moved here)
In ReCon's implementation of STT, a load that receives its value from an older store through store-to-load forwarding returns the data (\texttt{lsq\_unit.cc:2019}) before the STT taint check runs (\texttt{:2079}). A load with a tainted address therefore completes inside the speculative window without touching the cache. \toolname{} catches it because its checker asks whether the transmitter obtained its data, not whether it sent a cache request, and it appears in tests generated from the base model.
\urgentbr\urgentnote{\textbf{PROPOSED TEXT (UV2; merges this lead with the commented-out sec:bug-stlf):} In ReCon's implementation of STT~\cite{recon}, a load's taint is checked only on its cache path. When an older store in the store queue covers the load's address, \texttt{LSQUnit::read()} forwards the store's data (\texttt{lsq\_unit.cc:2019}) and schedules the load's writeback before it reaches the STT taint check (\texttt{:2079}), and the writeback handler checks only the Delay scheme. A speculative load with a tainted address therefore completes with the forwarded value inside the violation window, without touching the cache. STT treats forwarding as an implicit branch whose predicate, that the load aliases no older in-flight store, depends on the load's address, and delays its resolution until that predicate is untainted~\cite[Sec.~6.4.2]{yu2019stt}; it also writes the load's output only after the memory access would have completed~\cite[Sec.~6.5]{yu2019stt}. Under ReCon, the load's completion time instead reveals whether its secret-derived address matches the older store's address, which the attacker can choose. Fig.~\ref{fig:reconstl} shows test inst-014912: load i3 runs under a mispredicted branch, so every address built from its result is tainted; the store i4 and the transmitter load i5 use that address. ReCon's own check reports i5 as tainted and speculative at cycle~5, yet i5 completes at cycle~7, two cycles before its squash; by ReCon's \texttt{LSQUnit::read()}, it took i4's value from the store queue. \toolname{} catches it because its checker asks whether the transmitter obtained its data, not whether it sent a cache request. It appears in tests generated from the base model, and giving each store its own address removes all 80 such hits in 200 generated load tests. \textit{Why the wording:} the trace excerpt is grep-filtered and has no forwarding line, so the text says which part comes from the code. \textit{Source:} simspect project log 2026-11-08 21:30 ``fpr-fanout''; wiki/findings/bug/recon\_cache\_reach\_gate/README.md.}
\input{tikz/reconstl_v2_8pt}
\urgentbr\urgentnote{\textbf{FIGURE (Fig.~\ref{fig:reconstl}), merged from three notes.} (1) Its bar ``Tainted and speculative until the squash'' is the violation window: relabel it ``Violation window (until the squash)'', or add ``(the violation window)'' to the caption. (2) The store i4 is drawn as `? \textless{}- [r0]', the same form as a load: draw it as `[r0] \textless{}- ?' and add to the caption ``the store's address shares r0 by state assignment; the model has no store-to-load edge, so forwarding appears only in the trace''. (3) The trace figures use different letters (Fig.~\ref{fig:race} F/D/E, this one N/E, Fig.~\ref{fig:otb} N/E/C) and window styles; use one letter set (F fetch, D dispatch, I issue, C complete) and one window style (tan fill and frame) in all three. (4) The caption's ``took i4's value'' needs the same code-vs-trace wording as UV2's proposed text.}
\subsubsection{UV3: ReCon: only the oldest taint source}\label{sec:bug-otb}
% [lead for Sec.~\ref{sec:bug-otb}] (2026-11-10: the detail subsubsection is commented out below; its label moved here)
When an instruction merges two tainted operands, ReCon's implementation of STT keeps only the oldest taint source (\texttt{getOldestTaint}) rather than all of them. Once the branch guarding that older source resolves, the merged value counts as untainted although it still depends on a younger speculative load, and a transmitter that uses it accesses the cache. \toolname{} produces it when its test generator adds an older load under an already-resolved branch and merges it into the transmitter's address: 100 of 100 such tests leak, against 0 of 100 without the merge.
\urgentbr\urgentnote{\textbf{PROPOSED TEXT (UV3; merges this lead with the commented-out sec:bug-otb):} ReCon's taint tracker keeps a single taint source per register. When an instruction has several tainted operands, \texttt{propagateTaints} gives its destination the taint of only the oldest one (\texttt{getOldestTaint} picks the lowest sequence number), and \texttt{freeTaints} drops that taint as soon as its source is no longer speculative. A value computed from an older and a younger speculative load therefore counts as untainted once the older load's branch resolves, although the younger load is still speculative. STT instead untaints an instruction's operands only when their youngest root of taint reaches the visibility point~\cite[Sec.~7.1]{yu2019stt}. Fig.~\ref{fig:otb} shows the case: the ALU op i4 adds $r0$, loaded by i1 under branch i0, and $r1$, loaded by i3 under the mispredicted branch i2, and its result is the address of transmitter i5. Branch i0 resolves at cycle~0 and i5 issues at cycle~8, the same cycle as with no defense, while i2 stays unresolved until cycle~92. We found the bug with this hand-built test. \toolname{}'s Stage~2 reproduces it by merging an older load, under an already-resolved branch, into a transmitter's address: 100 of 100 such tests leak, against 0 of 100 without the merge, and in those tests the transmitter completes while the younger branch is still unresolved. In the main evaluation corpus, 11 of the 24 tests per branch direction whose transmitter takes such a merge leak (2\% sample). \textit{Still open:} the last sentence of the Fig.~\ref{fig:otb} caption (``built by hand; our bounded enumeration does not reach this shape'') is out of date, and the hand-built trace shows i5 only issuing (this build never logs a sent cache packet), so cite the generated tests for ``completes''. Fig.~\ref{fig:interleave}'s caption says interleaving builds this shape, but interleaving is off in every evaluation configuration and the OTB merge builds it (pending decision: redraw that figure around the merge, drop its OTB claim, or cut it). \textit{Source:} simspect wiki/design/VOCAB\_TABLES.md ``OTB merge''; URGENT\_TODO.md DECIDE ``OTB on two-stage tests'' (E2 update); SOURCE\_AUDIT\_RECON.md R-SA-8 ``gem5 check''.}
\input{tikz/recon_otb_v2_8pt}
\subsubsection{UV4: SpecLFB: ROB walk order} SpecLFB decides whether a speculative load is unsafe in two walks: the ROB walk marks a load that has no older unsafe load, and the LSQ walk flags loads in a branch's shadow. A load is held only if it is flagged and was not marked. A load reaches the ROB one cycle before it reaches the load queue (\texttt{renameToROBDelay}=1, \texttt{renameToIEWDelay}=2), so the ROB walk can mark a younger load before the LSQ walk has flagged the older one, and the mark is never cleared. Loads after the first speculative load can therefore go unprotected too, so the leak AMuLeT reports for the first load (A-UV6) reaches further than reported: in [157 and 162] of 1{,}500 tests per branch direction, only a later load leaks.
\urgentbr\urgentnote{\textbf{PROPOSED TEXT (UV4; the walk order counts as a new bug, your decision 2026-11-09).} Replace the last sentence with: ``Loads after the first speculative load can therefore go unprotected as well. This is a separate bug from A-UV6, which leaves only the first speculative load unprotected: in 157 and 162 of 1{,}500 tests per branch direction, only a later load leaks, which A-UV6 does not explain.'' Before claiming it is separate, check that AMuLeT's fix for the first-load exemption does not also close it. Keep ``six new ones'' in the abstract and intro. \textit{Source:} simspect URGENT\_TODO.md ``walk-order''; wiki/findings/audits/2026-11-08/SOURCE\_AUDIT\_CLEANUPSPEC\_SPECLFB.md.}
\subsubsection{UV5: SpecLFB: the shadow of a jmp}
% [new subsubsection]. (2026-11-10: marker moved to a comment)
SpecLFB's ROB walk keeps a single flag for whether all older branches have resolved. An unresolved branch clears it, but a younger direct \texttt{jmp} or \texttt{call} sets it again (\texttt{rob.cc:501} in the original artifact and in AMuLeT's copy), erasing the older branch's shadow, so loads after the jump count as safe and fill the L1D speculatively. The logic was ported from DOLMA~\cite{loughlin2021dolma}, which keeps a flag that stays set for the rest of the walk and does not have this bug. We confirm it on an unmodified build with a natural misprediction: the fill depends on the secret, and a one-line change that leaves the flag alone at direct jumps removes it without affecting UV6. It appears in 148 of 1{,}500 tests with the taken branch direction.
\urgentbr\urgentnote{\textbf{PROPOSED TEXT (UV5).} Replace the last sentence with: ``It appears in 148 of 1{,}500 taken-direction tests through the test's closing jump, and in both branch directions once our generator emits direct jumps (134 of 3{,}744 sampled tests). AMuLeT's own SpecLFB campaigns also trigger it, in 47 of their 248 violations, but attribute every one to A-UV6.'' Also: the 148 are tests whose jump is the harness's closing \texttt{jmp}; traced reruns tie a further 274 hits (loads re-sent after the ROB walk cleared them, with a jump in the shadow) to the same reset, and whether to count those is a pending decision. ``Without affecting UV6'' means AMuLeT's A-UV6; say ``A-UV6''. Check the line number against an unmodified copy (the copy we run, which has our hooks, has it at \texttt{rob.cc:538}). \textit{Source:} simspect project log 2026-11-09 20:10 ``bug-investigate-e14new''; wiki/state/EVAL\_MANIFEST.md (row E14; row T3; Footprint model experiment).}
\subsubsection{UV6: the visibility-point/squash race}\label{sec:bug-race}
% [Sec.~\ref{sec:bug-race}, moved last]. (2026-11-10: the detail subsubsection is commented out below; its label moved here)
STT clears a transmitter's protection when its branch resolves, without excluding instructions on the wrong path that is about to be squashed. When squashes take longer to reach issue (\texttt{commitToIEWDelay} $\geq$ [2--3]), the wrong-path transmitter accesses the cache in the gap between resolution and squash ([N] tests; none at stock timing). SPT-Fence shows the same race [pending a re-run on the fixed harness].
\urgentbr\urgentnote{\textbf{PROPOSED TEXT (UV6; merges this lead with the commented-out sec:bug-race):} STT roots a load's taint in \texttt{isAccess() \&\& !isUnsquashable()} and marks an instruction unsquashable once every older branch has resolved (\texttt{isPrevBrsResolved()}). While the mispredicted branch is unresolved, the wrong-path load is squashable, its result is tainted, and the transmitter that uses it is held. When the branch resolves, the younger wrong-path instructions become unsquashable at once, the load's taint drops, and the transmitter issues; but the squash caused by the same resolution reaches the issue stage only \texttt{commitToIEWDelay} cycles later, and in that gap the wrong-path transmitter accesses the cache. STT never requires the squash to land before it releases a transmitter. Fig.~\ref{fig:race} shows test inst-047231 on STT with slowed squash propagation (\texttt{commitToIEWDelay}~$=4$, \texttt{iewToCommitDelay}~$=30$, \texttt{squashWidth}~$=1$; \texttt{commitToIEWDelay}~$=3$ alone also triggers it). The transmitter i5 is ready at cycle~49 but held, because its address comes from the speculative load i3. At cycle~77 the mispredicted branch i2 resolves, i5 issues at cycle~78 and sends a cache packet at cycle~79, and the squash reaches the load queue only at cycle~81. At stock timing (\texttt{commitToIEWDelay}~$=1$) the squash arrives first and no packet is sent. SPT-Fence lifts its fence on the same resolution signal and races the same way: in 24 sampled race tests it leaks in none at stock timing, and in 3, 14 and 15 with \texttt{commitToIEWDelay} of 2, 3 and 4. \textit{Still open:} (1) the STT count: the only one, 23{,}966, predates the 2026-10-30 harness fix and includes 4{,}122 hits never re-checked; run a slowed-timing configuration as its own row, or describe inst-047231 without a count. (2) The window rule (decided 2026-11-08: a branch that has executed but whose squash has not landed is still unresolved) is not yet in the Pipeline section's window definition, and our checker still ends the window at resolution (simspect STAGE3\_gem5/gem5\_common.py, \texttt{resolve\_visible}), so the race is not yet re-scored. (3) The SPT numbers come from the old SPT\_6\_oneLP testset; re-run on the evaluation corpus before reporting them as a rate. \textit{Source:} simspect wiki/findings/bug/stt\_slowcomm\_vp\_squash\_race/ (provenance.json, example\_inst-047231.md); wiki/findings/audits/2026-11-08/SOURCE\_AUDIT\_SPT.md ``A5 gem5 check''; URGENT\_TODO.md DECIDED 2026-11-08.}
\input{tikz/squashrace_v2_8pt}
\urgentbr\urgentnote{\textbf{FIGURE (Fig.~\ref{fig:race}).} Fig.~\ref{fig:race} and Fig.~\ref{fig:checkerwindow} must draw the same violation window: from lc's retirement until fnc retires or a guarding branch's squash lands, whichever comes first. Here it runs to the squash at cycle~81: rename ``Committed window (lc retires to fnc retires)'' to ``Violation window'' and drop ``the committed window closes when fnc retires'' from the caption; Fig.~\ref{fig:checkerwindow} ends its window at the branch's resolution, so extend it to the squash. Also: fnc (i1) is not drawn although the caption ends the window at its retirement; E (issue) is marked at cycle~49 where i5 is held until cycle~78 (label 49 ``ready (held)''); and the caption is missing the slowed-timing settings in UV6's proposed text.}

\subsection{Comparison with prior work} Compared to AMuLeT, \toolname{} needs [N]$\times$ fewer tests on average to find the known bugs ([N] per bug, on the fixed harness). Part of the difference is targeting: every \toolname{} test contains a protected transmitter, against 17\% of AMuLeT's randomly synthesized tests.
\input{tikz/firstdet_bars_v2_8pt}
\urgentbr\urgentnote{\textbf{PROPOSED TEXT (Comparison with prior work; preliminary, numbers refresh from the final run; replaces the paragraph above).} We compare the number of tests each tool needs to find each previously known bug. The two tools look for different leaks: AMuLeT fuzzes for collateral channels by comparing executions on different inputs, while \toolname{} checks the transmitter's own transmission inside the violation window. For \toolname{}, a test is one generated program run once: we generate one corpus from the footprint model and run it on all three AMuLeT builds at AMuLeT's own configurations, drawing the cache state (warm, flushed or filled) and the branch-stall timing at random for each test, and report the expected number of tests to the first detection, $(N{+}1)/(K{+}1)$ for $K$ detections in $N$ tests. For AMuLeT we run its own fuzzer with its own settings (32 instances of 200 programs with 140 inputs each) and count test cases, one program on one input, per detection of that bug over all instances. \toolname{} needs at least $18\times$ fewer tests for every previously known bug, and $140\times$ to $2{,}100\times$ fewer for all but A-UV5, the rarest for both tools: \toolname{} detects it four times in all 50{,}232 tests of its corpus and AMuLeT four times in 896{,}000 test cases (95\% interval for the ratio, with both sides' uncertainty, $3.3\times$ to $96\times$; both counts use the same estimator, $(N{+}1)/(K{+}1)$) (Fig.~\ref{fig:firstdetv2}). Part of the difference is targeting: every \toolname{} test contains a transmitter in the protection set by construction, while 17\% of AMuLeT's randomly synthesized tests contain a speculative gadget inside the window. \textit{Notes:} (1) ``by construction'' matters: these are different criteria, and a test with a protected transmitter need not leak (STT gives 0 hits on our default tests); apply one criterion to both, or keep ``by construction''. (2) Do not use the June SpecLFB comparison (test \#6 vs.\ \#39, ``6.5$\times$''): it predates the harness fix and compared our checker's hits with AMuLeT's contract violations. (3) The jmp-shadow count for AMuLeT's own runs belongs in UV5. \textit{Source:} simspect wiki/state/EVAL\_MANIFEST.md, Footprint model experiment and rows T1--T5.}

\subsection{Timing and CPU usage}\label{sec:timing}
\urgentbr\urgentnote{\textbf{PROPOSED TEXT (timing; new paragraph, preliminary numbers).} We measure cost in CPU time, running one gem5 simulation per test at AMuLeT's microarchitectural configuration for each defense on a dedicated 64-core machine. A \toolname{} test takes 0.45 to 1.17 CPU-seconds of simulation and 1.6 to 3.7 CPU-seconds in total, including assembling, linking and checking the test, so one core runs about 1{,}000 to 2{,}200 tests per hour; the range comes from the defense and from the cache state each test starts with. Because every test targets a forbidden execution, a first detection comes quickly: on SpecLFB, \toolname{} first detects A-UV6 at test 47 when tests start with a warm cache (77 CPU-seconds in total) and at test 1 when they start with a flushed cache (3 CPU-seconds); on CleanupSpec it first detects A-UV4 at test 47 (130 CPU-seconds). Measured as total CPU time per detection, \toolname{} is [19--37]$\times$ cheaper than AMuLeT for A-UV1, A-UV4, A-UV6 and the jmp shadow, while AMuLeT finds A-UV3 [1.6--2.8]$\times$ faster, because only some of our tests contain a read-modify-write. \textit{Notes:} the per-detection multiples come from runs on a shared machine, so \toolname{}'s CPU times there are upper bounds; AMuLeT's paper reports a rate (campaign time divided by all violations, Table~4 p9), not a time to first detection, so compare with our measured AMuLeT runs, not with the paper. Refresh every number from the final unified footprint corpus. \textit{Source:} simspect wiki/state/EVAL\_MANIFEST.md, ``Timing results: SimSpect at AMuLeT parity'' (rows T2, T2c, T3, T3c) and ``SimSpect vs AMuLeT, CPU time per detection, per bug''.}
\paragraph{CPU usage}
\labbr\labnote{To be filled: CPU-hours for the whole evaluation per defense, split by stage (Stage~1a Alloy enumeration, Stage~1b flesh-out, Stage~2 concretization, Stage~3 gem5 and checking), peak memory per job, and wall time on the machines used. Where the numbers will come from: the per-test CPU records written by the timing driver (\texttt{results.jsonl}) and the evaluation rows in simspect wiki/state/EVAL\_MANIFEST.md.}

% Bug table: tikz/bugtable_v3_8pt (working table); refresh its rows before showing it (row updates as of 2026-11-09 are in the comments at the top of that file).
% The older table below, under \iffalse, is superseded by bugtable_v3_8pt.
% \input{tikz/bugtable_v3_8pt}% every confirmed and candidate bug, one spot (2026-11-03); uncomment to show
\iffalse
\begin{table*}[t]\centering\small
\setlength{\tabcolsep}{5pt}\renewcommand{\arraystretch}{1.2}
\begin{tabular}{@{}p{2.3cm}p{4.7cm}p{3.0cm}p{3.1cm}cl@{}}
\toprule
\textbf{Defense} & \textbf{Bug} & \textbf{First known from} & \textbf{Found by} & \textbf{In scope?} & \textbf{Verdict}\\
\midrule
\multicolumn{6}{@{}l}{\textbf{STT}}\\
& Visibility-point\,$\leftrightarrow$\,squash race (load channel) & This work & \textcolor{cclean}{\toolname{}} & $\checkmark$ & \textcolor{ctarget}{real}\\
& Pending-squash race --- secret-dependent squash signal & N.~Mosier (\texttt{cwfletcher/stt}\,@18e304f) & \textcolor{gray}{\textbf{neither} tool} & $\sim$ & \textcolor{ctarget}{real (fixed)}\\
& Tainted spec \emph{store} $\to$ D-TLB entry (KV3) & DOLMA; via AMuLeT & \textcolor{camulet}{AMuLeT} & $\times$ store/TLB & \textcolor{copen}{not a bug${}^{1}$}\\
\addlinespace[2pt]
\multicolumn{6}{@{}l}{\textbf{SPT-Fence}}\\
& Same visibility-point\,$\leftrightarrow$\,squash race (= STT) & This work & \textcolor{cclean}{\toolname{}} & $\checkmark$ & \textcolor{ctarget}{real}\\
\addlinespace[2pt]
\multicolumn{6}{@{}l}{\textcolor{gray}{\textbf{InvisiSpec \& CleanupSpec} \,{\footnotesize\itshape collateral transmitter-side channels --- not comprehensive-defense logic bugs, all out of \toolname{}'s scope}}}\\
& \textcolor{gray}{UV1, UV2, KV1 --- L1D eviction / MSHR contention / L1I spec-fetch} & \textcolor{gray}{AMuLeT, known} & \textcolor{camulet}{AMuLeT} & \textcolor{gray}{$\times$} & \textcolor{gray}{real}\\
& \textcolor{gray}{UV3--UV5 --- post-squash L1D residue / over-clean evicts non-spec line} & \textcolor{gray}{AMuLeT} & \textcolor{camulet}{AMuLeT} & \textcolor{gray}{$\times$} & \textcolor{gray}{real}\\
\addlinespace[2pt]
\multicolumn{6}{@{}l}{\textbf{SpecLFB}}\\
& First speculative load left unprotected (UV6) & AMuLeT & \textbf{Both} (\textcolor{camulet}{AMuLeT}@t39; \textcolor{cclean}{\toolname{}}@t6) & $\checkmark$ & \textcolor{ctarget}{real}\\
\bottomrule
\end{tabular}
\caption{Speculative-execution defense bugs, grouped by defense: where each was first
known, which tool exhibits it (\textcolor{cclean}{\toolname{}} / \textcolor{camulet}{AMuLeT}
/ Both / neither), whether it falls in \toolname{}'s leakage model, and the verdict. The
pending-squash race (a secret-dependent squash signal, fixed upstream by Mosier) is a real STT bug
that \emph{neither} automated tool found. $^{1}$not a bug: a tainted-store$\to$TLB leak is consistent behavior with STT's
specification, which does not recognize store transmitters, so it is outside the property STT claims to
enforce. The visibility-point\,$\leftrightarrow$\,squash race is a genuine logic bug: STT and SPT
release the transmitter's protection on the branch-resolution signal but do \emph{not} wait for the
squash to \emph{complete}, so the wrong-path transmitter reaches the cache in the gap before it is
killed.}
\label{tab:bugs}
\end{table*}
\fi

% ===== COMMENTED OUT 2026-11-10 (Anna): "Bugs found in the gem5 defenses" and its three subsubsections. =====
% Their content is merged into the PROPOSED TEXT notes under UV2, UV3 and UV6 above; the three \input{} figures
% and the labels sec:bug-race, sec:bug-stlf, sec:bug-otb moved up there. The text is kept here unchanged.
% \subsection{Bugs found in the gem5 defenses}
% The confirmed bugs fall into three mechanism classes: a
% \emph{visibility-point\,$\leftrightarrow$\,squash timing race}, shared by STT and SPT-Fence; a \emph{taint-check bypass through store-to-load forwarding} in ReCon; and the
% \emph{unprotected-first-load} exemption in SpecLFB (UV6).
%  In the STT race, the defense clears a
% transmitter's protection at branch \emph{resolution} with no guard excluding the instruction on the
% to-be-squashed wrong path, so the wrong-path transmitter touches the cache in the gap before the
% squash drains; we observe \textbf{23{,}966} race instances and \textbf{0 at stock timing}, where the
% defense holds.
% SPT-Fence exhibits the \emph{same} race --- one shared logic flaw across two defenses ---  we find no independent SPT bug.
% SpecLFB's UV6 leaves the
% first speculative load in the LSQ unprotected so that it performs its in-window cache access,
% demonstrated on \textbf{68/300} tests. We \emph{do not} find STT's tainted-store\,$\to$\,TLB bug, by
% design: \toolname{}'s leakage contract contains no store transmitter, matching STT's own contract of
% load and branch transmitters.
%
% On ReCon's gem5 implementation of STT~\cite{recon}, \toolname{} also exposes a bug in its taint tracking: a load whose address is tainted completes through store-to-load forwarding (Sec.~\ref{sec:bug-stlf}). A hand-built test in our notation shows a second: an instruction that merges two tainted operands keeps only the older operand's taint, so its result is untainted while it still depends on a speculative load (Sec.~\ref{sec:bug-otb}). We describe each bug's mechanism below.
%
% \subsubsection{The visibility-point\,$\leftrightarrow$\,squash race}% \label{sec:bug-race} moved to UV6
% \reviewtag{[TODO: Made this a subsubsection under Bugs found]}
% STT roots a load's taint in \texttt{isAccess() \&\& !isUnsquashable()} and sets
% \texttt{isUnsquashable} $\Leftrightarrow$ \texttt{isPrevBrsResolved()}. While the controlling
% mispredict is unresolved, the wrong-path producer load is squashable, so its result is tainted and
% the transmitter is held, and the defense works. The instant the branch \emph{resolves},
% \texttt{isPrevBrsResolved} propagates to the younger wrong-path instructions: the producer becomes
% unsquashable, its taint drops, the transmitter's address reads as ``clean,'' and the transmitter
% re-issues and touches the cache. But that same resolution is what \emph{squashes} the wrong-path
% load, and the squash is not instantaneous --- it propagates commit$\to$IEW over
% \texttt{commitToIEWDelay} cycles. In the gap between resolution (taint cleared) and squash (load
% killed), the transmitter performs a real cache access (\iffalse Fig.~\ref{fig:trace}, \fi Fig.~\ref{fig:race});% LAB-REVIEW CUT: old trace figure is in the hidden pending section
%
% STT never enforces squash-before-expose.
%
% \reviewtag{[TODO: Proposed this example paragraph, from the figure; check it]}
% Fig.~\ref{fig:race} shows the race on test inst-047231, on an STT build with \texttt{commitToIEWDelay}~$=4$.
% The transmitter i5 is ready to issue at cycle~49 but is held, because its address comes from the speculative load i3. At cycle~77 the mispredicted branch i2 resolves; STT immediately stops treating i3's result as tainted, so i5 issues at cycle~78 and sends a cache packet at cycle~79.
% The squash caused by the same resolution reaches the load queue only at cycle~81, two cycles too late. At stock timing (\texttt{commitToIEWDelay}~$=1$) the squash arrives first and no packet is sent.
%
% % figure fig:trace moved to sections/pendingfigures.tex (not yet redesigned), 2026-11-02
% (\input{tikz/squashrace_v2_8pt} moved to UV6)
%
% \subsubsection{ReCon: store-to-load forwarding skips the taint check}% \label{sec:bug-stlf} moved to UV2
% \reviewtag{[TODO: Proposed this subsubsection; check it]}
% ReCon's STT implementation checks a load's taint only on the load's cache path. When an older store in the store queue covers a load's address, \texttt{LSQUnit::read()} takes the store-to-load forwarding path instead: it copies the store's data, schedules the load's writeback, and returns before it reaches the taint check that would delay a tainted load, and the writeback handler checks only the Delay scheme, not STT. A speculative load whose address is tainted therefore completes with the forwarded value whenever an older in-flight store has the same address. Fig.~\ref{fig:reconstl} shows this on test inst-014912. Load i3 runs speculatively under a mispredicted branch, so its result, and every address built from it, is tainted; the store i4 and the transmitter load i5 use the same tainted address. ReCon's own check reports i5 as tainted and speculative at cycle~5, yet i5 takes i4's value from the store queue and completes at cycle~7, two cycles before its squash.
% STT treats store-to-load forwarding as an implicit branch whose predicate, that the load aliases no older in-flight store, depends on the load's address, and it delays the branch's resolution until that predicate is untainted~\cite[Sec.~6.4.2]{yu2019stt}. With STT's optimization the load issues only once its address is untainted, and its output is written only after the memory access completes, so that younger instructions see the same timing as without forwarding~\cite[Sec.~6.5]{yu2019stt}. Under ReCon, the load's completion time instead reveals whether its secret-derived address matches the older store's address, which the attacker can choose.
% (\input{tikz/reconstl_v2_8pt} moved to UV2)
%
% \subsubsection{ReCon: \texttt{getOldestTaint} drops a younger taint}% \label{sec:bug-otb} moved to UV3
% \reviewtag{[TODO: Proposed this subsubsection; check it]}
% ReCon's taint tracker records a single taint source per register. When an instruction has several tainted source operands, \texttt{propagateTaints} gives its destination the taint of only the oldest one (\texttt{getOldestTaint} picks the lowest sequence number), and \texttt{freeTaints} drops a register's taint as soon as that source is no longer speculative. A value computed from an older and a younger speculative load is therefore untainted when the older load's branch resolves, even though the younger load is still speculative. STT instead untaints an instruction's operands only when its youngest root of taint, the youngest access instruction its operands depend on, reaches the visibility point~\cite[Sec.~7.1]{yu2019stt}. Fig.~\ref{fig:otb} shows the case: the ALU op i4 adds $r0$, loaded by i1 under branch i0, and $r1$, loaded by i3 under the mispredicted branch i2, and its result is the address of transmitter i5. Branch i0 resolves at cycle~0, and i5 issues at cycle~8, the same cycle as with no defense, while i2 stays unresolved until cycle~92. We confirmed this with a hand-built two-branch test: our bounded Alloy enumeration does not reach this shape within its instance cap, and Stage~2 interleaving is meant to build it (Fig.~\ref{fig:interleave}).
% (\input{tikz/recon_otb_v2_8pt} moved to UV3)
% ===== end of commented-out "Bugs found in the gem5 defenses" =====

\labbr
\labnote{InvisiSpec, SpecLFB and CleanupSpec results still to be inserted here.}
\notebr

\iffalse % LAB-REVIEW CUT (2026-11-03): placeholder subsubsection, not written yet
\subsubsection{Defenses with their own load transmission mechanisms}\label{sec:bug-loadmech}
\reviewtag{[TODO placeholder: InvisiSpec, SpecLFB (UV6) and CleanupSpec hide or undo a speculative load's effect instead of delaying the load, so \toolname{} checks each with a defense-specific load transmission criterion. Describe each bug's mechanism here, with its criterion and its figure (the checker-window figure, Fig.~\ref{fig:checkerwindow}, is the InvisiSpec UV1 example).]}
\fi

% ===== COMMENTED OUT 2026-11-10 (Anna): "vs. AMuLeT: scope" contradicts Known Bugs Replicated (eight vs nine bugs; "exactly one" in scope).
% Its one-line contrast between the two tools' leakage models is folded into the Comparison with prior work note above. Text kept unchanged:
% \subsection{\toolname{} vs.\ AMuLeT: scope}
% The two tools use different leakage models. AMuLeT detects \emph{collateral shared-resource} leaks
% --- cache eviction, MSHR, L1I, TLB --- by input-differential fuzzing, whereas \toolname{} detects the
% \emph{transmitter's own transmission inside the speculative window
% } by targeted enumeration. Of the eight
% bugs AMuLeT reports on defenses, exactly one (SpecLFB UV6) falls in \toolname{}'s leakage model\iffalse (the ``in scope'' column of Table~\ref{tab:bugs})\fi; % LAB-REVIEW CUT: bug table hidden
% the rest are collateral shared-resource channels
% --- eviction, MSHR, I-cache, TLB, and post-squash residue --- outside the pipeline trace \toolname{}
% observes and outside its dataflow stimulus.
%
% Conversely, AMuLeT's random search does not target the
% transmitter-in-window timing races (the STT/SPT VP\,$\leftrightarrow$\,squash race) that \toolname{}
% reproduces deterministically.
% ===== end of commented-out "vs. AMuLeT: scope" =====

\iffalse % LAB-REVIEW CUT (2026-11-03): Anna: fixing this still; its figure is in the hidden pending section
\subsection{Test targeting} \anna{fixing this still}
Of AMuLeT's 500 random STT tests, only \textbf{17\%} contain a window-valid speculative transmitter
gadget; broken down by transmitter type, just \textbf{14\%} carry a leaky load and \textbf{5\%} a
leaky branch, and the rest cannot leak by construction (Fig.~\ref{fig:target}). By contrast,
\textbf{100\%} of \toolname{}'s tests carry a transmitter provably in the protection set, because
each test is enumerated \emph{to} place one there.
\fi


% figure fig:target moved to sections/pendingfigures.tex (not yet redesigned), 2026-11-02

\iffalse % LAB-REVIEW CUT (2026-11-03): sweeps predate the 2026-10-30 harness fix (re-run or drop); figure hidden
\subsection{Timing-window and fragility sweeps}
The VP\,$\leftrightarrow$\,squash race is \emph{timing-regime}: STT holds at stock squash latency
(\texttt{commitToIEWDelay}$=1$, 0 leaks) and leaks from $\geq 2$ (Fig.~\ref{fig:timing}a).
These are
micro-architectural \emph{implementation} parameters, not part of any defense's security argument, so
sweeping them tests whether the defense's \emph{logic} survives the legal timing space --- and it does
not. Three micro-architectural \emph{stall} knobs separately probe the speculative window; of these
only the unresolved-branch stall changes the verdict,
and it \emph{masks} the leak (never manufactures
one), while the resolved-branch stall is a flat plateau and the leak survives with the fnc stall off
(Fig.~\ref{fig:timing}b).
\urgentbr
\urgentnote{These stall-knob conclusions come from pre-fix sweeps: the old resolve-stall hook froze runs when the CPU slept, and the old BTB force was wrong on every build. So "masks the leak, never manufactures one" and "flat plateau" are known-unreliable results stated as fact. The 27-to-0 drop at s\_U $\geq$ 50 may be that freeze rather than masking, and s\_R (27 to 22) and the fnc stall (16 to 27) also change the hit count. \newline \textbf{PROPOSED TEXT} (edits the sentence just before this note): Drop the stall-knob sentence, or re-run it on the fixed harness first. \newline \textit{Source:} STATUS.md:28-30; project log:51-58; wiki/findings/bug/BRANCH\_ANNOTATION\_AND\_FORCE.md}
\notebr
\fi


% figure fig:timing moved to sections/pendingfigures.tex (not yet redesigned), 2026-11-02

% Superseded by "Timing and CPU usage" above (2026-11-10); kept hidden.
\iffalse % LAB-REVIEW CUT (2026-11-03): numbers predate the harness fix; figure hidden
\subsection{Detection efficiency vs.\ AMuLeT}
On the SpecLFB build --- the one bug both tools find (UV6) --- \toolname{} flags it at test
\textbf{\#6} versus AMuLeT's first contract violation at \textbf{\#39}: $6.5\times$ fewer tests, at
$\sim$7$\times$ the hit density (23\% vs.\ 3.3\%).
Isolating just the gem5 simulation
(excluding \toolname{}'s build and its measurement trace), the per-test cost is comparable (0.092
vs.\ 0.057\,s), so the per-test gem5 work is the same and \toolname{}'s targeting is what yields
$\sim$6.5$\times$ less total simulation to reach the bug (Fig.~\ref{fig:eff}).
\fi


% figure fig:eff moved to sections/pendingfigures.tex (not yet redesigned), 2026-11-02

\subsection{Every reported leak is gem5-validated}
\notebr
Every count in this section is a \emph{confirmed} real leak:
a transmitter that obtained its data
inside the speculative window
on the running defense build, decided by the per-transmitter checker
--- a real cache packet, or a store-forward that completes --- never a static or model-only signal,
and squash timing never disqualifies a load that already obtained its data.
% (2026-11-10: the closing list below is commented out; it repeated the bug subsections with pre-fix numbers, 23,966 and 68/300, and left out most bugs)
% On that bar, the
% leaks are
% the STT/SPT visibility-point\,$\leftrightarrow$\,squash race
% (\textbf{23{,}966} race instances; \textbf{0} at stock timing) and SpecLFB UV6 (\textbf{68/300}), and ReCon's taint-check bypass through store-to-load forwarding (inst-014912).
% \reviewtag{[TODO: Added ReCon here]}
\urgentbr\urgentnote{\textbf{PROPOSED TEXT (Every reported leak; merges three notes).} This bar does not hold yet, so reword it: ``Every load count in this section comes from the per-transmitter checker on the running defense build, and every footprint count from the defense's footprint criterion; counts marked $^\dagger$ come from earlier runs and are being repeated.'' Also change ``inside the speculative window'' to ``inside the violation window''. Why: (i) the A-UV1, A-UV3--A-UV6 and jmp-shadow counts come from the footprint checker, which counts a footprint left by any squashed load; (ii) the data-obtained gate asks whether any instance at the transmitter's PC sent a packet or completed, not whether this instance did inside the window (only the load's execute tick is tested against the window); (iii) on ReCon's builds the checker cannot see a sent packet at all (they never log one), so only loads that complete count there. The closing list is commented out: end the subsection here and point to the two bug subsections. \textit{Source:} simspect URGENT\_TODO.md DECIDE ``check\_ld's data-obtained gate is not tied to the window''; wiki/findings/bug/recon\_cache\_reach\_gate/README.md.}

\section{Related Work}
\notebr
\auditnote{cite}{related}{To cite background from here, point back to it instead of repeating it: e.g.\ ``Building on the axiomatic models of Sec.~\ref{sec:bg-axiomatic}, \ldots'' or ``Unlike the contract-based validators of Sec.~\ref{sec:bg-litmus}, \ldots''. Keep each paper's detailed description in one place.}
\notebr
\auditnote{cite}{related}{\textbf{RELATED-WORK PASS}: this is still bullets; write it as prose, about 0.75 page (1.5 columns). Proposed five parts, each ending with one sentence on how SimSpect differs: \textbf{Triaged 2026-11-03}: about 33 papers are kept as CITE notes below; the other 82 are in a commented CUT LIST at the end of each part (bib entries kept).}
\notebr
\auditnote{cite}{related}{(1) \emph{Validating defense implementations}: AMuLeT, Revizor, Hide and Seek, SpecDoctor, IntroSpectre, Scam-V, and SpecCheck (PACT 2023), the closest gem5-based tool, which you must cite and differentiate (a register-based FSM monitor in gem5 that flags potential transient data leakage from a generic attack description, not from a defense's own guarantee); at RTL, UPEC and Contract Shadow Logic; concurrent model checking of defense designs (S\&P 2026). SimSpect: a simulator prototype checked against the defense's own guarantee, with directed enumeration rather than random programs.}
\notebr
\auditnote{cite}{related}{(2) \emph{Leakage contracts and speculation security definitions}: Guarnieri et al.\ S\&P 2021, Spectector, Cheang et al., Cauligi et al.'s SoK (one cite standing in for the software-side formal work), LCM, SynthLC, LEAVE. SimSpect: our leakage contract is the transmitter component, as in SynthLC; Guarnieri's contract is an observation clause plus an execution clause. Cite trap: the bib key \texttt{synthlc} is Wang et al., CCS 2025, not SynthLC; SynthLC (Hsiao et al.) is in the RTL2M$\mu$PATH paper, so cite it as \cite{rtl2upath}.}
\notebr
\auditnote{cite}{related}{(3) \emph{Axiomatic memory models and litmus synthesis}: Herding Cats, Wickerson et al., RTLCheck; for litmus synthesis, Lustig et al., CheckMate, Trippel et al.'s MeltdownPrime/SpectrePrime, TransForm. The rest of the MCM lineage (the other Check tools, Simner et al.'s Arm semantics, x86-TSO) is proposed for Background (Sec.~\ref{sec:bg-axiomatic}); point back to it here. SimSpect: MCM-style minimal litmus enumeration brought to speculation-security defenses.}
\notebr
\auditnote{cite}{related}{(4) \emph{Execution models of speculation}: The Code That Never Ran, Cats vs.\ Spectre, InSpectre, Pensieve, RTL2MuPATH. SimSpect: an abstract snapshot of $\mu$arch state (resolved/committed flags) rather than a full operational semantics.}
\notebr
\auditnote{cite}{related}{(5) \emph{Spectre hardware defenses}: the defenses you test, and the classes of the rest (delay or taint-track, invisible speculation, undo, data-oblivious; their cites are in the intro's second paragraph).}
\notebr
% Bullets for now -- to be written into prose.

\subsection{Validating speculative-execution defenses}
\notebr
\auditnote{write}{writing}{Structure problems. Related Work is still bullet points, and AMuLeT is described four times (intro, II-B, V-B, VI-A), with the intro version misdescribing it (AMuLeT runs one program on two inputs with equal contract traces). The Results split the AMuLeT comparison across V-B, V-C and V-E, and V-F repeats the load gate from IV-C3. The paper has no Conclusion, and the Acknowledgments placeholder must go for double-blind review. \newline \textbf{PROPOSED TEXT} (edits the sentence just before this note): Write VI as prose. Keep one AMuLeT description in VI-A (move II-B there). Merge V-B, V-C and V-E into 'Comparison with AMuLeT' and fold V-F into V-A. Add a short Conclusion and delete the acknowledgments for submission. \newline \textit{Source:} related.tex:2 '\% Bullets for now'; intro.tex:21, background.tex:6, results.tex:122-138; fu-2025-amulet.pdf Def. 2.1; bare\_conf\_compsoc.tex:601.}
\notebr
\begin{itemize}\setlength{\itemsep}{1pt}
\item \textbf{AMuLeT}~\cite{amulet} --- the only prior automated validator of Spectre-defense
\notebr
  simulator prototypes; randomized differential testing over architecturally-equivalent
  program/input pairs, flagging a violation when speculative execution
\urgentbr
\urgentnote{\textbf{VOCAB}: ``flagging a violation'' $\to$ ``flagging a \emph{contract violation}'' (AMuLeT's own term, its Def.~2.1), and the same for ``every detected violation'' a few lines below, so neither reads as our model-level violation.}
\notebr
leaks more than
  architectural execution. Limitations \toolname{} addresses:
  \begin{itemize}\setlength{\itemsep}{0pt}
  \item generation is random, so it is agnostic to whether a test even contains the
    ingredients for Spectre leakage --- wasting runs and requiring manual root-causing of
    every detected violation;
  \item it is not specialized to common defense techniques or the bug classes implementing
    them invites, which targeted generation could pursue;
  \item it does not exploit the white-box simulator to bias scheduling toward likely-buggy
    windows;
  \item it conflates verifying that a defense prevents speculative leakage with discovering
    new unsafe instruction operands;
  \item it tests against a sequential model, so it cannot test defenses that protect data
    speculatively even after it has leaked architecturally~\cite{protean}.
  \end{itemize}
\end{itemize}
\notebr
\auditnote{cite}{related}{\textbf{CITE}~\cite{fadiheh2019upec}: UPEC formal RTL detection of transient-execution leaks; formal verification counterpart to SimSpect's simulation-based validation.}
\notebr
\auditnote{cite}{related}{\textbf{CITE}~\cite{hur2022specdoctor}: Differential fuzzing of RTL CPUs for transient-execution bugs; fuzzing alternative to SimSpect's model-guided tests.}
\notebr
\auditnote{cite}{related}{\textbf{CITE}~\cite{ghaniyoun2021introspectre}: Pre-silicon gadget fuzzing + RTL simulation to find transient execution vulnerabilities; close in spirit (simulation-based validation).}
\notebr
\auditnote{cite}{related}{\textbf{CITE}~\cite{nemati2020scamv}: Scam-V: validates abstract side-channel models against hardware via generated test programs; directly analogous model-vs-implementation test generation.}
\notebr
\auditnote{cite}{related}{\textbf{CITE}~\cite{mckevitt2023speccheck}: SpecCheck describes a generic transient-execution attack as a register-based FSM and implements it in gem5 to flag potential transient data leakage (mckevitt-2023-speccheck.pdf abstract). It is the closest gem5-based tool; the difference to state is that it checks a generic attack pattern, while \toolname{} tests a defense's simulated implementation against that defense's own guarantee.}
\notebr
\auditnote{cite}{related}{\textbf{CITE}~\cite{dong2026interplay}: Model checking of secure-speculation processor designs (S\&P 2026); formal verification of defenses, recent concurrent work.}
\notebr
\auditnote{cite}{related}{\textbf{CITE}~\cite{amulet} (already in the bib): Closest competitor: fuzz-tests gem5 secure-speculation defenses against contracts; SimSpect instead derives tests from axiomatic defense models.}
\notebr
\auditnote{cite}{related}{\textbf{CITE}~\cite{revizor} (already in the bib): Contract-based black-box CPU fuzzing; SimSpect replaces random generation with model-directed generation.}
\notebr
\auditnote{cite}{related}{\textbf{CITE}~\cite{oleksenko2023hideseekspectresefficient} (already in the bib): Efficient contract-based fuzzing (Revizor follow-up).}
\notebr
\auditnote{cite}{related}{\textbf{CITE}~\cite{contractshadow} (already in the bib): Formal RTL verification of secure speculation.}
\notebr
\auditnote{cite}{related}{\textbf{CUT LIST} (2026-11-03 triage): 14 more candidates for part (1) validating defenses are in the source comment just below this note. Their bib entries stay, so any of them can come back.}
\notebr
% ---- CUT from part (1) validating defenses in the 2026-11-03 triage (bib entries kept). To restore one, delete its leading "% ". ----
% \auditnote{cite}{related}{\textbf{CITE}~\cite{fadiheh2020formal}: UPEC-OO for out-of-order processors; formal transient-execution vulnerability detection in RTL.}\notebr
% \auditnote{cite}{related}{\textbf{CITE}~\cite{fadiheh2023exhaustive}: Exhaustive UPEC-based detection of transient side channels in RTL (IEEE TC).}\notebr
% \auditnote{cite}{related}{\textbf{CITE}~\cite{hur2021difuzzrtl}: Coverage-guided differential RTL CPU fuzzing; background for hardware fuzzing.}\notebr
% \auditnote{cite}{related}{\textbf{CITE}~\cite{buiras2021observation}: Scam-V follow-up validating side-channel models via observation refinement (MICRO 2021).}\notebr
% \auditnote{cite}{related}{\textbf{CITE}~\cite{weber2021osiris}: Fuzzing for microarchitectural timing side channels; transmitter-discovery related work.}\notebr
% \auditnote{cite}{related}{\textbf{CITE}~\cite{moghimi2020medusa}: Automated attack synthesis (Transynther) for MDS-style leaks; automated exploit generation related work.}\notebr
% \auditnote{cite}{related}{\textbf{CITE}~\cite{rajapaksha2023sigfuzz}: Fuzzing framework for microarchitectural timing side channels in RTL.}\notebr
% \auditnote{cite}{related}{\textbf{CITE}~\cite{oleksenko2020specfuzz}: Speculation-exposing fuzzing of software for Spectre gadgets; software-side testing contrast.}\notebr
% \auditnote{cite}{related}{\textbf{CITE}~\cite{wang2021oo7}: Binary analysis to detect and patch Spectre gadgets; software-side detection.}\notebr
% \auditnote{cite}{related}{\textbf{CITE}~\cite{solt2024cascade}: Cascade: CPU fuzzing via valid program generation; hardware fuzzing related work.}\notebr
% \auditnote{cite}{related}{\textbf{CITE}~\cite{xiao2020speechminer}: Framework to measure speculative execution vulnerabilities on real CPUs.}\notebr
% \auditnote{cite}{related}{\textbf{CITE}~\cite{zagieboylo2023specverilog}: SpecVerilog: IFC type system for secure speculation in HDL; design-time verification of defenses.}\notebr
% \auditnote{cite}{related}{\textbf{CITE}~\cite{weber2026trevex}: Black-box detection of data-flow transient execution vulnerabilities (S\&P 2026); recent testing work.}\notebr
% \auditnote{cite}{related}{\textbf{CITE}~\cite{blackbox} (already in the bib): Background on formal verification of secure processors.}\notebr

\subsection{Formal models of secure speculation}
\begin{itemize}\setlength{\itemsep}{1pt}
\item Hardware--software leakage contracts and speculative non-interference
  ~\cite{guarnieri2021contracts}
\urgentbr
\urgentnote{\textbf{VOCAB}: ``Hardware--software leakage contracts'' $\to$ ``Hardware--software contracts'' (Guarnieri et al.'s own name; you use leakage contract for the transmitter component, as SynthLC does). Also cite Spectector~\cite{guarnieri2020spectector} for speculative non-interference here: Spectector defines it (guarnieri-2020-spectector.pdf \S1); the contracts paper builds on it.}
\notebr
 --- the security properties
  \toolname{}'s contracts target. % \todo{add other contract/SNI works}
\item Axiomatic, happens-before microarchitectural execution models~\cite{herdingcats,lustig}
  --- the modeling style \toolname{} adapts to express leakage predicates.
\end{itemize}
\notebr
\auditnote{cite}{related}{\textbf{CITE}~\cite{guarnieri2021contracts}: Published S\&P 2021 hardware--software contracts (an observation clause and an execution clause per instruction). Map our three components onto its two clauses once, in the leakage-contracts paragraph proposed for Sec.~\ref{sec:bg-defenses}, and point back to it here; the protection set has no counterpart in a contract.}
\notebr
\auditnote{cite}{related}{\textbf{CITE}~\cite{guarnieri2020spectector}: Defines speculative non-interference (SNI) and checks it on software by symbolic execution (guarnieri-2020-spectector.pdf \S1). Cite it as the end-to-end relational property that \toolname{}'s per-execution leakage predicate is contrasted with (as the intro does), not as what a leak means in \toolname{}.}
\notebr
\auditnote{cite}{related}{\textbf{CITE}~\cite{cheang2019formal}: Early formal definition of secure speculation as an information-flow property, trace-property-dependent observational determinism (TPOD), not a trace property; background for what it means for a defense to prevent speculative leakage. (From the 2026-11-09 audit's reading of ePrint 2019/310; the PDF is not in the local library.)}
\notebr
\auditnote{cite}{related}{\textbf{CITE}~\cite{disselkoen2019code}: Models speculative evaluation attacks with a pomset-style (relaxed-memory-inspired) semantics; closest prior bridge between MCM-style semantics and Spectre.}
\notebr
\auditnote{cite}{related}{\textbf{CITE}~\cite{ponce2022cats}: Published S\&P 2022 version of Cats vs. Spectre (replaces CoRR entry); an axiomatic (cat) model of speculative execution, the closest axiomatic analogue to SimSpect's Alloy models.}
\notebr
\auditnote{cite}{related}{\textbf{CITE}~\cite{guanciale2020inspectre}: InSpectre: formal framework modeling microarchitectural features (OoO, speculation) and analyzing countermeasures; closest formal-model-of-defense prior work.}
\notebr
\auditnote{cite}{related}{\textbf{CITE}~\cite{cauligi2022sok}: SoK of formal foundations for software Spectre defenses; survey anchor for the formal-models subsection.}
\notebr
\auditnote{cite}{related}{\textbf{CITE}~\cite{lcms} (already in the bib): Leakage containment models: axiomatic contracts in the uhb/Alloy style; SimSpect's most direct conceptual parent.}
\notebr
\auditnote{cite}{related}{\textbf{CITE}~\cite{synthlc} (already in the bib): Wang et al., CCS 2025, synthesis of sound and precise leakage contracts for open-source RISC-V processors. Despite its key, this entry is not SynthLC: SynthLC (Hsiao et al., MICRO 2024) is in the RTL2M$\mu$PATH paper, \texttt{rtl2upath}. Consider renaming this key.}
\notebr
\auditnote{cite}{related}{\textbf{CITE}~\cite{leave} (already in the bib): Verifies RTL against leakage contracts.}
\notebr
\auditnote{cite}{related}{\textbf{CUT LIST} (2026-11-03 triage): 11 more candidates for part (2) contracts and security definitions are in the source comment just below this note. Their bib entries stay, so any of them can come back.}
\notebr
% ---- CUT from part (2) contracts and security definitions in the 2026-11-03 triage (bib entries kept). To restore one, delete its leading "% ". ----
% \auditnote{cite}{related}{\textbf{CITE}~\cite{cauligi2020constanttime}: Speculative constant-time semantics (Pitchfork); a software-side formal speculation model contrasting with SimSpect's hardware-defense models.}\notebr
% \auditnote{cite}{related}{\textbf{CITE}~\cite{mohr2024synthesizing}: Synthesizes leakage contracts for RISC-V processors from test cases; contrasts with SimSpect, which starts from a defense model and generates tests.}\notebr
% \auditnote{cite}{related}{\textbf{CITE}~\cite{mosier2024serberus}: Trippel-group compile-time Spectre protection for crypto code built on contract reasoning; related formal-speculation work from the same line as lcms.}\notebr
% \auditnote{cite}{related}{\textbf{CITE}~\cite{fabian2022combinations}: Semantics combining multiple speculation sources; relevant to SimSpect modeling multiple speculation primitives in one contract.}\notebr
% \auditnote{cite}{related}{\textbf{CITE}~\cite{patrignani2021exorcising}: Secure-compilation view of speculative non-interference; formal speculation background.}\notebr
% \auditnote{cite}{related}{\textbf{CITE}~\cite{vassena2021blade}: Type-based automatic insertion of protections against speculative leaks (Blade); formal speculation semantics with transmitter/protect notions similar to SimSpect's protection set.}\notebr
% \auditnote{cite}{related}{\textbf{CITE}~\cite{barthe2021highassurance}: Speculative constant-time for verified crypto (Jasmin); software-side formal speculation.}\notebr
% \auditnote{cite}{related}{\textbf{CITE}~\cite{shivakumar2023declassified}: Shows declassification interacts with speculation; relevant to SimSpect's distinction between speculative and non-speculative (declassified) data.}\notebr
% \auditnote{cite}{related}{\textbf{CITE}~\cite{shivakumar2023typing}: Type system for Spectre v1 in high-speed crypto (selective SLH); software formal speculation work.}\notebr
% \auditnote{cite}{related}{\textbf{CITE}~\cite{daniel2021haunter}: Relational symbolic execution for Spectre-PHT/STL (Binsec/Haunted); software-level speculative leak detection including store-to-load forwarding.}\notebr
% \auditnote{cite}{related}{\textbf{CITE}~\cite{mcilroy2019spectre}: Formal model of speculation showing software mitigations are insufficient; motivation for hardware defenses that SimSpect validates.}\notebr
\auditnote{cite}{related}{\textbf{PROPOSED NEW SUBSECTION} (3): axiomatic memory models and model formalization.}
\notebr
\auditnote{cite}{related}{\textbf{CITE}~\cite{manerkar2017rtlcheck}: Verifies RTL against axiomatic uhb models using litmus tests; analogous to SimSpect checking implementations against axiomatic models.}
\notebr
\auditnote{cite}{related}{\textbf{CITE}~\cite{memalloy} (already in the bib): Alloy-based comparison of MCMs and distinguishing litmus tests.}
\notebr
\auditnote{cite}{related}{\textbf{CITE}~\cite{herdingcats} (already in the bib): cat axiomatic MCM framework.}
\notebr
\auditnote{cite}{related}{\textbf{CUT LIST} (2026-11-03 triage): 23 more candidates for part (3) axiomatic models are in the source comment just below this note. Their bib entries stay, so any of them can come back. Most of these are already proposed in the Background notes (Sec.~\ref{sec:bg-axiomatic}), so here just point back to that section.}
\notebr
% ---- CUT from part (3) axiomatic models in the 2026-11-03 triage (bib entries kept). To restore one, delete its leading "% ". ----
% \auditnote{cite}{related}{\textbf{CITE}~\cite{lustig2014pipecheck}: Introduces microarchitectural happens-before (uhb) axiomatic pipeline models; foundation of the Check-tools lineage SimSpect builds on.}\notebr
% \auditnote{cite}{related}{\textbf{CITE}~\cite{lustig2016coatcheck}: Extends uhb modeling to the hardware-OS interface; part of the Check suite of axiomatic microarchitecture models.}\notebr
% \auditnote{cite}{related}{\textbf{CITE}~\cite{manerkar2015ccicheck}: uhb graphs for the coherence-consistency interface; axiomatic microarchitecture modeling.}\notebr
% \auditnote{cite}{related}{\textbf{CITE}~\cite{trippel2017tricheck}: Full-stack (HLL/ISA/uarch) MCM verification with litmus tests; methodological ancestor of SimSpect's model-to-test pipeline.}\notebr
% \auditnote{cite}{related}{\textbf{CITE}~\cite{lustig2015armor}: Formal ordering-specification (MOSTs) for translating between MCMs; axiomatic MCM modeling.}\notebr
% \auditnote{cite}{related}{\textbf{CITE}~\cite{manerkar2020realitycheck}: Modular/hierarchical axiomatic microarchitectural MCM verification (preprint).}\notebr
% \auditnote{cite}{related}{\textbf{CITE}~\cite{manerkar2018pipeproof}: Proves uhb microarchitectural models correct for all programs, not just litmus tests; contrasts bounded test generation with unbounded proof.}\notebr
% \auditnote{cite}{related}{\textbf{CITE}~\cite{norman2023pipesynth}: Synthesizes uhb microarchitectural axioms; related to deriving axiomatic models automatically.}\notebr
% \auditnote{cite}{related}{\textbf{CITE}~\cite{hsiao2021rtl2uspec}: rtl2uspec synthesizes uhb models from RTL; related work on obtaining axiomatic hardware models.}\notebr
% \auditnote{cite}{related}{\textbf{CITE}~\cite{batty2011mathematizing}: Axiomatic formalization of the C++11 memory model; canonical axiomatic MCM work.}\notebr
% \auditnote{cite}{related}{\textbf{CITE}~\cite{sarkar2011power}: Operational POWER model validated against hardware with litmus tests; model-vs-hardware testing analogous to SimSpect model-vs-gem5.}\notebr
% \auditnote{cite}{related}{\textbf{CITE}~\cite{pulte2018simplifying}: Axiomatic and operational ARMv8 models; canonical axiomatic MCM work.}\notebr
% \auditnote{cite}{related}{\textbf{CITE}~\cite{flur2016modelling}: Operational ARMv8 concurrency+ISA model tested against hardware.}\notebr
% \auditnote{cite}{related}{\textbf{CITE}~\cite{simner2020ifetch}: Simner et al.: relaxed semantics for instruction fetch in Armv8-A, extending axiomatic models beyond plain loads/stores.}\notebr
% \auditnote{cite}{related}{\textbf{CITE}~\cite{simner2022vm}: Simner et al.: relaxed virtual-memory semantics for Armv8-A; axiomatic system-level modeling.}\notebr
% \auditnote{cite}{related}{\textbf{CITE}~\cite{simner2025exceptions}: Simner et al. (ISCA 2025): precise exceptions in relaxed architectures; axiomatic modeling at an architecture venue, closest Simner paper to 'relaxed exception semantics'.}\notebr
% \auditnote{cite}{related}{\textbf{CITE}~\cite{armstrong2019sail}: Sail ISA semantics for ARMv8-A/RISC-V/CHERI; formal ISA substrate used by axiomatic tools.}\notebr
% \auditnote{cite}{related}{\textbf{CITE}~\cite{alglave2010fences}: Generic axiomatic framework for weak memory models with fences (precursor of herd/cat).}\notebr
% \auditnote{cite}{related}{\textbf{CITE}~\cite{lahav2017repairing}: RC11 repair of C11 axiomatic model; example of axiomatic model debugging/refinement.}\notebr
% \auditnote{cite}{related}{\textbf{CITE}~\cite{alglave2018linux}: Linux-kernel memory model in cat; axiomatic model of a real-world specification.}\notebr
% \auditnote{cite}{related}{\textbf{CITE}~\cite{owens2009x86tso}: x86-TSO; the axiomatic/operational model of the ISA SimSpect targets (x86 gem5).}\notebr
% \auditnote{cite}{related}{\textbf{CITE}~\cite{memsynth} (already in the bib): Synthesis of axiomatic memory models.}\notebr
% \auditnote{cite}{related}{\textbf{CITE}~\cite{jackson2002alloy} (already in the bib): The modeling language SimSpect uses.}\notebr
\auditnote{cite}{related}{\textbf{PROPOSED NEW SUBSECTION} (4): execution models of speculation.}
\notebr
\auditnote{cite}{related}{\textbf{CITE}~\cite{rtl2upath} (already in the bib): Synthesizes uPATH models and leakage signatures from RTL; complementary bottom-up modeling of speculation. Includes SynthLC (leakage-contract synthesis).}
\notebr
\auditnote{cite}{related}{\textbf{CITE}~\cite{pensieve} (already in the bib): models early-stage defense designs and checks them against speculative non-interference by bounded model checking (yang-2023-pensieve.pdf abstract). It is the prior work closest to our early-stage framing, so say the difference: it checks a model of the design against a generic property, while \toolname{} tests the simulated implementation against the defense's own guarantee.}
\notebr
\auditnote{cite}{related}{\textbf{CUT LIST} (2026-11-03 triage): 3 more candidates for part (4) execution models of speculation are in the source comment just below this note. Their bib entries stay, so any of them can come back.}
\notebr
% ---- CUT from part (4) execution models of speculation in the 2026-11-03 triage (bib entries kept). To restore one, delete its leading "% ". ----
% \auditnote{cite}{related}{\textbf{CITE}~\cite{canella2019systematic}: Systematic taxonomy of transient execution attacks and defenses; framing for the attack/defense space SimSpect covers.}\notebr
% \auditnote{cite}{related}{\textbf{CITE}~\cite{primeandprobe} (already in the bib): Cache side-channel transmitter background.}\notebr
% \auditnote{cite}{related}{\textbf{CITE}~\cite{gem5} (already in the bib): Simulator substrate.}\notebr

\subsection{Spectre hardware defenses}
\begin{itemize}\setlength{\itemsep}{1pt}
\item The hardware defenses \toolname{} expresses and tests: STT~\cite{yu2019stt}, SpecShield,
  NDA~\cite{nda}, DoM, SDO~\cite{sdo}, SPT~\cite{spt}
\reviewtag{[TODO: SPO was a typo for SDO]}%
, and related schemes~\cite{doppelganger,sdo,protean}.
\urgentbr
\urgentnote{This list is out of date: \toolname{} has Alloy models only for STT, ReCon and SPT, and tests InvisiSpec, CleanupSpec and SpecLFB with STT's test set; SpecShield, NDA, DoM and SDO are not expressed or tested. \textbf{PROPOSED TEXT} (edits the item just before this note): The hardware defenses \toolname{} expresses: STT~\cite{yu2019stt}, ReCon~\cite{recon} and SPT~\cite{spt}; it also tests InvisiSpec~\cite{yan2018invisispec}, CleanupSpec~\cite{saileshwar2019cleanupspec} and SpecLFB~\cite{cheng2024speclfb} with STT's test set and a defense-specific leakage criterion. \textit{Source:} simspect STAGE1\_alloy/models/; specs/leakage/*.json.}
\notebr
  % \todo{tighten/complete the defense citations}
\end{itemize}
\notebr
\auditnote{cite}{related}{Defense breaks and SoKs, good motivation for systematic validation.}
\notebr
\auditnote{cite}{related}{\textbf{CITE}~\cite{hu2021sok}: SoK of hardware defenses against speculative execution attacks; survey anchor for the defense list.}
\notebr
\auditnote{cite}{related}{\textbf{CITE}~\cite{li2022unxpec}: Attack on undo-based defenses (CleanupSpec); example of a defense flaw found after publication, motivating systematic validation.}
\notebr
\auditnote{cite}{related}{\textbf{CITE}~\cite{behnia2021interference}: Speculative interference attacks break invisible-speculation schemes; motivates validating defenses against their stated threat model.}
\notebr
\auditnote{cite}{related}{\textbf{CITE}~\cite{chakraborty2022timed}: Timed speculative attacks via store-to-load forwarding bypass cache-based countermeasures; relevant to SimSpect's STLF/SLF findings.}
\notebr
\auditnote{cite}{related}{\textbf{CUT LIST} (2026-11-03 triage): 2 more candidates for part (5) hardware defenses and breaks are in the source comment just below this note. Their bib entries stay, so any of them can come back.}
\notebr
% ---- CUT from part (5) hardware defenses and breaks in the 2026-11-03 triage (bib entries kept). To restore one, delete its leading "% ". ----
% \auditnote{cite}{related}{\textbf{CITE}~\cite{aimoniotis2022doin}: DOIN!: attack on speculative delay schemes via inclusive caches; another post-publication defense break.}\notebr
% \auditnote{cite}{related}{\textbf{CITE}~\cite{spectre} (already in the bib): The attack class defenses target.}\notebr

\subsection{Litmus testing and test generation}
\begin{itemize}\setlength{\itemsep}{1pt}
\item Memory-consistency litmus-test methodology~\cite{herdingcats} that \toolname{}'s
  minimal-execution enumeration builds on. % \todo{add Spectre litmus benchmark + fuzzing cites}
\end{itemize}
\notebr
\auditnote{cite}{related}{\textbf{CITE}~\cite{trippel2018meltdownprime}: Trippel et al. synthesize Spectre/Meltdown variants with CheckMate-style microarchitectural happens-before models; precursor of SimSpect's model-driven exploit synthesis.}
\notebr
\auditnote{cite}{related}{\textbf{CITE}~\cite{hossain2020transform}: Published ISCA 2020 TransForm (replaces CoRR `transform`); Alloy-based axiomatic transistency models with synthesized litmus tests, the most direct methodological precedent for SimSpect.}
\notebr
\auditnote{cite}{related}{\textbf{CITE}~\cite{lustig} (already in the bib): Alloy-based exhaustive minimal litmus-test synthesis; SimSpect's generation methodology descends from it.}
\notebr
\auditnote{cite}{related}{\textbf{CITE}~\cite{checkmate} (already in the bib): uhb-based security litmus test synthesis; SimSpect's direct precursor.}
\notebr
\auditnote{cite}{related}{\textbf{CUT LIST} (2026-11-03 triage): 6 more candidates for part litmus testing and test generation are in the source comment just below this note. Their bib entries stay, so any of them can come back.}
\notebr
% ---- CUT from part litmus testing and test generation in the 2026-11-03 triage (bib entries kept). To restore one, delete its leading "% ". ----
% \auditnote{cite}{related}{\textbf{CITE}~\cite{armstrong2021isla}: Isla: combines full ISA semantics with axiomatic concurrency models to evaluate litmus tests; tool-level analogue of SimSpect's concretization.}\notebr
% \auditnote{cite}{related}{\textbf{CITE}~\cite{maradhaim2010generating}: Generates litmus tests that distinguish memory models; the original model-driven litmus generation idea SimSpect applies to defenses.}\notebr
% \auditnote{cite}{related}{\textbf{CITE}~\cite{maradhaim2011litmus}: DAC 2011 follow-up: litmus tests for comparing memory consistency models.}\notebr
% \auditnote{cite}{related}{\textbf{CITE}~\cite{alglave2011litmus}: litmus tool for running tests on real hardware; analogue of SimSpect's gem5 test harness.}\notebr
% \auditnote{cite}{related}{\textbf{CITE}~\cite{perple} (already in the bib): Litmus-test execution on hardware.}\notebr
% \auditnote{cite}{related}{\textbf{CITE}~\cite{mcmutants} (already in the bib): Mutation testing of MCM test suites; analogous to evaluating test-suite quality.}\notebr
\auditnote{cite}{related}{\textbf{CUT LIST} (2026-11-03 triage): 23 more candidates for part extra hardware defenses (not in the intro's lean list) are in the source comment just below this note. Their bib entries stay, so any of them can come back.}
\notebr
% ---- CUT from part extra hardware defenses (not in the intro's lean list) in the 2026-11-03 triage (bib entries kept). To restore one, delete its leading "% ". ----
% \auditnote{cite}{related}{\textbf{CITE}~\cite{aimoniotis2024janus}: JANUS [delay, RL-guided]: reinforcement-learning policy for selective speculative defense. (Mechanism tag inferred from title/abstract snippet, not the full paper; confirm before citing in a mechanism-specific sentence.)}\notebr
% \auditnote{cite}{related}{\textbf{CITE}~\cite{castes2024sharing}: Core-gapped confidential VMs [partition/isolation]: blocks cross-domain transient attacks by not sharing cores (system-level). (Mechanism tag inferred from title/abstract snippet, not the full paper; confirm before citing in a mechanism-specific sentence.)}\notebr
% \auditnote{cite}{related}{\textbf{CITE}~\cite{chen2025mustt}: muSTT [taint-tracking]: microarchitecture design for speculative taint tracking. (Mechanism tag inferred from title/abstract snippet, not the full paper; confirm before citing in a mechanism-specific sentence.)}\notebr
% \auditnote{cite}{related}{\textbf{CITE}~\cite{choudhary2023declassiflow}: Declassiflow [static analysis, software]: models non-speculative knowledge to relax hardware defenses like STT/SPT; software side of the defense space.}\notebr
% \auditnote{cite}{related}{\textbf{CITE}~\cite{cui2023ssft}: Speculative Secret Flow Tracking [taint-tracking]: tracks speculatively accessed secret flows (JCST). (Mechanism tag inferred from title/abstract snippet, not the full paper; confirm before citing in a mechanism-specific sentence.)}\notebr
% \auditnote{cite}{related}{\textbf{CITE}~\cite{daniel2023prospect}: ProSpeCT [taint-tracking/delay, provable]: tracks secrets and blocks their speculative leakage for constant-time code (Proteus RISC-V).}\notebr
% \auditnote{cite}{related}{\textbf{CITE}~\cite{fustos2019spectreguard}: SpectreGuard [taint/secret annotation, delay]: software-marked secret regions block forwarding of speculative secrets.}\notebr
% \auditnote{cite}{related}{\textbf{CITE}~\cite{ganjisaffar2025specasan}: SpecASan [delay on memory-safety violation]: uses memory tagging to block speculative out-of-bounds accesses.}\notebr
% \auditnote{cite}{related}{\textbf{CITE}~\cite{green2023safebet}: SafeBet [delay with safe-access prediction] (preprint): allows only provably safe speculative accesses. (Mechanism tag inferred from title/abstract snippet, not the full paper; confirm before citing in a mechanism-specific sentence.)}\notebr
% \auditnote{cite}{related}{\textbf{CITE}~\cite{hajiabadi2024levioso}: Levioso [delay, compiler-informed]: compiler-provided dependence info restricts delays to truly dependent instructions. (Mechanism tag inferred from title/abstract snippet, not the full paper; confirm before citing in a mechanism-specific sentence.)}\notebr
% \auditnote{cite}{related}{\textbf{CITE}~\cite{hajiabadi2025cassandra}: Cassandra [sequential-execution enforcement]: replaces branch prediction for crypto code with recorded control flow. (Mechanism tag inferred from title/abstract snippet, not the full paper; confirm before citing in a mechanism-specific sentence.)}\notebr
% \auditnote{cite}{related}{\textbf{CITE}~\cite{jin2023specterminator}: SpecTerminator [delay by instruction class]: blocks speculative side channels on RISC-V based on transmitter classes.}\notebr
% \auditnote{cite}{related}{\textbf{CITE}~\cite{kim2024perspective}: Perspective [OS-level speculation control]: principled framework for secure speculation in the OS (HW/OS co-design). (Mechanism tag inferred from title/abstract snippet, not the full paper; confirm before citing in a mechanism-specific sentence.)}\notebr
% \auditnote{cite}{related}{\textbf{CITE}~\cite{kiriansky2018dawg}: DAWG [partition]: way-partitioned caches isolate protection domains.}\notebr
% \auditnote{cite}{related}{\textbf{CITE}~\cite{kvalsvik2025shadowbinding}: ShadowBinding [taint-tracking/delay microarchitecture]: realizable in-core implementations of secure speculation schemes like STT/NDA. (Mechanism tag inferred from title/abstract snippet, not the full paper; confirm before citing in a mechanism-specific sentence.)}\notebr
% \auditnote{cite}{related}{\textbf{CITE}~\cite{li2019conditional}: Conditional Speculation [delay]: blocks only security-dependent speculative loads (with cache-hit filters).}\notebr
% \auditnote{cite}{related}{\textbf{CITE}~\cite{schwarz2020context}: ConTExT [taint/secret annotation, delay]: memory marked secret is never used speculatively (HW+SW).}\notebr
% \auditnote{cite}{related}{\textbf{CITE}~\cite{taram2019csf}: Context-Sensitive Fencing [delay/fence]: microcode-injected fences on demand.}\notebr
% \auditnote{cite}{related}{\textbf{CITE}~\cite{tran2020clearing}: Clearing the Shadows [delay, HW/SW co-design]: compiler-assisted reduction of delay-scheme overhead.}\notebr
% \auditnote{cite}{related}{\textbf{CITE}~\cite{trujillo2026register}: Register hiding + ShadowCFI [limit secret reachability]: restricts which secrets transient execution can reach (S\&P 2026). (Mechanism tag inferred from title/abstract snippet, not the full paper; confirm before citing in a mechanism-specific sentence.)}\notebr
% \auditnote{cite}{related}{\textbf{CITE}~\cite{wu2025rcp}: Reversible Coherence Protocol [undo/reversible coherence]: published TDSC version of the ReversiSpec/RCP preprint (arXiv 2006.16535). (Mechanism tag inferred from title/abstract snippet, not the full paper; confirm before citing in a mechanism-specific sentence.)}\notebr
% \auditnote{cite}{related}{\textbf{CITE}~\cite{zhao2020invarspec}: InvarSpec [delay relaxation]: program analysis identifies speculation-invariant instructions safe to execute early.}\notebr
% \auditnote{cite}{related}{\textbf{CITE}~\cite{zhao2022pinned}: Pinned Loads [delay relaxation]: pins loads early so they become non-speculative sooner (reduces delay of secure schemes).}\notebr

\section{Discussion}
\iffalse % LAB-REVIEW: discussion to-do list hidden (2026-11-03)
{\color{cfigtodo}\footnotesize\noindent[fig todo: major discussion points to include]
\begin{itemize}
\item \textbf{Bounded exhaustiveness}: complete up to the instruction bound where Stage~1 reaches UNSAT (so far the STT and ReCon models); SPT's Stage~1 stops at its cap, so SPT is not bounded-exhaustive; Stage~1.5 is non-exhaustive by design.
% STAGE NAMES (2026-11-11): in the item above, Stage~1 is now Stage~1a and Stage~1.5 is Stage~1b (Fig.~pipeline renamed).
\item \textbf{Trust base}: each defense's specification is translated by hand into a speculation contract, protection set and leakage function; SpecLFB, InvisiSpec and CleanupSpec currently reuse STT's model; note the per-defense exceptions in the shared skeleton.
\item \textbf{Transmission mechanisms} are defense-specific and hand-written (JSON criteria); STT and Recon still use the generic load gate.
\item \textbf{Harness}: probes never change defense logic, but the speculation and timing directives change prediction and timing on purpose; runs that do not realize the modeled speculation are invalid, never hits.
\item \textbf{Window semantics}: a branch counts as unresolved until its squash lands (decided 2026-11-08), so a transmission after resolution but before the squash (the race) is a violation.
\item \textbf{Threat models differ} across implementations of the same defense (STT-Spectre vs.\ Recon-STT); results are per implementation.
\item \textbf{Complementary to random testing}: \toolname{} replicates the AMuLeT bugs that break a defense's stated guarantee (A-UV1, A-UV3--A-UV6); the others (MSHR contention, instruction-cache fetch, cleanup timing, the store\,$\to$\,TLB leak) lie outside the stated guarantee or outside our model.
\item \textbf{Generality}: other simulators or RTL; more speculation primitives (store-to-load forwarding done; e.g.\ value prediction next).
\auditnote{claim}{claim, low}{Store-to-load forwarding is not done as a primitive. The STL model, stl\_bypass mode and hooks exist and passed controls on 60 tests, but no STL window checker or STL campaign exists, and the ReCon STLF bug came from ordinary forwarding in a branch-speculation test.  \textbf{PROPOSED TEXT} (edits the sentence just before this note): more speculation primitives (store-to-load forwarding in progress; e.g.\textbackslash{} value prediction next)  \textit{Source:} project log 2026-10-30 stl-spec (60-test controls); project log 2026-10-28 16:30 ('The window checker is proposed, not built'); URGENT\_TODO.md fig (d) DECIDE}
\end{itemize}}
\fi
- Sentence about 

A defense's execution contract conservatively specifies which instructions are speculative, and therefore protected. The space of \emph{real} leakage, however, is the intersection of that contract with the speculation a given microarchitecture actually performs (gem5's \emph{speculation primitives}). This leaves a design choice: do we test a contract against the broad space it specifies, or against its projection onto realizable speculation? Testing the realizable intersection is more charitable to implementers --- it surfaces only bugs that can actually manifest --- but is arguably incomplete, since a correct implementation should uphold the \emph{entire} promised contract. The same tension recurs in whether to require misprediction or to also exercise correctly-predicted but still-speculative windows
\urgentbr
\urgentnote{\textbf{VOCAB}: Here ``still-speculative windows'' means stretches where a correctly predicted branch is still unresolved, not the checker's violation window. Reword rather than rename, e.g.\ ``to also exercise correctly predicted branches while they are still unresolved'', so that ``window'' keeps one meaning.}%
\notebr
, and in whether a conservative window (protecting every instruction after an unresolved branch, whether or not it is data-dependent) must be honored even where no real dataflow path to a transmitter exists. \toolname{} can target either interpretation --- for several defenses we maintain both the raw and the intersected execution contract
\notebr
\auditnote{fact}{fact, low}{Unsupported: every model in the evaluation uses the branch-only execution contract, and no defense has both a raw and an intersected contract in use. Broader STT contracts exist only as a draft awaiting approval (store shadows, for ReCon's STT) or unused (Futuristic). \newline \textbf{PROPOSED TEXT} (replaces the clause just before this note): \textbackslash{}toolname\{\} can target either interpretation (for STT we have also written broader execution contracts, which we leave to future work) \newline \textit{Source:} simspect URGENT\_TODO.md DECIDE ``approve the draft STAGE1\_alloy/models/proposed/STT\_6\_shadows.als''; wiki/design/FUTURISTIC\_MODEL.md.}
\notebr
--- and we default to the realizable intersection, since it is what yields demonstrable, root-caused bugs on the running prototype.

\anna{2. Alloy abstract LILT of the STT load load, 3. the enumerated, annotated veersion 4. a really minimalist version of my pipelineout of running on gem5 with the relevant windows circled and completion check occuring -> bug}

\urgentbr\urgentnote{\textbf{PROPOSED TEXT (checked guarantees and scope).} \toolname{} checks each defense against the guarantee its authors state: each footprint criterion lists only the structures and state that the defense's paper claims to protect, together with the sentence that claims it, and where we test less than a paper claims we say so (SpecLFB's protection after faulting instructions, below). For InvisiSpec, that is the data-cache hierarchy, meaning the occupancy, replacement and coherence state of every level, plus the data TLB (\cite[\S III-B]{yan2018invisispec}). Two of the vulnerabilities AMuLeT reports against InvisiSpec fall outside this statement. Speculative instruction fetches into the L1 instruction cache (A-KV1 in~\cite{amulet}) are excluded by the paper's footnote~2: ``For simplicity, we do not describe the protection of the instruction cache or instruction TLB.'' MSHR contention (A-UV2 in~\cite{amulet}) is excluded explicitly in the security analysis: ``contention on other resources, such as MSHRs or execution units, could slow down the receiver, but such side channels are considered out of scope in this paper'' (\cite[\S VII]{yan2018invisispec}). We therefore treat these as limits of the defense's stated scope rather than implementation bugs, and our checker does not flag them, whereas the speculative L1D eviction (A-UV1), which violates the stated data-cache guarantee, is flagged. The same rule applies to the other footprint defenses. CleanupSpec states its guarantee for the data-cache hierarchy (L1-D and L2/LLC); it acknowledges that cleanup time varies and leaves a constant-time cleanup to future work (\cite{saileshwar2019cleanupspec}, \S 4(b)), so a cleanup-latency difference such as unXpec (A-KV2) is not checked. CleanupSpec's failures to undo a speculative install (A-UV3--A-UV5) do violate its stated guarantee, and our checker flags them. SpecLFB excludes TLB and port-contention channels and loads that hit in the cache~\cite{cheng2024speclfb}. \textit{Sources:} InvisiSpec fn.~2 p3 and \S VII p10 of the original MICRO version; CleanupSpec \S 4(b) p8 and \S 2.1 p3; SpecLFB \S 4 p635 and \S 9 p643; criterion files specs/leakage/\{invisispec,cleanupspec,speclfb\}.json.}

\urgentbr\urgentnote{\textbf{PROPOSED TEXT (STT and store addresses).} AMuLeT also reports that a speculative store with a tainted address installs a D-TLB entry under STT (A-KV3 in~\cite{amulet}), a leak first identified by DOLMA~\cite{loughlin2021dolma}. STT ties its protection of translation state to delaying loads (``for loads, delaying execution implies delaying the TLB lookup'', \cite[footnote~2]{yu2019stt}) and ``does not require tracking taint in any part of the memory system (TLB, caches, or memory)''. Our STT model's leakage function contains load addresses and branch predicates, the instructions STT delays, and not store addresses, so \toolname{} does not generate tests aimed at A-KV3. STT also assumes that stores act only at retirement: ``We assume that stores trigger a cache coherence invalidation only on retirement, or else are defined as transmitters'' (\cite{yu2019stt}, p5), and ``In isolation, a store creates neither explicit or implicit channels because we assume stores are performed at retirement'' (p10). Our model makes the same assumption, so A-KV3, a store that translates its address before retirement, lies outside what we check. \textit{Note:} the first quote's alternative (``or else are defined as transmitters'') lets a reader argue that a store acting early is an STT transmitter; the project's decision (simspect URGENT\_TODO.md, ``store addresses under STT'') keeps A-KV3 out of scope, so word this as our model's assumption, not as STT's exclusion.}

\urgentbr\urgentnote{\textbf{PROPOSED TEXT (structure sizes).} AMuLeT amplifies contention-dependent leaks by testing with smaller microarchitectural structures; in its InvisiSpec campaign it reduces the L1D from 8 ways to 2 and the MSHRs from 256 to 2 (\cite[\S 4.5.1, Table~6]{amulet}). We apply the same idea to every footprint defense and run with a small L1D (1--2\,kB, 2--4 ways). Small structures let a short litmus test reach the full-set and eviction conditions under which these bugs appear, and they do not change the defense logic. All timing and throughput comparisons with AMuLeT instead use exactly AMuLeT's structure sizes and configuration.}

\urgentbr\urgentnote{\textbf{PROPOSED TEXT (cache sizes for the taint defenses).} We do not shrink caches for the taint-tracking defenses (STT, ReCon, SPT): their violation test asks whether a transmitter obtained its data inside the violation window, which does not depend on cache pressure, so smaller caches would amplify nothing. Any new footprint bug found on the small configuration is re-run once at AMuLeT's exact configuration before we report it.}
\urgentbr\urgentnote{\textbf{PROPOSED TEXT (why the pending-squash bug is not reached).} A fix in the upstream STT repository (December 2025) addresses a bug that \toolname{} does not reach: STT delays a tainted branch's squash to hide its direction, but when an older tainted branch and a younger untainted branch mispredict in the same cycle, the untainted branch's squash is delayed too, so its redirect timing reveals the tainted branch's direction~\cite{yu2019stt}. Two things keep it out of reach. First, the observable is the squash signal of an \emph{untainted} branch, which is not a transmitter in our STT specification (load addresses and the predicates of tainted branches); our security predicate calls such a gadget secure, and capturing it needs a squash-signal channel in the model, which we have drafted but not evaluated. Second, even with that channel, a test of our size cannot make the signal depend on the secret: tainting the older branch requires a speculative load, which requires a still older unresolved branch, and gem5 holds both pending squashes until that oldest branch resolves, so the younger branch's redirect happens when the window closes whatever the secret is. Comparing a tainted with an untainted older branch shows the younger branch being held (17 of 20 generated gadgets), but varying only the secret changes nothing observable (0 of 12). Reaching it likely needs the two branches to resolve in the same cycle while the window is opened by something other than an older branch, which we leave to future work.}
\urgentbr\urgentnote{\textbf{PROPOSED TEXT (what we leave out of scope, and why this is closer to each paper).} \toolname{} judges each prototype against the guarantee its own paper states. A report then means that the implementation breaks a promise its authors made, which is a bug they can fix in the code. Checking every defense against a generic leakage contract, as AMuLeT does~\cite{amulet}, instead mixes such bugs with channels the authors chose not to defend, which no change to the implementation would remove. Of the reports we leave out, three kinds follow from the papers or from our specifications. First, channels the defense's paper excludes: InvisiSpec's MSHR contention (A-UV2) and instruction cache (A-KV1), above. Second, the timing of a defense's own actions: CleanupSpec acknowledges that cleanup time varies and leaves a constant-time cleanup to future work (\cite{saileshwar2019cleanupspec}, \S 4(b)); unXpec~\cite{li2022unxpec} shows the variation is exploitable, and AMuLeT observes it (A-KV2) through the instruction cache (\cite{amulet}, p12), which CleanupSpec's threat model excludes (\cite{saileshwar2019cleanupspec}, \S 2.1). Third, transmitters our specifications do not include: speculative stores (STT's TLB leak, A-KV3, above), and the squash signal of an untainted branch (the pending-squash bug, above). \textit{Pending decision} (simspect URGENT\_TODO.md, ``store addresses under STT''): this list used to include the tainted-address store channel in ReCon's implementation (a tainted store address resolves and squashes an aliasing younger load). STT claims to close exactly that channel (its memory-dependence example on p5, and p8: delay squashing ``until the load's address and the addresses of all older stores become untainted''), so it may belong under what we do not yet cover instead; it is left out of both lists until decided. \textit{Check before use:} that ``generic leakage contract'' is a fair description of how AMuLeT chooses its contract per defense.}

\urgentbr\urgentnote{\textbf{PROPOSED TEXT (what we do not yet cover).} Some reports fall within a paper's claim but outside what \toolname{} currently generates; we list them as coverage limits rather than scope. SpecLFB's paper protects loads that follow a possibly faulting instruction, and source reading shows its implementation does not; our speculation contract models only branch and store-address speculation, so \toolname{} cannot yet test this. Source reading also finds candidate bugs that need tests our generator cannot yet produce: interactions between cores (in InvisiSpec, a speculative request that downgrades another core's copy of a line), instructions outside our vocabulary (SPT's untaint rules for flags, conditional moves and 8- and 16-bit writes, and ReCon's floating-point loads, which are never marked speculative), and shapes beyond our bound (a squash in the middle of the window while an older shadow is still open). Each needs an extension of the model or the vocabulary, not of the method.}

\urgentbr\urgentnote{\textbf{PROPOSED TEXT (sketch: why forcing predictions and stretching windows is sound).} The harness changes two things in the simulated core (Speculation Primitive Instrumentation): where the predictors send execution, by forcing each annotated branch to mispredict or to predict correctly and by forcing a load to be predicted not to alias an older store, and when an unresolved branch may resolve or the final instruction may commit, by stalling them for a set number of cycles. (a) This does not change the defense under test. The hooks act only on predictor outputs and on when resolution or commit proceeds; they add prediction overrides, stalls and trace prints, and never modify the defense's own state or code paths (taint, protection flags, speculative buffers, cleanup records). The defense therefore sees an ordinary misprediction and a branch whose condition arrives late, inputs it must already handle. (b) It does not create leakage that a regular predictor could not. In the Spectre threat model the attacker trains the predictors~\cite{spectre}, so any prediction we force is one an attacker can induce, and a late resolution corresponds to a branch condition that misses in the cache, which an attacker can arrange by evicting it. The defenses aim to be secure for any prediction and any speculation window [cite each threat model], so a leak inside a forced, stretched window is one the defense set out to prevent. Two limits remain. A stretched window can make a leak easier to hit than natural timing would, so we confirm new bugs without the hooks where we can (the SpecLFB \texttt{jmp} bug reproduces on an unmodified build with a natural misprediction). Changing a pipeline parameter, such as the delay between commit and issue, changes the microarchitecture itself, so bugs that need it (the visibility-point/squash race) are reported together with that setting.}

\urgentbr\urgentnote{\textbf{PROPOSED TEXT (targets without white-box access; from the old draft, oldtechnicaldetails.tex:42 and :101).} \toolname{} uses gem5's white-box access to force predictions and hold resolution (Gem5 Execution and Instrumentation). A target without that access could realize the same annotations the way an attacker does (argument (b) of the note above): train the predictors before the test so that each annotated branch mispredicts or predicts as its annotation says, and widen the window by evicting the data a branch condition depends on so that the branch resolves late. Detection carries over to RTL simulation, which still exposes every signal; on silicon a hit would have to be measured, for example with a cache-timing probe. \textit{Notes:} none of this is implemented. It is a generality point for the next-steps paragraph or the Conclusion, not a result, and the RTL and silicon sentence is a suggestion, not something we have checked.}

\section{Conclusion}
\urgentbr\urgentnote{\textbf{PROPOSED TEXT (Conclusion; paste as \textbackslash{}section\{Conclusion\} here).} Hardware Spectre defenses are prototyped in simulators that can measure their performance but not their security guarantee. \toolname{} closes that gap: a defense is specified as a speculation contract, a protection set and a leakage function, a synthesizer turns the specification into a compact suite of litmus tests that each exhibit a forbidden execution, and an executor runs the suite on the gem5 prototype and checks whether the execution is prevented. On six prototypes, \toolname{} found all five previously known bugs that violate their defense's own guarantee, six new bugs (see Results), in [N]$\times$ fewer tests than random testing on the same defenses. Three limits bound these claims. Each specification is translated from its paper by hand, so a mistake in the translation is a mistake in the test suite, and the protection sets are the least certain part. Enumeration is bounded: Stage~1a is exhaustive up to its instruction bound only for the STT and ReCon models, and SPT's stops at its cap. A hit is a leak only under the window rule and harness we fix (a branch counts as unresolved until its squash lands, and probes never change defense logic), and we call a hit a defect in the prototype only after triage. Because the specification is separate from the tests, checking a new defense or a fix changes only the three definition clauses. Next steps: describe further defenses and speculation sources (for example value prediction), run the executor on RTL, and turn the manual triage that separated checker artifacts from real bugs into automatic checks. \textit{Notes:} the ``six new bugs'' counts the SpecLFB walk order as new (Anna, 2026-11-09); ``[N]$\times$'' is the same placeholder as the abstract and intro (AMuLeT parity run T5 is blocked); the next-steps sentence is a suggestion from the Discussion to-do list, not a result.}

\begin{acks}
This document is derived from previous conferences, in particular ISCA 2025.
\end{acks}

\bibliographystyle{ACM-Reference-Format}
\bibliography{sections/references}

\end{document}
